\chapter{Implementation}
\label{chap:implementation}

This chapter describes the practical implementation of the theoretical concepts and system architectural components discussed in the previous chapters. The system was realized by leveraging the technologies detailed in Chapter~\ref{chap:background}, with a primary focus on the custom simulation engine built upon the Godot Engine framework and its integration with the BDI reasoning pipeline. The complete codebase for the project is publicly available on GitHub at \url{https://github.com/MatVecchi/war_game}.

% =========================================================================
\section{Game Environment and World Construction}
\label{sec:game_environment}

To establish a realistic strategic scenario for evaluating multi-agent interactions, the spatial topology, provincial boundaries, navigable path networks, resource outputs, and geopolitical borders were modeled based on the historical "The Great War" map template from the grand strategy game \textit{Supremacy 1914}. 

To maintain architectural clarity and computational feasibility while preserving geopolitical complexity, the original $31$-nation world map was scaled down to a representative subset of six sovereign nations: \textbf{Finland}, \textbf{Lithuania}, \textbf{Norway}, \textbf{Sweden}, \textbf{Ukraine}, and \textbf{Russia}.

\begin{figure}[H]
    \centering
    \includegraphics[width=0.9\textwidth]{imgs/supremacy_map.png}
    \caption{World map of the \textit{Supremacy 1914} "The Great War" template}
    \label{fig:world_map_original}
\end{figure}

These specific nations were selected to provide a diverse array of geopolitical configurations, incorporating both terrestrial (land-based) and maritime borders. This spatial heterogeneity enables comprehensive testing of land unit mobility while laying the structural foundation for future extensions involving naval operations. 

During the development phase, high-resolution cartographic reference assets of the original map were imported into the Godot Engine. The national and provincial boundaries were precisely vector-traced to maximize geometric fidelity relative to the source material.

\begin{figure}[H]
    \centering
    \includegraphics[width=0.9\textwidth]{imgs/map_recreation.png}
    \caption{Vector-traced world map representation showcasing national boundaries, provincial capitals.}
    \label{fig:world_map_reference}
\end{figure}

\begin{figure}[H]
    \centering
    \includegraphics[width=0.9\textwidth]{imgs/province_recreation.png}
    \caption{Vector-traced province representation showcasing resource allocations, and inter-provincial transit networks.}
    \label{fig:province_reference}
\end{figure}

As illustrated in Figure~\ref{fig:province_reference}, the environment incorporates several key structural elements:
\begin{itemize}
    \item \textbf{National Capitals and Provincial Seats:} National capitals are visually highlighted via high-contrast label backgrounds. The central node of each province (represented by an icon connecting the path network) acts as the strategic provincial seat. To retain control over a province, a nation must safeguard this specific node against enemy occupation.
    \item \textbf{Resource Output Allocation:} Every province produces a specific strategic commodity, denoted by an iconic indicator positioned above its spatial label (e.g., the \textit{Kharkov} province producing wood for the Ukrainian nation).
    \item \textbf{Transit Path Networks:} Inter-provincial connectivity is explicitly defined by color-coded path vectors connecting adjacent provincial seats. These paths represent the constrained, graph-based routes along which both allied and hostile military divisions can traverse.
    \item \textbf{Military Division Representations:} Armies are rendered as discrete unit tokens positioned dynamically on the map, paired with a specialized User Interface (UI) overlay displaying real-time operational telemetry, including hit points, soldier count, and a unique identifier relative to the owning nation.
\end{itemize}

\begin{figure}[H]
    \centering
    \includegraphics[width=0.3\textwidth]{imgs/army_recreation.png}
    \caption{Example of an army.}
    \label{fig:army_representation}
\end{figure}

\begin{figure}[H]
    \centering
    \includegraphics[width=0.9\textwidth]{imgs/army_travelling.png}
    \caption{Example of an army travelling towards an objective.}
    \label{fig:army_travelling}
\end{figure}

\subsection{Sensing Mechanics and Fog of War}
To enforce bounded perception—a core constraint of realistic BDI agent modeling—each military division is endowed with a multi-tiered operational Vision Radius that dictates environmental information retrieval:
\begin{enumerate}
    \item \textbf{Full Reconnaissance Zone ($0 \le d \le 90\text{ px}$):} Within a primary radius of $90$ pixels, the owning agent possesses full sensory fidelity, acquiring complete telemetry regarding foreign armies (e.g., exact unit count, health metrics, and owner identity).
    \item \textbf{Obfuscated Detection Zone ($90 < d \le 135\text{ px}$):} Within a secondary radius spanning from $90$ to $135$ pixels, the agent operates under partial sensory occlusion. It registers the presence and spatial coordinates of foreign forces but remains oblivious to their specific unit compositions or health metrics.
    \item \textbf{Blind Zone ($d > 135\text{ px}$):} Beyond a $135$-pixel radius, foreign units are completely occluded by the Fog of War. The agent retains access only to static, public environmental characteristics, such as provincial names and baseline resource outputs.
\end{enumerate}

\begin{figure}[H]
    \centering
    \includegraphics[width=0.8\textwidth]{imgs/russia_view.png}
    \caption{Example of Russia's point of view.}
    \label{fig:russia_view}
\end{figure}

\begin{figure}[H]
    \centering
    \includegraphics[width=0.8\textwidth]{imgs/global_view.png}
    \caption{Same example as Figure~\ref{fig:russia_view} but with the spectator point of view.}
    \label{fig:global_view}
\end{figure}

% =========================================================================
\section{Spectator Camera System and User Interaction}
\label{sec:camera_controls}

To facilitate real-time monitoring and empirical evaluation of the BDI agents during simulation runs, a dedicated spectator camera controller was implemented. Integrated within the Godot Engine input pipeline, the spectator interface exposes three core interaction capabilities: continuous spatial panning, logarithmic viewport zooming, and discrete provincial selection.

These interaction mechanics are implemented within the \texttt{Camera.gd} script attached to the \texttt{Camera2D} node. The controller relies on Godot's native event-driven method \texttt{\_input(event)}, which intercepts hardware input events and dispatches them to dedicated handler routines.

\subsection{Camera Control Handlers}
The spectator interface relies on three primary handler functions:
\begin{itemize}
    \item \textbf{\texttt{\_handleZoom(s)}:} Modifies the scale factor of the \texttt{Camera2D} node. The updated zoom vector is computed according to Equation~\ref{eq:zoom_calculation}:
    \begin{equation}
        \label{eq:zoom_calculation}
        \mathbf{Z}_{\text{new}} = \text{clamp}\Big(\mathbf{Z}_{\text{current}} \times \big(1.0 + (\gamma \times s)\big), \mathbf{Z}_{\text{min}}, \mathbf{Z}_{\text{max}}\Big)
    \end{equation}
    where $\gamma = 0.1$ represents the zoom speed constant, and $s \in \{-1.0, 1.0\}$ indicates the directional input scroll. To preserve visual spatial coherence, the global camera position is dynamically re-centered toward the mouse cursor's world coordinates. Furthermore, to optimize visual readability and rendering performance, Level of Detail (LOD) thresholds automatically toggle the visibility of province borders ($\mathbf{Z}_{x} \ge 1.5$) and army telemetry UI elements ($\mathbf{Z}_{x} \ge 2.0$).
    
    \item \textbf{\texttt{\_handlePan(s)}:} Manages the spatial drag state via a boolean flag (\texttt{dragging}). When initiated, it caches the initial cursor and camera transformation vectors, allowing the main engine loop (\texttt{\_process}) to smoothly translate the viewport relative to mouse displacement.
    
    \item \textbf{\texttt{\_unhandled\_input(event)}:} A native Godot fallback handler invoked only when an input event is not consumed by higher-level UI canvas controls (e.g., interface buttons). This prevents accidental province selections when interacting with overlay menus. It invokes the helper function \texttt{\_select\_province\_at\_mouse()}, which performs a point-in-polygon raycasting query using \texttt{Geometry2D.is\_point\_in\_polygon()}. If an underlying province polygon intersects the cursor's local coordinates, the \texttt{GameController} selects that territory; otherwise, active selections are cleared.
\end{itemize}

\subsection{GDScript Implementation}
The implementation details of the spectator camera controller described above are structured in Listing~\ref{lst:camera_script}.

\begin{lstlisting}[language=python, caption={GDScript implementation of the Camera2D controller handling panning, dynamic zooming, Level of Detail (LOD) rendering, and raycasted province selection.}, label={lst:camera_script}]

var dragging: bool = false
var dragged: bool = false
var input_enabled: bool = true
var drag_start_mouse_pos: Vector2 = Vector2.ZERO
var drag_start_camera_pos: Vector2 = Vector2.ZERO
var border_shown: bool = true
var army_shown: bool = true

const MIN_ZOOM: float = 1.0
const MAX_ZOOM: float = 8.0
const ZOOM_SPEED: float = 0.1
const SHOW_BORDERS_ZOOM: float = 1.5
const SHOW_ARMY_COUNT_ZOOM: float = 2.0

func _handleZoom(s: float) -> void:
    if s != 1.0 and s != -1.0: 
        return
    
    var mouse_world_before = get_global_mouse_position()
    zoom = (zoom * (1.0 + (ZOOM_SPEED * s))).clamp(
        Vector2(MIN_ZOOM, MIN_ZOOM), 
        Vector2(MAX_ZOOM, MAX_ZOOM)
    )
    var mouse_world_after = get_global_mouse_position()
    global_position += mouse_world_before - mouse_world_after

    # Dynamic Level of Detail (LOD) for Province Borders
    if zoom.x >= SHOW_BORDERS_ZOOM and not border_shown:
        var provinces: Array[Node] = get_tree().get_nodes_in_group("provinces")
        for province in provinces:
            if province is not Province: continue
            (province as Province).province_border.show()
            (province as Province).control.show()
        border_shown = true
    elif zoom.x < SHOW_BORDERS_ZOOM and border_shown:
        var provinces: Array[Node] = get_tree().get_nodes_in_group("provinces")
        for province in provinces:
            if province is not Province: continue
            (province as Province).province_border.hide()
            (province as Province).control.hide()
        border_shown = false
        
    # Dynamic Level of Detail (LOD) for Army Telemetry UI
    if zoom.x >= SHOW_ARMY_COUNT_ZOOM and not army_shown:
        var armies: Array[Node] = get_tree().get_nodes_in_group("armies")
        for army in armies:
            if army is not Army: continue
            (army as Army).control.show()
        army_shown = true
    elif zoom.x < SHOW_ARMY_COUNT_ZOOM and army_shown:
        var armies: Array[Node] = get_tree().get_nodes_in_group("armies")
        for army in armies:
            if army is not Army: continue
            (army as Army).control.hide()
        army_shown = false

func _handlePan(s: bool) -> void:
    if not dragging and s:
        drag_start_mouse_pos = get_viewport().get_mouse_position()
        drag_start_camera_pos = position
        dragging = true
    
    if dragging and not s:
        dragging = false

func _process(_delta: float) -> void:
    if dragging:
        var moveVector = get_viewport().get_mouse_position() - drag_start_mouse_pos
        if moveVector == Vector2.ZERO: 
            return
        dragged = true
        position = drag_start_camera_pos - moveVector * (1.0 / zoom.x)

func _input(_event: InputEvent) -> void:
    if input_enabled:
        if Input.is_action_just_pressed("camera_pan"):
            _handlePan(true)
            dragged = false
        _handleZoom(Input.get_axis("camera_zoom_out", "camera_zoom_in"))
    if Input.is_action_just_released("camera_pan"):
        _handlePan(false)
        dragged = false

func _unhandled_input(event: InputEvent) -> void:
    if input_enabled:
        if event.is_action_pressed("select_province") and not dragging:
            _select_province_at_mouse()

func _select_province_at_mouse() -> void:
    var global_mouse_pos = get_global_mouse_position()
    var provinces = get_tree().get_nodes_in_group("provinces")
    var clicked_province: Province = null
    
    for province in provinces:
        if province is not Province: continue
        var local_pos = province.to_local(global_mouse_pos)
        if Geometry2D.is_point_in_polygon(local_pos, province.collision_polygon.polygon):
            clicked_province = province as Province
            break
            
    var game: GameController = get_tree().current_scene as GameController
    if clicked_province != null:
        game.select_province(clicked_province)
    else:
        game.on_clear_province()
\end{lstlisting}

% =========================================================================
\section{The Provincial Subsystem}
\label{sec:provincial_subsystem}

Within the architectural design of the simulation, each discrete geopolitical territory is modeled through a specialized \texttt{Province} component. Rather than acting merely as visual containers, provinces serve as autonomous localized subsystems responsible for topological relationship discovery, dynamic boundary rendering, spatial graph connectivity, modular infrastructure management, and military force recruitment.

By centralizing these functions at the provincial level, the game engine enforces a clean separation of concerns: macro-level geopolitical shifts (such as conquests or diplomatic cessions) automatically cascade down into localized geometric and algorithmic state updates without requiring centralized, map-wide recalculations.

% -------------------------------------------------------------------------
\subsection{Dynamic Border Generation and Computational Topology}
\label{subsec:dynamic_borders}

A central visual requirement of grand strategy simulations is the clear demarcation of sovereign national borders. While internal provincial boundaries within the same nation are drawn with subtle, thin lines to preserve map readability, frontiers separating different sovereign states require pronounced, high-contrast visual demarcations. To achieve this in a dynamic environment where territorial control shifts in real time, the system implements an algorithmic border-generation pipeline that evaluates provincial geometry and dynamically constructs boundary lines.

The national border rendering process is encapsulated within the \texttt{set\_borders()} routine. This method is executed during initial scene composition and subsequently triggered whenever a province changes ownership. Upon invocation, the routine performs the following sequential operations:

\begin{enumerate}
    \item \textbf{Scene Tree Cleansing:} The method sweeps the province's children, removing any existing \texttt{Line2D} instances previously generated for borders. It explicitly ignores core static child nodes, such as the base \texttt{ProvinceBorder} or inter-provincial transit lines (\texttt{PathLine}), ensuring that only obsolete dynamic frontiers are cleared from memory.
    \item \textbf{Foreign Neighbor Filtering:} The province iterates across its cached collection of adjacent territories stored in the \texttt{neighbors} array. Adjacent provinces sharing the same parent node (i.e., belonging to the same sovereign nation) are skipped, isolating only foreign-controlled neighbors.
    \item \textbf{Coordinate Space Normalization:} Provincial collision geometry in Godot is stored locally within each node's coordinate space. Because two adjacent provinces possess distinct local origin vectors, their polygon vertices cannot be compared directly. To resolve this geometric discrepancy, the vertices of the foreign neighbor's \texttt{CollisionPolygon2D} are transformed through a compound coordinate mapping:
    \begin{equation}
        \mathbf{P}_{\text{local}} = \mathbf{T}_{\text{calling}}^{-1} \Big( \mathbf{T}_{\text{neighbor}} ( \mathbf{P}_{\text{neighbor}} ) \Big)
    \end{equation}
    where $\mathbf{T}_{\text{neighbor}}$ translates the neighbor's local vertices into global world coordinates, and $\mathbf{T}_{\text{calling}}^{-1}$ maps those global coordinates directly into the calling province's local space.
    \item \textbf{Shared Edge Collinearity Detection:} Once both geometric polygons are aligned within the identical reference frame, the helper function \texttt{get\_shared\_edges()} identifies overlapping perimeter segments.
    \item \textbf{Line Instantiation and Rendering:} For each shared edge detected, a new \texttt{Line2D} node is dynamically instantiated, configured with a thickened stroke ($\text{width} = 1.3\text{ px}$), a semi-transparent white tint, and smooth rounded joints (\texttt{LINE\_JOINT\_ROUND}). The node is attached to the scene tree and appended to the active \texttt{borders} tracking array.
\end{enumerate}

\begin{lstlisting}[language=python, caption={Dynamic border generation and shared edge extraction routines in GDScript.}, label={lst:dynamic_borders}]
func set_borders() -> void:
    borders.clear()
    for i in get_children():
        if i is not Line2D or i.name == "ProvinceBorder" or i.name.begins_with("PathLine"): 
            continue
        remove_child(i)

    for j in neighbors:
        if j.get_parent() == get_parent(): 
            continue
        var poly1 := PackedVector2Array()
        for p in j.get_node("CollisionPolygon2D").polygon:
            poly1.append(to_local(j.get_node("CollisionPolygon2D").to_global(p)))

        var poly2 := PackedVector2Array()
        for p in collision_polygon.polygon:
            poly2.append(p)

        var inter = get_shared_edges(poly1, poly2)

        var l = Line2D.new()
        l.width = 1.3
        l.default_color = Color(1.0, 1.0, 1.0, 0.784)
        l.joint_mode = Line2D.LINE_JOINT_ROUND
        for polygon_ in inter:
            for p in polygon_:
                l.add_point(p)
        add_child(l)
        borders.append(l)

func get_shared_edges(poly1: PackedVector2Array, poly2: PackedVector2Array, epsilon: float = 0.1) -> Array:
    var shared := []

    for i in range(poly1.size()):
        var a1 = poly1[i]
        var a2 = poly1[(i + 1) % poly1.size()]

        for j in range(poly2.size()):
            var b1 = poly2[j]
            var b2 = poly2[(j + 1) % poly2.size()]

            if a1.distance_to(b1) < epsilon and a2.distance_to(b2) < epsilon:
                shared.append([a1, a2])
            elif a1.distance_to(b2) < epsilon and a2.distance_to(b1) < epsilon:
                shared.append([a1, a2])

    return shared
\end{lstlisting}

The collinear edge extraction logic implemented in \texttt{get\_shared\_edges()} accounts for winding order variations across adjacent polygons. In computational geometry, two adjacent polygons sharing a common edge typically define that edge in opposite vertex orders (clockwise versus counter-clockwise). The algorithm iterates through every segment pair $(a_1, a_2)$ from the first polygon and $(b_1, b_2)$ from the second polygon. An intersection is identified if either:
\begin{equation}
    \big(\|a_1 - b_1\| < \epsilon \land \|a_2 - b_2\| < \epsilon\big) \quad \lor \quad \big(\|a_1 - b_2\| < \epsilon \land \|a_2 - b_1\| < \epsilon\big)
\end{equation}
where $\epsilon = 0.1\text{ px}$ acts as an error tolerance threshold to compensate for floating-point inaccuracies during vector transformations.

% -------------------------------------------------------------------------
\subsection{Asynchronous Topological Neighbor Discovery}
\label{subsec:neighbor_discovery}

To eliminate manual hardcoding of province adjacency matrices and enable modular map editing, the simulation incorporates an automated neighbor discovery routine, \texttt{\_get\_neighbors()}. This method leverages Godot's built-in 2D physics engine to evaluate spatial adjacency at startup.

Because adjacent provincial polygons touch along their borders without naturally penetrating one another, standard unscaled physics queries often fail to report bounding intersections. To resolve this, the discovery algorithm temporarily scales the province's collision polygon outward relative to its geometric center:

\begin{enumerate}
    \item \textbf{Centroid Calculation:} The arithmetic centroid $\mathbf{C}$ of the polygon is computed by averaging the coordinates of all its $N$ vertices:
    \begin{equation}
        \mathbf{C} = \frac{1}{N} \sum_{k=1}^{N} \mathbf{P}_k
    \end{equation}
    \item \textbf{Centroid-Relative Scaling:} Each vertex $\mathbf{P}_k$ is projected away from the centroid along its directional displacement vector scaled by a factor $s = 1.1$:
    \begin{equation}
        \mathbf{P}'_k = \mathbf{C} + (\mathbf{P}_k - \mathbf{C}) \times s
    \end{equation}
    \item \textbf{Physics Server Synchronization:} The expanded polygon is assigned to the \texttt{CollisionPolygon2D}. Because Godot processes spatial collisions across discrete physics steps, the execution is paused asynchronously across two consecutive physics ticks using \texttt{await get\_tree().physics\_frame}. This ensures the 2D physics server fully updates its broad-phase collision grid.
    \item \textbf{Collision Extraction and Reversion:} Overlapping \texttt{Area2D} nodes intersecting the inflated volume are retrieved via \texttt{get\_overlapping\_areas()} and added to the neighbor collection. Finally, the collision polygon is restored to its original unscaled geometry and redrawn.
\end{enumerate}

\begin{lstlisting}[language=python, caption={Asynchronous topological neighbor discovery via centroid-based polygon expansion.}, label={lst:neighbor_discovery}]
func _get_neighbors(scale_factor: float) -> Array[Area2D]:
    var polygon: PackedVector2Array = collision_polygon.polygon

    var center = Vector2.ZERO
    for point in polygon:
        center += point
    center /= polygon.size()

    var scaled_polygon = PackedVector2Array()
    for point in polygon:
        var direction = point - center
        scaled_polygon.append(center + direction * scale_factor)

    collision_polygon.polygon = scaled_polygon
    await get_tree().physics_frame
    await get_tree().physics_frame

    var res: Array[Area2D] = []
    for collision in get_overlapping_areas():
        res.append(collision)

    collision_polygon.polygon = polygon
    collision_polygon.queue_redraw()
    await get_tree().process_frame

    return res
\end{lstlisting}

% -------------------------------------------------------------------------
\subsection{Modular Infrastructure and Unit Mobilization Lifecycle}
\label{subsec:unit_recruitment}

Beyond topological management, each province acts as an infrastructure-dependent mobilization center for military divisions. Unit generation operates on an asynchronous timer loop governed by an internal \texttt{Timer} node (\texttt{army\_timer}). Rather than assuming unconditional continuous recruitment across all territories, mobilization speed and active execution are strictly gated by the presence and operational state of local military infrastructure.

\subsubsection{Infrastructure-Driven Mobilization Timing}
The total recruitment percentage is calculated by querying all operational infrastructure via the helper method \texttt{get\_total\_troop\_recruit\_percentage()}. This method iterates through the province's \texttt{buildings} collection, aggregating the \texttt{recruit\_increase\_percentage} exclusively from structures whose operational state evaluates to \texttt{FUNCTIONING}:
\begin{equation}
    P_{\text{recruit}} = \sum_{B \in \text{Buildings}} \Big( \text{Boost}(B) \cdot \mathbb{I}_{\text{functioning}}(B) \Big)
\end{equation}

The active recruitment duration $t_{\text{wait}}$ is computed dynamically in \texttt{\_calculate\_army\_generation\_wait()} by combining infrastructural speed boosts with provincial morale $M \in [0, 100]$:
\begin{equation}
    t_{\text{wait}} = \left(t_{\text{default}} \times \frac{P_{\text{recruit}}}{100.0}\right) + t_{\text{default}} \times \left(1.0 - \frac{M}{100.0}\right)
\end{equation}
where $t_{\text{default}} = 720.0\text{ s}$ denotes the baseline interval constant. If a province lacks operational recruitment infrastructure ($P_{\text{recruit}} = 0$), the recruitment subsystem automatically enters an idle state.

During the periodic \texttt{\_on\_global\_timer\_timeout()} audit, the system inspects $P_{\text{recruit}}$. If $P_{\text{recruit}} = 0$ and the recruitment timer is currently running, \texttt{army\_timer.stop()} is invoked. Conversely, if operational recruitment facilities are detected ($P_{\text{recruit}} > 0$) while the timer is halted, \texttt{army\_timer.start()} is triggered.

\subsubsection{Dynamic Timer Recalculation}
When provincial morale or infrastructure statuses shift mid-cycle, the active recruitment countdown must update without losing current execution progress. The \texttt{\_update\_army\_timer()} routine computes the active cycle's progress ratio:
\begin{equation}
    r = 1.0 - \frac{t_{\text{remaining}}}{t_{\text{current\_wait}}}
\end{equation}
The timer is then assigned the newly evaluated duration $t_{\text{new}}$ and restarted with a scaled remaining interval:
\begin{equation}
    t_{\text{remaining}}' = \max\Big(0.1, t_{\text{new}} \times (1.0 - r)\Big)
\end{equation}

\subsubsection{Recruitment Execution and BDI State Synchronization}
Upon timer completion, the signal handler \texttt{\_on\_army\_timer\_timeout()} executes the mobilization logic:

\begin{itemize}
    \item \textbf{Stationed Division Reinforcement:} The system checks whether a friendly division is currently stationed at the provincial capital coordinates. If present, its soldier count is incremented by one unit, and its hit points are scaled proportionally to provincial morale:
    \begin{equation}
        \Delta \text{HP} = \text{HP}_{\text{unit}} \times \frac{M}{100.0}
    \end{equation}
    Following this state modification, the updated division data is immediately transmitted to the external BDI reasoning engine via VEsNA through \texttt{dispatch\_my\_army\_update(el)}.
    \item \textbf{New Division Commissioning:} If the provincial capital is unoccupied, the province requests the central \texttt{GameController} to instantiate a new \texttt{Army} division at its capital node initialized with a single soldier.
\end{itemize}

\begin{lstlisting}[language=python, caption={Modular recruitment speed calculation and timer lifecycle routines in GDScript.}, label={lst:unit_recruitment}]
func get_total_troop_recruit_percentage() -> int:
    var total_recruit_percentage = 0
    for building in buildings:
        if buildings[building].state == Enums.BuildingStates.FUNCTIONING:
            total_recruit_percentage += buildings[building].recruit_increase_percentage
    return total_recruit_percentage

func _calculate_army_generation_wait() -> float:
    var total_recruit_percentage = get_total_troop_recruit_percentage()
    return (DEFAULT_ARMY_GENERATION_WAIT * float(total_recruit_percentage / 100.0)) + DEFAULT_ARMY_GENERATION_WAIT * (1.0 - (float(province_morale) / 100.0))

func _update_army_timer(val: float) -> void:
    if army_timer == null: return
    current_army_timer = val

    if army_timer.is_stopped():
        army_timer.wait_time = val
    else:
        var progress_ratio: float = 0.0
        if army_timer.wait_time > 0:
            progress_ratio = 1.0 - (army_timer.time_left / army_timer.wait_time)
        army_timer.wait_time = val
        var remaining_time: float = max(0.1, val * (1.0 - progress_ratio))
        army_timer.start(remaining_time)

func _on_global_timer_timeout() -> void:
    target_morale = min(100, 100 - _get_total_morale_penalties() + _get_total_morale_bonuses())

    var recruit_percentage = get_total_troop_recruit_percentage()
    if recruit_percentage == 0 and not army_timer.is_stopped():
        army_timer.stop()
    elif recruit_percentage > 0 and army_timer.is_stopped():
        army_timer.start()

func _on_army_timer_timeout() -> void:
    current_army_timer = _calculate_army_generation_wait()
    army_timer.wait_time = current_army_timer
    for el: Army in get_parent().armies:
        if el.global_position != to_global(get_city_position()): continue
        el.soldier_number += 1
        el.life_points += el.LIFE_PER_UNIT * (float(province_morale) / 100.0)
        el.queue_redraw()
        var agent = get_parent()
        if agent is Agent:
            agent.dispatch_my_army_update(el)
        return

    var game: GameController = get_tree().current_scene as GameController
    game.add_army_to_province(self, null, 1)
\end{lstlisting}

% -------------------------------------------------------------------------
\subsection{Resource Generation and Economic Production}
\label{subsec:resource_generation}

In parallel with unit recruitment, each province acts as an economic extraction center that fuels the nation's broader macro-management capabilities. Commodity and treasury extraction rates are calculated daily through two specialized methods: \texttt{get\_daily\_resource\_count()} and \texttt{get\_daily\_gold\_count()}. These production routines scale deterministically as a function of the province's current morale $M \in [0, 100]$ and its double-yield status $\delta_{\text{double}} \in \{1, 2\}$.

The mathematical formulations governing daily resource and gold yields are defined as follows:
\begin{align}
    \text{Yield}_{\text{commodity}} &= \big((M \times 24) + 750\big) \times \delta_{\text{double}} \\
    \text{Yield}_{\text{gold}} &= \begin{cases} 
        M \times 29, & \text{if } \delta_{\text{double}} = 2 \\ 
        M \times 57, & \text{if } \delta_{\text{double}} = 1 
    \end{cases}
\end{align}

Under this formulation, commodity output scales linearly with local stability while maintaining a baseline production threshold of $750$ units, ensuring that even severely destabilized provinces provide a minimal trickle of raw materials. Conversely, gold revenue scales purely in proportion to morale, rewarding players for domestic stability. In double-material provinces ($\delta_{\text{double}} = 2$), commodity extraction is doubled, while baseline gold revenue is halved to reflect industrial specialization over commercial taxation.

\begin{lstlisting}[language=python, caption={Resource and treasury daily extraction functions in GDScript.}, label={lst:resource_generation}]
func get_daily_resource_count() -> int:
    var resource_to_add: int = (province_morale * 24) + 750

    if double_material:
        resource_to_add *= 2
    return resource_to_add

func get_daily_gold_count() -> int:
    if double_material:
        return province_morale * 29
    else:
        return province_morale * 57
\end{lstlisting}

% -------------------------------------------------------------------------
\subsection{The Socio-Political Morale Subsystem}
\label{subsec:morale_subsystem}

A core determinant of provincial performance is its morale value $M$, which dynamically governs unit mobilization rates, economic yields, and civil rebellion risks. Rather than applying static values, the province continuously computes an ideal target morale equilibrium ($M_{\text{target}} \in [0, 100]$) toward which the active morale converges over successive simulation days:
\begin{equation}
    M_{\text{target}} = \min\Big(100, \max\big(0, 100 - \sum \text{Penalties} + \sum \text{Bonuses}\big)\Big)
\end{equation}

\subsubsection{Object-Oriented Modifier Architecture and Infrastructure Modularity}
To ensure structural modularity and clean separation of concerns, the simulation adopts an object-oriented Strategy Pattern where external logistical and geopolitical strains are encapsulated within independent classes inheriting from the abstract base class \texttt{MoraleModifier}. This class specifies baseline attributes (\texttt{MAX\_CAP}, \texttt{MIN\_CAP}, and \texttt{DEFAULT\_MODIFIER}) alongside a standardized polymorphic interface:
\begin{itemize}
    \item \textbf{\texttt{apply\_modifier(province)}:} Computes and returns the signed integer penalty value applicable to the given province.
    \item \textbf{\texttt{is\_modifier\_applied(province)}:} Evaluates whether the underlying preconditions for modifier activation are met.
\end{itemize}

\begin{lstlisting}[language=python, caption={Abstract MoraleModifier interface in GDScript.}, label={lst:morale_modifier_base}]
extends Object
class_name MoraleModifier

var MAX_CAP: Variant = null
var MIN_CAP: Variant = null
var DEFAULT_MODIFIER: Variant = null

func apply_modifier(province: Province) -> int:
    push_error("apply_modifier must be implemented")
    return -1

func is_modifier_applied(province: Province) -> bool:
    push_error("is_modifier_applied must be implemented")
    return false
\end{lstlisting}

In contrast to external geopolitical maluses, positive morale bonuses are decoupled from static modifier classes and linked directly to operational \texttt{Building} components. This architectural separation ensures that civil infrastructure improvements dynamically register their positive morale contributions without requiring dedicated wrapper classes.

\subsubsection{Concrete Geopolitical Maluses}
Drawing inspiration from the grand strategy mechanics of \textit{Supremacy 1914}, four concrete malus classes were implemented to reflect administrative, martial, and expansionist strains:

\begin{enumerate}
    \item \textbf{Capital Distance Penalty (\texttt{CapitalDistanceModifier}):} Models administrative friction and logistical decay away from the center of governance. If the sovereign nation possesses no operational capital building, a severe flat penalty of $-40\%$ is imposed. When a capital exists at Euclidean distance $d_{\text{cap}} = \|\mathbf{P}_{\text{city}} - \mathbf{P}_{\text{capital}}\|$, the penalty scales linearly up to a maximum cap of $-35\%$:
    \begin{equation}
        \text{Penalty}_{\text{dist}} = -\min\left(35, \left\lfloor \frac{d_{\text{cap}}}{2000.0} \times 35 \right\rfloor\right)
    \end{equation}

    \item \textbf{Enemy Occupation Penalty (\texttt{EnemyArmiesModifier}):} Simulates civilian terror and martial disruption caused by hostile troops within provincial borders. The penalty is directly proportional to the total number of foreign combatants $N_{\text{hostiles}}$ belonging to nations with which the owner is actively at war, capped at $-30\%$:
    \begin{equation}
        \text{Penalty}_{\text{armies}} = -\min(30, N_{\text{hostiles}})
    \end{equation}

    \item \textbf{Hostile Frontier Penalty (\texttt{EnemyNeighborsModifier}):} Reflects psychological vulnerability along contested borders. A flat penalty of $-5\%$ is levied for each adjacent border province $P_{\text{neighbor}} \in \text{Neighbors}$ controlled by an active enemy nation:
    \begin{equation}
        \text{Penalty}_{\text{neighbors}} = -5 \times \sum_{P \in \text{Neighbors}} \mathbb{I}_{\text{enemy}}(P)
    \end{equation}

    \item \textbf{Territorial Overextension Penalty (\texttt{ExpansionModifier}):} Implements anti-snowballing balance mechanics by penalizing rapid overexpansion. Let $r = \frac{N_{\text{owned}}}{N_{\text{global}}}$ denote the ratio of owned territories relative to the global map. For $r \le 0.06$ (early game footprint), no penalty is incurred; for $r \ge 0.50$ (hegemonic scale), the maximum penalty of $-35\%$ is enforced. In intermediate ranges, the penalty interpolates linearly:
    \begin{equation}
        \text{Penalty}_{\text{exp}} = \begin{cases}
            0, & \text{if } r \le 0.06 \\
            -35, & \text{if } r \ge 0.50 \\
            -\left\lfloor 35 \times \frac{r - 0.06}{0.50 - 0.06} \right\rfloor, & \text{otherwise}
        \end{cases}
    \end{equation}
\end{enumerate}

\begin{lstlisting}[language=python, caption={Concrete MoraleModifier implementations for capital distance, enemy presence, hostile frontiers, and expansion penalties.}, label={lst:concrete_morale_modifiers}]
class_name CapitalDistanceModifier
extends MoraleModifier

const NO_CAPITAL_PENALTY = 40
const MAX_DISTANCE = 2000.0

func _init() -> void:
    MAX_CAP = 35

func apply_modifier(province: Province) -> int:
    var parent = province.get_parent()
    if not parent is Agent: return 0

    var capital: Province = null
    for child in parent.get_children():
        if child is Province and Enums.Buildings.CAPITAL in child.buildings:
            capital = child
            break

    if capital == null:
        return -NO_CAPITAL_PENALTY

    var dist = province.to_global(province.get_city_position()).distance_to(capital.to_global(capital.get_city_position()))
    var penalty = int((dist / MAX_DISTANCE) * MAX_CAP)
    return -min(penalty, MAX_CAP)

func is_modifier_applied(province: Province) -> bool:
    return apply_modifier(province) != 0


class_name EnemyArmiesModifier
extends MoraleModifier

func _init() -> void:
    MAX_CAP = 30

func apply_modifier(province: Province) -> int:
    var parent = province.get_parent()
    if not parent is Agent: return 0

    var enemy_soldiers = 0
    var players = province.get_tree().get_nodes_in_group("players")
    for p in players:
        if p is Agent and p in parent.player_in_war:
            for army in p.armies:
                if is_instance_valid(army) and army.province_location == province:
                    enemy_soldiers += army.soldier_number

    return -min(MAX_CAP, enemy_soldiers)

func is_modifier_applied(province: Province) -> bool:
    return apply_modifier(province) != 0


class_name EnemyNeighborsModifier
extends MoraleModifier

func _init() -> void:
    DEFAULT_MODIFIER = 5

func apply_modifier(province: Province) -> int:
    var parent = province.get_parent()
    if not parent is Agent: return 0

    var enemy_neighbors_count = 0
    for neighbor in province.neighbors:
        if neighbor is Province:
            var neighbor_owner = neighbor.get_parent()
            if neighbor_owner is Agent and neighbor_owner in parent.player_in_war:
                enemy_neighbors_count += 1

    return -DEFAULT_MODIFIER * enemy_neighbors_count

func is_modifier_applied(province: Province) -> bool:
    return apply_modifier(province) != 0


class_name ExpansionModifier
extends MoraleModifier

const MIN_EXPANSION = 0.06
const MAX_EXPANSION = 0.5

func _init() -> void:
    MAX_CAP = 35

func apply_modifier(province: Province) -> int:
    var parent = province.get_parent()
    if not parent is Agent: return 0

    var all_provinces = province.get_tree().get_nodes_in_group("provinces")
    var total_count = all_provinces.size()
    if total_count == 0: return 0

    var owned_count = 0
    for child in parent.get_children():
        if child is Province:
            owned_count += 1

    var ratio = float(owned_count) / float(total_count)

    if ratio <= MIN_EXPANSION:
        return 0
    elif ratio >= MAX_EXPANSION:
        return -MAX_CAP
    else:
        var penalty = int(float(MAX_CAP) * (ratio - MIN_EXPANSION) / (MAX_EXPANSION - MIN_EXPANSION))
        return -penalty

func is_modifier_applied(province: Province) -> bool:
    return apply_modifier(province) != 0
\end{lstlisting}

\subsubsection{Dynamic Modifier Aggregation and Building Bonus Evaluation}
Within each \texttt{Province} instance, external maluses are registered in a typed lookup dictionary (\texttt{all\_maluses}) mapped via \texttt{ProvinceMoraleMaluses} enumerations.

\begin{lstlisting}[language=python, caption={Dictionary instantiations and enumerations for provincial morale maluses.}, label={lst:modifier_dictionaries}]
var all_bonuses: Dictionary = {}
var all_maluses: Dictionary = {
    Enums.ProvinceMoraleMaluses.CAPITAL_DISTANCE: CapitalDistanceModifier.new(),
    Enums.ProvinceMoraleMaluses.EXPANSION: ExpansionModifier.new(),
    Enums.ProvinceMoraleMaluses.ENEMY_NEIGHBORS: EnemyNeighborsModifier.new(),
    Enums.ProvinceMoraleMaluses.ENEMY_ARMIES: EnemyArmiesModifier.new(),
}

enum ProvinceMoraleBonuses {
    CAPITAL,
}

enum ProvinceMoraleMaluses {
    CAPITAL_DISTANCE,
    EXPANSION,
    ENEMY_NEIGHBORS,
    ENEMY_ARMIES,
}
\end{lstlisting}

During active evaluation passes, the method \texttt{\_update\_active\_morale\_modifiers()} synchronizes the active modifier pools:
\begin{enumerate}
    \item It evaluates all registered \texttt{all\_maluses}, verifying precondition validity via \texttt{is\_modifier\_applied(self)} before inserting them into \texttt{morale\_maluses}.
    \item It iterates across the \texttt{buildings} collection, dynamically registering structures that declare a non-zero \texttt{morale\_increase\_percentage} and satisfy the \texttt{FUNCTIONING} operational state check into \texttt{morale\_bonuses} under a synthetic \texttt{"building\_\%s"} key format.
\end{enumerate}

The aggregated score is subsequently resolved through \texttt{\_get\_total\_morale\_penalties()} and \texttt{\_get\_total\_morale\_bonuses()}. Specifically, \texttt{\_get\_total\_morale\_bonuses()} distinguishes between custom modifier strategies and infrastructure contributions, directly adding the building's \texttt{morale\_increase\_percentage} to the total positive offset.

\begin{lstlisting}[language=python, caption={Active modifier resolution and building bonus aggregation routines in GDScript.}, label={lst:modifier_aggregation}]
func _update_active_morale_modifiers() -> void:
    morale_bonuses.clear()
    morale_maluses.clear()
    for key in all_bonuses:
        var modifier = all_bonuses[key]
        if modifier.is_modifier_applied(self):
            morale_bonuses[key] = modifier
    for building in buildings:
        if buildings[building].morale_increase_percentage != 0 and buildings[building].state == Enums.BuildingStates.FUNCTIONING:
            morale_bonuses["building_%s" % building] = buildings[building]
    for key in all_maluses:
        var modifier = all_maluses[key]
        if modifier.is_modifier_applied(self):
            morale_maluses[key] = modifier

func _get_total_morale_penalties() -> int:
    _update_active_morale_modifiers()
    var total_penalties = 0
    for key in morale_maluses:
        var modifier = morale_maluses[key]
        total_penalties += abs(modifier.apply_modifier(self))
    return total_penalties

func _get_total_morale_bonuses() -> int:
    _update_active_morale_modifiers()
    var total_bonuses = 0
    for key in morale_bonuses:
        var modifier = morale_bonuses[key]
        if key.begins_with("building_"):
            total_bonuses += modifier.morale_increase_percentage
        else:
            total_bonuses += abs(modifier.apply_modifier(self))
    return total_bonuses
\end{lstlisting}

% =========================================================================
\subsection{Infrastructure and Modular Building Subsystem}
\label{subsec:building_subsystem}

Within each province, territorial advancement and socio-economic progression are regulated by a modular infrastructure system. In grand strategy simulations such as \textit{Supremacy 1914}, provincial infrastructure serves as the primary strategic lever for long-term domestic advancement. Players construct specialized buildings—such as Recruiting Centers, Barracks, Workshops, Fortresses, Harbors, and Railways—to unlock advanced mechanized divisions, accelerate troop recruitment, reinforce defensive combat absorption, augment raw resource extraction yields, and accelerate movement velocity across logistics corridors.

To incorporate these multi-dimensional effects within the BDI experimental framework while maintaining code modularity, the simulation models infrastructure via an object-oriented Strategy Pattern anchored by the abstract base class \texttt{Building}. This class defines the standardized state machine, lifecycle methods, construction costs, and operational attributes common to all infrastructure components.

\subsubsection{The Abstract \texttt{Building} Base Specification}
The lifecycle state of any building instance is governed by the strongly typed enumeration \texttt{BuildingStates}:

\begin{lstlisting}[language=python, caption={Building state enumeration in GDScript.}, label={lst:building_states_enum}]
enum BuildingStates {
    FUNCTIONING,
    DAMAGED,
    NOT_BUILT,
}
\end{lstlisting}

Every concrete subclass inheriting from \texttt{Building} encapsulates a comprehensive set of structural, economic, and military attributes:
\begin{itemize}
    \item \textbf{\texttt{state} (\texttt{BuildingStates}):} Tracks the operational state of the structure.
    \item \textbf{\texttt{max\_lp} and \texttt{lp} (\texttt{int}):} Represent the maximum and current hit points of the building, allowing structural damage to be modeled during hostile artillery bombardments or post-conquest attrition.
    \item \textbf{\texttt{morale\_increase\_percentage} (\texttt{int}):} Positive percentage offset contributed directly to the province's target morale equilibrium ($M_{\text{target}}$).
    \item \textbf{\texttt{recruit\_increase\_percentage} (\texttt{int}):} Percentage acceleration applied to provincial recruitment countdowns.
    \item \textbf{\texttt{damage\_absorbtion\_percentage} (\texttt{float}):} Combat defensive mitigation provided to stationed divisions during hostile assaults.
    \item \textbf{\texttt{resource\_production\_increase\_percentage} (\texttt{float}):} Percentage multiplier applied to daily provincial raw commodity extraction.
    \item \textbf{\texttt{troop\_speed\_increase\_percentage} (\texttt{float}):} Traversal velocity bonus applied to friendly divisions moving through the province's transit paths.
    \item \textbf{\texttt{onboarding\_time\_reduce\_percentage} (\texttt{float}):} Embarkation latency reduction for amphibious or maritime transit operations.
    \item \textbf{\texttt{build\_cost} and \texttt{daily\_cost} (\texttt{Dictionary}):} Multi-commodity upfront construction costs and ongoing hourly maintenance deductions.
    \item \textbf{\texttt{build\_time} (\texttt{float}):} Temporal duration (in minutes) required to complete construction.
\end{itemize}

The \texttt{Building} base class exposes helper routines including \texttt{has\_enough\_resources(player)}, which verifies whether a nation holds sufficient raw materials and treasury reserves to commission the building, alongside structural mutation methods: \texttt{set\_state(new\_state)}, \texttt{add\_lp(delta)}, \texttt{set\_lp(val)}, and \texttt{set\_max\_lp()}.

\subsubsection{Concrete Building Implementations}
While advanced operational parameters are architecturally established to facilitate future expansion (such as naval docks and logistical railways), the core simulation environment implements two foundational structures:
\begin{enumerate}
    \item \textbf{\texttt{Capital}:} Represents the administrative, political, and cultural seat of a sovereign nation. Each nation is restricted to possessing exactly one active capital building at any time. When operational, the capital anchors the distance metric used by all other provinces to compute the \texttt{CapitalDistanceModifier} penalty. Furthermore, its presence triggers custom UI containers (\texttt{province\_name\_capital\_container}) that display highlighted national capital labels.
    \item \textbf{\texttt{RecruitCenter}:} Represents military enlistment facilities. Unit recruitment is strictly contingent upon possessing a functioning recruitment center. The structure provides a baseline \texttt{recruit\_increase\_percentage}, ensuring that \texttt{get\_total\_troop\_recruit\_percentage()} evaluates above zero and permits the province's \texttt{army\_timer} to remain active.
\end{enumerate}

% =========================================================================
\subsection{Hierarchical Multi-Scale Clock Architecture}
\label{subsec:hierarchical_timers}

To simulate the passage of time realistically while preserving computational efficiency, the provincial execution loop connects to three global clock nodes: \texttt{GlobalTimer}, \texttt{GlobalHourTimer}, and \texttt{GlobalDayTimer}. 

This hierarchical separation decouples high-frequency state polling from heavy periodic computations (such as daily economic yields and civil rebellion checks). Furthermore, it allows the entire simulation to dynamically scale its playback speed (e.g., accelerating training runs for BDI agents) while maintaining invariant proportionality between ticks, simulation hours, and simulation days.

\subsubsection{1. Short-Tick Loop: \texttt{GlobalTimer}}
The \texttt{GlobalTimer} executes short-interval periodic audits. In the provincial scope, its primary responsibility is auditing infrastructure integrity and synchronizing the recruitment timer state:
\begin{enumerate}
    \item \textbf{Target Morale Recalculation:} It re-evaluates the province's equilibrium target morale:
    \begin{equation}
        M_{\text{target}} = \min\Big(100, \max\big(0, 100 - \sum \text{Penalties} + \sum \text{Bonuses}\big)\Big)
    \end{equation}
    where $\sum \text{Bonuses}$ dynamically aggregates positive percentage offsets from all functioning buildings.
    \item \textbf{Recruitment Gating:} It queries \texttt{get\_total\_troop\_recruit\_percentage()}. If no functioning recruitment infrastructure exists ($P_{\text{recruit}} = 0$) and the \texttt{army\_timer} is active, it halts recruitment via \texttt{army\_timer.stop()}. If recruitment infrastructure is present ($P_{\text{recruit}} > 0$) and the timer was halted, \texttt{army\_timer.start()} is invoked.
\end{enumerate}

\subsubsection{2. Hourly Economic Loop: \texttt{GlobalHourTimer}}
The \texttt{GlobalHourTimer} manages continuous economic accounting. Upon each hourly tick, the province calculates and dispatches fractional economic yields and upkeep deductions directly to the owning agent's treasury:
\begin{align}
    \Delta \text{Resource} &= \left\lfloor \frac{\text{Yield}_{\text{commodity}}}{24.0} \right\rfloor \\
    \Delta \text{Gold} &= \left\lfloor \frac{\text{Yield}_{\text{gold}}}{24.0} \right\rfloor \\
    \Delta \text{Upkeep}(R) &= -\left\lfloor \frac{\text{Cost}(R)}{24.0} \right\rfloor, \quad \forall R \in \text{Resources}
\end{align}
where $\text{Cost}(R)$ is defined by \texttt{DEFAULT\_DAILY\_COST}, establishing baseline daily upkeep expenditures across all strategic commodities:

\begin{lstlisting}[language=python, caption={Default daily provincial maintenance costs across strategic commodities.}, label={lst:daily_cost_dict}]
const DEFAULT_DAILY_COST: Dictionary[Enums.Resources, int] = {
    Enums.Resources.WOOD: 400,
    Enums.Resources.IRON: 400,
    Enums.Resources.OIL: 267,
    Enums.Resources.GAS: 267,
    Enums.Resources.COAL: 267,
    Enums.Resources.GRAIN: 400,
    Enums.Resources.FISH: 400
}
\end{lstlisting}

\subsubsection{3. Daily Progression Loop: \texttt{GlobalDayTimer}}
The \texttt{GlobalDayTimer} executes macro-level geopolitical transitions, gradual morale convergence, telemetry synchronization, and civil rebellion evaluations.

\paragraph{Morale Convergence Dynamics}
Rather than transitioning instantaneously to $M_{\text{target}}$, current provincial morale $M$ adjusts gradually each day by a bounded step function based on the delta $\Delta = M_{\text{target}} - M$:
\begin{equation}
    \Delta M_{\text{base}} = \left\lceil \frac{\Delta}{100.0} \times 15.0 \right\rceil
\end{equation}
If $\Delta M_{\text{base}} < 3$, the morale increments directly by $\Delta M_{\text{base}}$. Otherwise, a stochastic perturbation $\epsilon_{\text{rng}} \in \{-3, \dots, +3\}$ is sampled to simulate localized political noise:
\begin{equation}
    M_{t+1} = \min\big(100, M_t + \Delta M_{\text{base}} + \epsilon_{\text{rng}}\big)
\end{equation}

Following morale adjustment, \texttt{\_calculate\_army\_generation\_wait()} and \texttt{\_update\_army\_timer()} recompute recruitment durations, and an asynchronous telemetry frame is dispatched to VEsNA to update the BDI agent's Belief Base.

\paragraph{Civil Rebellion and Riot Mechanics}
If provincial morale falls below the critical riot threshold ($\theta_{\text{riot}} = 33\%$), the province conducts a garrison stability audit. Stationed military defense strength ($D_{\text{stat}}$) is evaluated against a required garrison threshold ($D_{\text{req}}$):
\begin{equation}
    D_{\text{stat}} = \sum_{A \in \text{StationedArmies}} \left( \text{Soldiers}(A) \times \text{DefUnit} \times \frac{\text{HP}(A)}{\text{Soldiers}(A) \times \text{HPUnit}} \right)
\end{equation}
\begin{equation}
    D_{\text{req}} = \begin{cases}
        5.0, & \text{if } 30 \le M < 33 \\
        5.0 + (30.0 - M) \times 6.0, & \text{if } M < 30
    \end{cases}
\end{equation}

If $D_{\text{stat}} < D_{\text{req}}$, a riot probability check is performed:
\begin{equation}
    P_{\text{riot}} = 1.0 - \frac{M}{33.0}
\end{equation}

If a sampled random scalar $r \sim \mathcal{U}(0, 1) < P_{\text{riot}}$, the method \texttt{\_trigger\_riot()} is invoked. The revolting province cedes territorial control to a random adjacent non-allied nation, causing immediate realignment of regional borders and topological graphs.

\begin{lstlisting}[language=python, caption={Global timer signal handlers managing multi-scale temporal execution in GDScript.}, label={lst:global_timer_handlers}]
func _on_global_timer_timeout() -> void:
    target_morale = min(100, 100 - _get_total_morale_penalties() + _get_total_morale_bonuses())

    var recruit_percentage = get_total_troop_recruit_percentage()
    if recruit_percentage == 0 and not army_timer.is_stopped():
        army_timer.stop()
    elif recruit_percentage > 0 and army_timer.is_stopped():
        army_timer.start()

func _on_global_hour_timer_timeout() -> void:
    var agent: Agent = get_parent()
    agent.add_resource(province_material, int(get_daily_resource_count() / 24.0))
    agent.add_resource(Enums.Resources.GOLD, int(get_daily_gold_count() / 24.0))
    for res in daily_cost:
        agent.add_resource(res, -int(daily_cost[res] / 24.0))
    agent.queue_redraw()

func _on_global_day_timer_timeout() -> void:
    if conquered:
        conquered = false
        return
    var delta: float = target_morale - province_morale
    if delta == 0: return
    var adding: int = ceil(delta / 100.0 * 15.0)
    if adding < 3:
        province_morale += adding
    else:
        var adding2: int = rng.randi_range(-3, 3)
        province_morale += adding + adding2
    province_morale = min(100, province_morale)
    current_army_timer = _calculate_army_generation_wait()
    _update_army_timer(current_army_timer)
    queue_redraw()

    var agent = get_parent()
    if agent is Agent:
        var msg: Dictionary = {
            "type": "sight",
            "sender": agent.name,
            "receiver": "vesna",
            "data": {
                "type": "province_morale_update",
                "id": name,
                "morale": province_morale
            }
        }
        agent.send_vesna_message(msg)

    # Civil Rebellion Evaluation
    if province_morale < RIOT_THRESHOLD:
        var stationary_defense = 0.0
        var player = get_parent()
        if player is Agent:
            for army in player.armies:
                if is_instance_valid(army) and not army.travelling and army.province_location == self:
                    if army.soldier_number > 0:
                        stationary_defense += (army.soldier_number * army.DEFENSE_PER_UNIT) * army.life_points / (army.soldier_number * army.LIFE_PER_UNIT)

        var required_defense = 0.0
        if province_morale >= RIOT_BREAKPOINT:
            required_defense = RIOT_BASE_DEFENSE
        else:
            required_defense = RIOT_BASE_DEFENSE + (float(RIOT_BREAKPOINT) - province_morale) * RIOT_PENALTY_MULTIPLIER

        if stationary_defense < required_defense:
            var riot_chance = 1.0 - (float(province_morale) / float(RIOT_THRESHOLD))
            if rng.randf() < riot_chance:
                _trigger_riot()

func _trigger_riot() -> void:
    var player = get_parent()
    if player is not Agent: return

    var possible_nations: Array[Agent] = []
    for neighbor in neighbors:
        if neighbor is Province:
            var neighbor_owner = neighbor.get_parent()
            if neighbor_owner is Agent and neighbor_owner != player and neighbor_owner not in possible_nations:
                possible_nations.append(neighbor_owner)

    if possible_nations.is_empty():
        var players = get_tree().get_nodes_in_group("players")
        for p in players:
            if p is Agent and p != player:
                possible_nations.append(p)

    if possible_nations.is_empty(): return

    var new_owner = possible_nations[rng.randi_range(0, possible_nations.size() - 1)]
    var game: GameController = get_tree().current_scene as GameController
    game.conquer_province(self, new_owner)
\end{lstlisting}

% =========================================================================
\subsection{Auxiliary Utilities, Rendering Pipelines, and Provincial Lifecycle}
\label{subsec:provincial_utilities_lifecycle}

To complement the core economic, topological, and infrastructural systems, the \texttt{Province} component incorporates a dedicated suite of visual rendering utilities, UI event handlers, graph-path query routines, and a unified initialization sequence executed within Godot's node lifecycle.

\subsubsection{Dynamic Texture Mapping and Graphical Selection}
Visual state feedback in the simulation environment must immediately reflect ownership, economic outputs, and spectator selection:

\begin{itemize}
    \item \textbf{Resource Texture Dispatching (\texttt{\_get\_sprite\_key()}):} To ensure high rendering performance without redundant node allocations, the provincial commodity display is resolved dynamically. The method \texttt{\_get\_sprite\_key()} inspects the exported attributes \texttt{province\_material} and \texttt{double\_material}, mapping them to the corresponding \texttt{Enums.ResourcesSprites} key. This key indexes into the preloaded dictionary \texttt{MATERIAL\_TEXTURES}, assigning the appropriate sprite texture directly to the \texttt{TextureRect} container above the provincial capital.
    \item \textbf{Alpha-Composited Province Highlighting (\texttt{toggle\_selection()}):} When an observer selects a province via mouse input, the system alters the fill color of the underlying \texttt{Polygon2D} node without overwriting the nation's identifying color palette. This is achieved via two-stage alpha-blending:
    \begin{align}
        \mathbf{C}_{\text{rgb}} &= \mathbf{C}_{\text{bottom}} \times (1.0 - \alpha_{\text{top}}) + \mathbf{C}_{\text{top}} \times \alpha_{\text{top}} \\
        \alpha_{\text{out}} &= \alpha_{\text{bottom}} + \alpha_{\text{top}} \times (1.0 - \alpha_{\text{bottom}})
    \end{align}
    where $\mathbf{C}_{\text{bottom}}$ represents the sovereign nation's primary RGB tint, and $\mathbf{C}_{\text{top}}$ is a semi-transparent black modulation layer ($\text{RGBA} = (0.0, 0.0, 0.0, 0.114)$). Deselecting the province restores $\mathbf{C}_{\text{bottom}}$, and \texttt{queue\_redraw()} is invoked to synchronize the visual canvas item.
\end{itemize}

\begin{lstlisting}[language=python, caption={Dynamic sprite texture resolution and alpha-composited selection blending routines in GDScript.}, label={lst:texture_and_selection}]
func _get_sprite_key() -> Enums.ResourcesSprites:
    match province_material:
        Enums.Resources.WOOD:
            return Enums.ResourcesSprites.WOOD_DOUBLE if double_material else Enums.ResourcesSprites.WOOD
        Enums.Resources.IRON:
            return Enums.ResourcesSprites.IRON_DOUBLE if double_material else Enums.ResourcesSprites.IRON
        Enums.Resources.OIL:
            return Enums.ResourcesSprites.OIL_DOUBLE if double_material else Enums.ResourcesSprites.OIL
        Enums.Resources.GAS:
            return Enums.ResourcesSprites.GAS_DOUBLE if double_material else Enums.ResourcesSprites.GAS
        Enums.Resources.COAL:
            return Enums.ResourcesSprites.COAL_DOUBLE if double_material else Enums.ResourcesSprites.COAL
        Enums.Resources.GRAIN:
            return Enums.ResourcesSprites.GRAIN_DOUBLE if double_material else Enums.ResourcesSprites.GRAIN
        Enums.Resources.FISH:
            return Enums.ResourcesSprites.FISH_DOUBLE if double_material else Enums.ResourcesSprites.FISH
        _:
            return Enums.ResourcesSprites.WOOD

func toggle_selection() -> bool:
    selected = not selected
    if selected:
        fill.color = blend(get_parent().color, Color(0.0, 0.0, 0.0, 0.114))
    else:
        fill.color = get_parent().color
    queue_redraw()
    return selected

func blend_under(bottom: Color, top: Color) -> Color:
    var a = top.a
    return bottom * (1.0 - a) + top * a

func blend(bottom: Color, top: Color) -> Color:
    var out = blend_under(bottom, top)
    out.a = bottom.a + top.a * (1.0 - bottom.a)
    return out
\end{lstlisting}

\subsubsection{Spatial Coordinates and Path Graph Extraction}
For military pathfinding and army navigation, the province exposes geometric interrogation methods:
\begin{itemize}
    \item \textbf{Provincial Seat Positioning:} The local coordinates of the urban center are retrieved via \texttt{get\_city\_position()}, while \texttt{get\_city\_global\_position()} returns its global vector in the simulation world space using \texttt{to\_global(city\_node.position)}.
    \item \textbf{Graph Path Edge Resolution (\texttt{get\_province\_paths()}):} Military units traverse along explicit \texttt{Curve2D} resources enclosed within child \texttt{Path2D} nodes. To determine available transit corridors, \texttt{get\_province\_paths()} scans all nodes registered in the \texttt{"paths"} group. It extracts any trajectory curve whose origin point $\mathbf{P}_0 = \text{curve.get\_point\_position}(0)$ or terminal point $\mathbf{P}_n = \text{curve.get\_point\_position}(N-1)$ coincides with the provincial city coordinates, pairing the path reference with its corresponding anchor point.
    \item \textbf{Match Termination Handling (\texttt{\_on\_match\_over()}):} When the global match concludes, the game controller emits a termination signal, prompting \texttt{\_on\_match\_over()} to halt \texttt{army\_timer} and freeze all military production.
\end{itemize}

\begin{lstlisting}[language=python, caption={Spatial positioning, path extraction, and match lifecycle handlers.}, label={lst:spatial_and_lifecycle_helpers}]
func get_city_position() -> Vector2:
    return city_node.position

func get_city_global_position() -> Vector2:
    return to_global(city_node.position)

func get_province_paths() -> Array:
    var global_paths = get_tree().get_nodes_in_group("paths")
    var paths = []
    for el in global_paths:
        if el is not Path2D: continue
        var curve: Curve2D = (el as Path2D).curve
        var start = curve.get_point_position(0)
        var end = curve.get_point_position(curve.point_count - 1)

        if start == get_city_position():
            paths.append([el, start])
        elif end == get_city_position():
            paths.append([el, end])
    return paths

func _on_match_over() -> void:
    army_timer.stop()
\end{lstlisting}

\subsubsection{Lifecycle Initialization and Scene Setup (\texttt{\_ready()})}
The entire configuration sequence of a province is consolidated within the \texttt{\_ready()} virtual method, executed automatically when the node enters the active scene tree:

\begin{enumerate}
    \item \textbf{Group and Signal Registration:} The province registers into the global \texttt{"provinces"} node group and connects the timeout signals of \texttt{GlobalTimer}, \texttt{GlobalHourTimer}, and \texttt{GlobalDayTimer} to their respective handlers. The \texttt{army\_timer} instance is added as an internal child node and bound to \texttt{\_on\_army\_timer\_timeout()}.
    \item \textbf{Visual Boundary and Collision Geometry Construction:} The collision vertices from \texttt{collision\_polygon.polygon} are copied to \texttt{fill.polygon} to generate the territorial background mesh. Concurrently, vertices are added point-by-point to the \texttt{province\_border} line renderer.
    \item \textbf{Control Layer Isolation:} The helper function \texttt{\_set\_mouse\_filter\_ignore\_recursive()} recursively traverses the child UI hierarchy under \texttt{city\_node}, applying \texttt{Control.MOUSE\_FILTER\_IGNORE}. This ensures mouse clicks propagate directly to provincial raycasting routines without being consumed by static UI labels.
    \item \textbf{Asynchronous Topology and Boundary Discovery:} The province awaits neighbor detection via \texttt{await \_get\_neighbors(1.1)} and triggers \texttt{set\_borders()} to construct foreign frontier lines.
    \item \textbf{Startup Building Instantiation:} Based on the exported \texttt{startup\_buildings} array, concrete building instances (\texttt{Capital} and \texttt{RecruitCenter}) are allocated, initialized to the \texttt{FUNCTIONING} state, and stored in the \texttt{buildings} dictionary.
    \item \textbf{Recruitment Initialization:} Recruitment interval latency is evaluated using \texttt{\_calculate\_army\_generation\_wait()}. If the computed duration is positive ($t_{\text{wait}} > 0$, indicating functioning recruitment infrastructure), \texttt{army\_timer.start()} is called.
    \item \textbf{Dynamic Transit Path Visualization:} The province iterates over all child \texttt{Path2D} components, programmatically instantiating a subtle, semi-transparent \texttt{Line2D} overlay ($\text{width} = 0.7\text{ px}$, $\text{RGBA} = (0.0, 0.0, 0.0, 0.1)$) along each curve trajectory to render roads connecting provincial seats on the global map.
\end{enumerate}

\begin{lstlisting}[language=python, caption={Complete provincial initialization sequence implemented in the \texttt{\_ready()} routine.}, label={lst:province_ready_full}]
func _ready() -> void:
    add_to_group("provinces")
    timer = get_node("/root/Game/GlobalTimer")
    timer.timeout.connect(_on_global_timer_timeout)
    hour_timer = get_node("/root/Game/GlobalHourTimer")
    hour_timer.timeout.connect(_on_global_hour_timer_timeout)
    day_timer = get_node("/root/Game/GlobalDayTimer")
    day_timer.timeout.connect(_on_global_day_timer_timeout)
    army_timer.timeout.connect(_on_army_timer_timeout)
    add_child(army_timer)

    fill.polygon = collision_polygon.polygon
    daily_cost = DEFAULT_DAILY_COST

    province_border.hide()
    control.hide()
    _set_mouse_filter_ignore_recursive(city_node)
    resource_texture.texture = MATERIAL_TEXTURES[_get_sprite_key()]

    for el in collision_polygon.polygon:
        province_border.add_point(el)

    var game: GameController = get_tree().current_scene as GameController
    game.match_over.connect(_on_match_over)

    neighbors = await _get_neighbors(1.1)
    set_borders()

    target_morale = min(100, 100 - _get_total_morale_penalties() + _get_total_morale_bonuses())

    # Startup Infrastructure Initialization
    if Enums.Buildings.CAPITAL in startup_buildings:
        var capital = Capital.new()
        capital.state = Enums.BuildingStates.FUNCTIONING
        buildings.get_or_add(Enums.Buildings.CAPITAL, capital)
    if Enums.Buildings.RECRUIT_CENTER in startup_buildings:
        var recruit_center = RecruitCenter.new()
        recruit_center.state = Enums.BuildingStates.FUNCTIONING
        buildings.get_or_add(Enums.Buildings.RECRUIT_CENTER, recruit_center)

    current_army_timer = _calculate_army_generation_wait()
    _update_army_timer(current_army_timer)

    if current_army_timer > 0:
        army_timer.start()

    # Programmatic Transit Road Rendering
    for child in get_children():
        if child is not Path2D: continue
        var curve: Curve2D = (child as Path2D).curve
        var path_line: Line2D = Line2D.new()
        path_line.name = "PathLine_" + child.name
        path_line.width = 0.7
        path_line.default_color = Color(0.0, 0.0, 0.0, 0.1)
        for i in range(curve.point_count):
            var pos: Vector2 = curve.get_point_position(i)
            path_line.add_point(pos)
        add_child(path_line)

    queue_redraw()

func _set_mouse_filter_ignore_recursive(node: Node) -> void:
    if node is Control:
        (node as Control).mouse_filter = Control.MOUSE_FILTER_IGNORE
    for child in node.get_children():
        _set_mouse_filter_ignore_recursive(child)
\end{lstlisting}

% =========================================================================
\subsection{Structural Node Hierarchy and Scene Composition}
\label{subsec:province_node_hierarchy}

To encapsulate the complex behavioral, topological, and graphical requirements of a territory into a reusable component, every province is implemented as a dedicated composite scene in the Godot Engine. The scene hierarchy organizes discrete specialized nodes to handle spatial collision detection, vector surface rendering, perimeter outlining, urban telemetry overlays, and military pathfinding networks.

\subsubsection{Hierarchical Node Architecture}
The internal scene tree of a \texttt{Province} is structured around a nested object hierarchy:

\begin{figure}[H]
    \centering
    \includegraphics[width=0.5\textwidth]{imgs/structure_of_a_province.png} 
    \caption{Example of scene tree of the province St. Petersburg}
    \label{fig:structure_of_a_province}
\end{figure}

\begin{itemize}
    \item \textbf{Root Node (\texttt{Area2D}):} The base of the province scene graph is an \texttt{Area2D} node, whose instance name directly defines the geopolitical designation of the territory (e.g., \textit{"Kharkov"}). It serves as the primary spatial sensory interface, emitting decoupled observer signals (\texttt{area\_entered} and \texttt{area\_exited}) to intercept when military units cross provincial boundaries and trigger corresponding BDI perceptual events.
    \item \textbf{Collision Boundary (\texttt{CollisionPolygon2D}):} Direct child of the root \texttt{Area2D}, encapsulating the array of geometric vertices that delineate the exact physical footprint of the province. It provides the computational foundation for broad-phase physics detection, neighbor discovery routines, and spectator mouse raycasting.
    \item \textbf{Surface Fill Mesh (\texttt{Polygon2D}):} Configured with the identical vertex coordinates of the collision polygon, this node programmatically renders the continuous 2D colored mesh representing the province's surface, reflecting the owning nation's color palette alongside dynamic selection tinting.
    \item \textbf{Perimeter Outliner (\texttt{ProvinceBorder}):} A dedicated \texttt{Line2D} component that traces the perimeter vertices of the polygon to visually render the thin domestic borders separating adjacent provinces.
    \item \textbf{Provincial Capital Anchor (\texttt{City} / \texttt{Node2D}):} An instantiated child scene representing the urban capital seat of the territory. The \texttt{City} composite scene contains:
    \begin{itemize}
        \item A root \texttt{Node2D} establishing the local spatial origin of the provincial seat.
        \item Two \texttt{Sprite2D} / \texttt{TextureRect} nodes dedicated to rendering the structural icon of the settlement and the specific commodity icon produced by the territory.
        \item A \texttt{Control} UI overlay that dynamically displays the name label of the province, utilizing specialized container styling (\texttt{PanelContainer}) with high-contrast darkened backgrounds to visually demarcate national capitals.
    \end{itemize}
    \begin{figure}[H]
        \centering
        \includegraphics[width=0.5\textwidth]{imgs/structure_of_a_city.png} 
        \caption{Scene tree of a city node}
        \label{fig:structure_of_a_city}
    \end{figure}
    \begin{figure}[H]
        \centering
        \includegraphics[width=0.5\textwidth]{imgs/visual_representation_of_a_city.png} 
        \caption{Visual representation of a city node}
        \label{fig:visual_representation_of_a_city}
    \end{figure}
    \item \textbf{Transit Corridor Network (\texttt{Path2D} Array):} A collection of child \texttt{Path2D} nodes containing mathematical \texttt{Curve2D} resources. These curves represent all navigable road segments connecting the urban seat to neighboring provinces, providing the explicit kinematic tracks utilized by moving army divisions.
\end{itemize}

\subsubsection{Parametric Provincial Specialization}
While all provinces share this uniform structural node composition, individual territories are differentiated by their spatial geometries, their urban seat positions, their internal path network layouts, and a series of engine-exported configuration attributes:

\begin{itemize}
    \item \textbf{\texttt{province\_material} (\texttt{Enums.Resources}):} Specifies the primary commodity generated by the province (e.g., Wood, Iron, Oil, Gas, Coal, Grain, or Fish).
    \item \textbf{\texttt{double\_material} (\texttt{bool}):} Flags whether the province operates under a specialized double-production regime, directly altering economic extraction yields.
    \item \textbf{\texttt{startup\_buildings} (\texttt{Array[Enums.Buildings]}):} Defines the pre-existing infrastructure instances automatically provisioned in an operational state (\texttt{FUNCTIONING}) at the start of the match without incurring initial construction resource costs.
\end{itemize}

\begin{lstlisting}[language=python, caption={Exported configuration parameters and node references for provincial scene specialization in GDScript.}, label={lst:province_exports}]
class_name Province
extends Area2D

# Exported Configuration Parameters
@export var province_material: Enums.Resources
@export var double_material: bool
@export var startup_buildings: Array[Enums.Buildings]

# Node-Tree Component References
@onready var fill: Polygon2D = $Polygon2D
@onready var collision_polygon: CollisionPolygon2D = $CollisionPolygon2D
@onready var province_border: Line2D = $ProvinceBorder

@onready var city_node: Node2D = $City if has_node("City") else $Node2D
@onready var control: Control = city_node.get_node("Control")
@onready var resource_texture: TextureRect = city_node.get_node("Control/ResourceTexture")
@onready var province_name_capital_container: PanelContainer = city_node.get_node("Control/SubViewportContainer/SubViewport/CapitalLabelContainer")
@onready var province_name_capital_label: Label = city_node.get_node("Control/SubViewportContainer/SubViewport/CapitalLabelContainer/MarginContainer/ProvinceName")
@onready var province_name_label: Label = city_node.get_node("Control/SubViewportContainer/SubViewport/ProvinceNameLabel")
\end{lstlisting}

\section{The Military Subsystem: Army Structure and Custom Pathfinding}
\label{sec:military_subsystem}

In the simulation environment, military units are represented by instances of the \texttt{Army} scene. To support dynamic creation, scalable deployment, and real-time execution as outlined in the system architecture, each army division is modeled as a reusable, self-contained scene component within the Godot Engine. This enables the simulation to instantiate, configure, and manipulate armies dynamically at runtime based on incoming BDI practical reasoning commands.

% -------------------------------------------------------------------------
\subsection{Army Scene Composition and Component Hierarchy}
\label{subsec:army_composition}

The structural composition of the \texttt{Army} scene is rooted in a \texttt{PathFollow2D} node. This fundamental architectural design choice is dictated by the simulation kinematics: army divisions must move strictly along the pre-defined graph edges (\texttt{Path2D} components) connecting provincial capitals across the world map. By inheriting from \texttt{PathFollow2D}, each division can manipulate its spatial location along a parent curve via the \texttt{progress} and \texttt{progress\_ratio} properties.

Nested beneath the root \texttt{PathFollow2D} node, the scene hierarchy organizes several specialized graphical, informational, and sensory sub-nodes:

\begin{figure}[H]
    \centering
    \includegraphics[width=0.5\textwidth]{imgs/structure_of_an_army.png} 
    \caption{Scene tree of an army node}
    \label{fig:structure_of_an_army}
\end{figure}

\begin{itemize}
    \item \textbf{Visual Representation Nodes (\texttt{Line2D}, \texttt{Mark}, \texttt{MarkSelected}, and \texttt{Troop}):} A composite visual layer responsible for rendering the army's on-screen presence. The \texttt{Troop} sprite provides the physical unit token, while \texttt{Mark} and \texttt{MarkSelected} visually highlight the army's position and active selection state. Additionally, \texttt{Line2D} trajectory lines project the planned path across the map.
    \item \textbf{Telemetry and State Display (\texttt{Control}):} A specialized user interface overlay that renders real-time division metrics on screen. This canvas includes visual indicators for the current soldier headcount, the division's total remaining life points, and the unique \texttt{army\_id} assigned relative to the owning nation.
    \item \textbf{Combat Interception Area (\texttt{Area2D}):} A dedicated 2D physical collision trigger. It intercepts when an army collides with other unit instances in space, allowing the engine to detect when opposing hostile divisions make contact and must initiate an engagement.
    \item \textbf{Concentric Sensory Detection Fields (\texttt{PerfectVision} and \texttt{ObfuscatedVision}):} Two specialized \texttt{Area2D} nodes with circular collision perimeters that realize the multi-tiered Fog of War model:
    \begin{itemize}
        \item \textbf{\texttt{PerfectVision}:} Configured with a baseline radius of $R = 90\text{ px}$. When an enemy army enters this perimeter, the owning agent acquires complete, unobstructed tactical information regarding the foreign force (e.g., exact soldier count, life points, and identity).
        \item \textbf{\texttt{ObfuscatedVision}:} Configured by scaling the baseline vision radius by a factor of $1.5$, yielding an extended detection radius of $R_{\text{obf}} = 90 \times 1.5 = 135\text{ px}$. When an enemy division is within this outer perimeter but outside the perfect vision circle, the agent is merely alerted to the existence and coordinates of the foreign force, while its specific soldier count and health metrics remain obscured.
    \end{itemize}
\end{itemize}

\begin{figure}[H]
    \centering
    \includegraphics[width=0.5\textwidth]{imgs/visual_representation_of_an_army.png} 
    \caption{Visual representation of an army node}
    \label{fig:visual_representation_of_an_army}
\end{figure}

% -------------------------------------------------------------------------
\subsection{Custom Weighted Graph Pathfinding: The \texttt{GPA2D} Framework}
\label{subsec:gpa2d_pathfinding}

The Godot Engine provides a native graph search framework through its built-in \texttt{AStar2D} class. In standard applications, \texttt{AStar2D} allows developers to define a network of 2D waypoints (in this context, provincial capital seats) interconnected by topological graph edges (represented by the provincial \texttt{Path2D} nodes). 

However, Godot's default \texttt{AStar2D} implementation presents a fundamental architectural limitation: it permits the assignment of custom traversal weight scales (\texttt{weight\_scale}) strictly to the graph \textit{nodes} (vertices), rather than allowing independent weight configurations on individual \textit{connections} (edges).

In a grand strategy military simulation, this limitation is critical. The movement speed and travel friction of an army division are directly determined by the geopolitical sovereign status of the specific territory being traversed:
\begin{itemize}
    \item \textbf{Domestic Territory:} An army traveling within its own nation's borders moves at $100\%$ nominal velocity ($\text{multiplier} = 1.0$).
    \item \textbf{Allied Territory:} An army marching across the sovereign territory of a formal ally moves at a reduced speed ($\text{multiplier} = 0.70$).
    \item \textbf{Hostile / War Territory:} An army advancing through enemy provinces encounters severe transit friction and guerrilla resistance, reducing its speed to just $35\%$ ($\text{multiplier} = 0.35$).
\end{itemize}

If graph pathfinding were computed purely based on topological node distances or node-only weights, an army calculating the shortest route between two domestic provinces might erroneously choose a path cutting across hostile enemy territory because the geometric distance is marginally shorter, completely ignoring the fact that the severe speed penalty ($0.35\times$) would make the actual travel duration substantially longer. 

To overcome this issue and ensure realistic, time-optimized route calculations, a custom weighted $A^*$ pathfinding engine---named \textbf{\texttt{GPA2D}} (\textit{Graph Pathfinding Algorithm 2D})---was implemented from scratch.

\subsubsection{Object-Oriented Topological Structure: \texttt{GPANode} and \texttt{GPAEdge}}
To achieve high code maintainability, clean encapsulation, and structural clarity, the internal topology of \texttt{GPA2D} is modeled through two dedicated helper classes: \texttt{GPANode} and \texttt{GPAEdge}.

\begin{itemize}
    \item \textbf{\texttt{GPANode}:} Encapsulates a distinct waypoint on the map. Each node maintains a unique integer identifier (\texttt{id}), a 2D spatial coordinate vector (\texttt{position}), a vertex weight modifier (\texttt{weight\_scale}), a boolean flag to exclude disabled nodes from search routines (\texttt{disabled}), and an internal collection of outbound directed connections (\texttt{connections}) typed as \texttt{Array[GPAEdge]}.
    \item \textbf{\texttt{GPAEdge}:} Models a directed connection bridging the current node to an adjacent target waypoint. It encapsulates the destination node identifier (\texttt{\_to\_id}) and an explicit numerical traversal cost (\texttt{\_weight}).
\end{itemize}

\begin{lstlisting}[language=python, caption={Structural definitions of GPANode and GPAEdge helper classes in GDScript.}, label={lst:gpa_node_edge}]
class GPANode:
    var id: int
    var position: Vector2
    var weight_scale: float
    var disabled: bool
    var connections: Array[GPAEdge]
    
    func _init(new_id: int, new_position: Vector2, new_weight_scale: float) -> void:
        id = new_id
        position = new_position
        weight_scale = new_weight_scale
        disabled = false
        connections = []
        
    func add_connection(to_id: int, weight: float) -> void:
        for el in connections:
            if el.get_to_id() == to_id: return
        connections.append(GPAEdge.new(to_id, weight))

class GPAEdge:
    var _to_id: int
    var _weight: float
    
    func _init(to_id: int, weight: float = 1.0):
        _to_id = to_id
        _weight = weight
        
    func get_to_id() -> int:
        return _to_id
        
    func get_weight() -> float:
        return _weight
        
    func set_to_id(to_id: int) -> void:
        _to_id = to_id
        
    func set_weight(weight: float) -> void:
        _weight = weight
\end{lstlisting}

\subsubsection{The \texttt{GPA2D} Interface and Core Methods}
To minimize refactoring overhead and ensure frictionless interoperability with Godot's existing ecosystem, \texttt{GPA2D} mirrors the public method signatures and conventions of the native \texttt{AStar2D} class:

\begin{itemize}
    \item \textbf{\texttt{add\_point(id, position, weight\_scale)}:} Instantiates and registers a new \texttt{GPANode} within the internal graph array at the specified index, storing its 2D coordinates and initial weight multiplier.
    \item \textbf{\texttt{get\_point\_position(id)}:} Retrieves the 2D spatial coordinate vector of the waypoint matching the specified ID.
    \item \textbf{\texttt{set\_point\_weight\_scale(id, weight\_scale)}:} Dynamically updates the traversal weight scale associated with a specific graph node.
    \item \textbf{\texttt{connect\_points(id, to\_id, weight, bidirectional)}:} Establishes an edge between two waypoints, scaling the baseline Euclidean distance by the supplied edge \texttt{weight}. If the \texttt{bidirectional} flag is \texttt{true}, a reciprocal connection is automatically established in the reverse direction.
    \item \textbf{\texttt{set\_connection\_weight(from\_id, to\_id, weight)}:} Modifies the edge weight of an existing connection between two nodes in either direction.
    \item \textbf{\texttt{get\_closest\_point(to\_position, include\_disabled)}:} Iterates across all active waypoints to compute Euclidean spatial distances, returning the unique ID of the node situated closest to the query coordinate vector.
    \item \textbf{\texttt{get\_id\_path(from\_id, to\_id, allow\_partial\_path)}:} Executes the weighted $A^*$ search algorithm between origin and destination nodes. If the destination is unreachable and \texttt{allow\_partial\_path} is enabled, the method returns the optimal sub-route leading to the nearest accessible node based on the minimum heuristic score.
\end{itemize}

\begin{lstlisting}[language=python, caption={The GPA2D graph pathfinding implementation and weighted search loop in GDScript.}, label={lst:gpa2d_methods}]
class_name GPA2D
extends Object

var graph: Array[GPANode] = []

func add_point(id: int, position: Vector2, weight_scale: float) -> void:
    graph.insert(id, GPANode.new(id, position, weight_scale))
    
func connect_points(id: int, to_id: int, weight: float = 1.0, bidirectional: bool = true) -> void:
    graph[id].add_connection(to_id, weight * (graph[id].position - graph[to_id].position).length())
    if bidirectional:
        graph[to_id].add_connection(id, weight * (graph[id].position - graph[to_id].position).length())

func get_closest_point(to_position: Vector2, include_disabled: bool = false) -> int:
    var closest_point = null
    for el in graph:
        if closest_point == null and (not el.disabled or include_disabled):
            closest_point = el
            continue
        if el.position.distance_to(to_position) < closest_point.position.distance_to(to_position) and (not el.disabled or include_disabled):
            closest_point = el
    if closest_point == null:
        return -1
    return closest_point.id

func get_id_path(from_id: int, to_id: int, allow_partial_path: bool = false) -> PackedInt64Array:
    if from_id < 0 or from_id >= graph.size() or to_id < 0 or to_id >= graph.size():
        return PackedInt64Array()
    if graph[from_id] == null or graph[to_id] == null or graph[from_id].disabled:
        return PackedInt64Array()

    var open_set: Array[int] = [from_id]
    var came_from: Dictionary = {}
    
    var g_score: Dictionary = {}
    g_score[from_id] = 0.0
    
    var f_score: Dictionary = {}
    f_score[from_id] = _heuristic(from_id, to_id)
    
    var closest_node: int = from_id
    var min_h_score: float = f_score[from_id]

    while open_set.size() > 0:
        var current_index: int = 0
        var current: int = open_set[0]
        var lowest_f: float = f_score.get(current, INF)
        
        for i in range(1, open_set.size()):
            var node_id = open_set[i]
            var f = f_score.get(node_id, INF)
            if f < lowest_f:
                lowest_f = f
                current = node_id
                current_index = i
                
        if current == to_id:
            return _reconstruct_path(came_from, current)
            
        open_set.remove_at(current_index)
        var current_node: GPANode = graph[current]
        
        for edge in current_node.connections:
            var neighbor_id: int = edge.get_to_id()
            var neighbor_node: GPANode = graph[neighbor_id]
            
            if neighbor_node.disabled:
                continue
                
            var tentative_g: float = g_score[current] + (edge.get_weight() * neighbor_node.weight_scale)
            
            if tentative_g < g_score.get(neighbor_id, INF):
                came_from[neighbor_id] = current
                g_score[neighbor_id] = tentative_g
                
                var h: float = _heuristic(neighbor_id, to_id)
                f_score[neighbor_id] = tentative_g + h
                
                if h < min_h_score:
                    min_h_score = h
                    closest_node = neighbor_id
                    
                if not open_set.has(neighbor_id):
                    open_set.append(neighbor_id)
                    
    if allow_partial_path and closest_node != from_id:
        return _reconstruct_path(came_from, closest_node)
        
    return PackedInt64Array()

func _heuristic(from_id: int, to_id: int) -> float:
    return graph[from_id].position.distance_to(graph[to_id].position)

func _reconstruct_path(came_from: Dictionary, current: int) -> PackedInt64Array:
    var path: PackedInt64Array = PackedInt64Array()
    path.append(current)
    while came_from.has(current):
        current = came_from[current]
        path.append(current)
    path.reverse()
    return path
\end{lstlisting}

\subsubsection{Visual Graph and Path Debugging Utilities}
To facilitate the verification and testing of complex multi-nation graph topologies during development, \texttt{GPA2D} incorporates integrated visual debugging routines:

\begin{itemize}
    \item \textbf{Graph Topology Inspection (\texttt{debug\_graph}):} Dynamically instantiates a visual debug canvas containing circular markers at waypoint coordinates, numerical node ID labels, and red \texttt{Line2D} connection segments illustrating all active edges. The debug overlay can be rendered indefinitely or destroyed after a specified timeout.
    \item \textbf{Path Solution Visualization (\texttt{debug\_path}):} Evaluates and visually renders the solution path between a start node and a target node. It highlights the origin waypoint in blue, subsequent path nodes in green, and connects them with prominent trajectory lines to visually validate the algorithm's routing choices.
\end{itemize}

\begin{lstlisting}[language=python, caption={Visual graph inspection and path debugging routines in GPA2D.}, label={lst:gpa2d_debug}]
func debug_graph(tree: SceneTree, time: float = -1.0) -> void:
    var graph_container: Node = Node.new()
    for node in graph:
        var l: Line2D = Line2D.new()
        var t: Label = Label.new()
        var circle: Line2D = _debug_create_circle(node.position, 5.0, Color(1.0, 0.0, 0.0, 1.0))
        t.position = node.position
        t.text = str(node.id)
        t.scale = Vector2(0.5, 0.5)
        circle.add_child(t)
        graph_container.add_child(circle)
        
        l.width = 1.0
        l.default_color = Color(1.0, 0.0, 0.0, 1.0)
        l.add_point(node.position)
        for c in node.connections:
            l.add_point(graph[c.get_to_id()].position)
        graph_container.add_child(l)
        
    tree.current_scene.add_child(graph_container)
    if time != -1.0:
        await tree.create_timer(time).timeout
        tree.current_scene.remove_child(graph_container)

func debug_path(start: int, target: int, tree: SceneTree, time: float = -1.0) -> void:
    var path_container: Node = Node.new()
    var path = get_id_path(start, target)
    for point in path:
        var node = graph[point]
        var l: Line2D = Line2D.new()
        var color: Color = Color(0.0, 0.0, 1.0, 1.0) if point == start else Color(0.0, 1.0, 0.0, 1.0)
        var circle: Line2D = _debug_create_circle(node.position, 5.0, color)
        
        l.width = 1.0
        l.default_color = Color(0.0, 1.0, 0.0, 1.0)
        l.add_point(node.position)
        
        path_container.add_child(l)
        path_container.add_child(circle)
    tree.current_scene.add_child(path_container)
    if time != -1.0:
        await tree.create_timer(time).timeout
        tree.current_scene.remove_child(path_container)

func _debug_create_circle(center: Vector2, radius: float, color: Color, points: int = 64) -> Line2D:
    var pts: Array[Vector2] = []
    var line: Line2D = Line2D.new()
    line.width = 1.0
    line.default_color = color
    line.closed = true
    for i in range(points):
        var angle = TAU * float(i) / points
        pts.append(center + Vector2(cos(angle), sin(angle)) * radius)
    line.points = pts
    return line
\end{lstlisting}

% =========================================================================
\subsection{Army Movement and Multi-Leg Traversal Subsystem}
\label{subsec:army_movement}

The spatial relocation of a military division is triggered by the invocation of the \texttt{order(city)} method. Dispatched from the nation controller in response to deliberative BDI intentions, this routine coordinates multi-segment trajectory calculation, scene re-parenting, path interpolation, dynamic trajectory visualization, and geopolitical transition checks across sovereign borders.

% -------------------------------------------------------------------------
\subsubsection{March Order Dispatching and Redundancy Filtering}
When an agent commands an army to march toward a target destination (\texttt{city}), the division first performs an idempotency check to verify whether the incoming command is structurally redundant:

\begin{lstlisting}[language=python, caption={Order redundancy evaluation in GDScript.}, label={lst:order_redundancy}]
if travelling and path != null and len(path) > 0:
    var last_node = path[len(path) - 1]
    var point_path: Path2D = player.graph_to_path[last_node][0]
    var path_point: int = player.graph_to_path[last_node][1]
    var current_dest = point_path.to_global(point_path.curve.get_point_position(path_point))
    var new_dest = city.to_global(city.get_city_position())
    if current_dest.distance_to(new_dest) < 1.0:
        return
\end{lstlisting}

If the division is actively moving and the spatial distance between its current final waypoint and the requested provincial seat is less than $1.0\text{ px}$, the execution terminates immediately. This filtering mechanism prevents unnecessary path re-evaluations and avoids resetting trajectory states when the BDI deliberative loop re-asserts existing march intentions.

If the order is valid, the shortest weighted route is computed via the nation's \texttt{GPA2D} instance (\texttt{player.routes}):
\begin{enumerate}
    \item The nearest topological graph waypoints to the army's current world position and the target province's capital coordinates are extracted using \texttt{get\_closest\_point()}.
    \item The custom weighted $A^*$ routine \texttt{get\_id\_path()} computes an ordered array of waypoint identifiers (\texttt{points}). If no path exists, the command aborts.
    \item The army identifies the initial \texttt{Path2D} node corresponding to the first leg of the route. If the army is not currently attached to this node, dynamic scene re-parenting is performed via \texttt{reparent(new\_path)}.
    \item The initial curve distance and movement orientation (\texttt{dir} $\in \{-1, 1\}$) are resolved, updating the \texttt{progress} and \texttt{direction} attributes.
    \item Trajectory indicator lines are updated, the intermediate destination vector is broadcast via the \texttt{movement\_updated} signal, and an updated operational telemetry packet is dispatched to the BDI reasoning engine via \texttt{dispatch\_my\_army\_update(self)}.
\end{enumerate}

\begin{lstlisting}[language=python, caption={Complete implementation of the \texttt{order()} method in GDScript.}, label={lst:order_full}]
func order(city: Province) -> void:
    if travelling and path != null and len(path) > 0:
        var last_node = path[len(path) - 1]
        var point_path: Path2D = player.graph_to_path[last_node][0]
        var path_point: int = player.graph_to_path[last_node][1]
        var current_dest = point_path.to_global(point_path.curve.get_point_position(path_point))
        var new_dest = city.to_global(city.get_city_position())
        if current_dest.distance_to(new_dest) < 1.0:
            return

    var start: int = player.routes.get_closest_point(global_position)
    var target: int = player.routes.get_closest_point(city.to_global(city.get_city_position()))
    var points: PackedInt64Array = player.routes.get_id_path(start, target, false)
    if len(points) == 0: return
    var new_path: Path2D = player.graph_to_path[points[0]][0]
    var point: int = player.graph_to_path[points[0]][1]
    var dir = new_path.get_point_direction(point)
    if get_parent() != new_path:
        reparent(new_path)
        
    var distance = new_path.curve.get_point_position(point).distance_to(position)
    if dir == -1:
        self.progress = distance
    else:
        self.progress = new_path.curve.get_point_position(point).distance_to(new_path.curve.get_point_position(0)) - distance
    direction = dir
    path = points
    nodes_travelled = 0
    travelling = true
    target_province = city

    direction_line.default_color = MOVE_LINE_COLOR
    direction_line_partial.default_color = ATTACK_LINE_COLOR
    _update_direction_lines()

    var next_pos = _get_next_path_position()
    if next_pos != Vector2.ZERO:
        emit_signal("movement_updated", next_pos)
    if is_instance_valid(player):
        player.dispatch_my_army_update(self)
\end{lstlisting}

% -------------------------------------------------------------------------
\subsubsection{Dynamic Trajectory Line Rendering}
To provide clear visual feedback for human observers and debugging sessions, the active march corridor is rendered through two distinct \texttt{Line2D} canvas items:

\begin{itemize}
    \item \textbf{Global Route Projection (\texttt{direction\_line}):} Configured with a green tint (\texttt{MOVE\_LINE\_COLOR}), this line connects the army's current world position through all remaining waypoints up to the final destination ($\{ \text{path}[k] \mid k \in [\text{nodes\_travelled}, |\text{path}|-1] \}$).
    \item \textbf{Immediate Leg Vector (\texttt{direction\_line\_partial}):} Configured with a red tint (\texttt{ATTACK\_LINE\_COLOR}), this segment isolates the immediate trajectory spanning from the army's current coordinates to the next intermediate curve waypoint.
\end{itemize}

Whenever a division advances to a new path leg or reaches its destination, \texttt{\_update\_direction\_lines()} clears and rebuilds the vertex arrays:

\begin{lstlisting}[language=python, caption={Dynamic trajectory line reconstruction routine in GDScript.}, label={lst:update_direction_lines}]
func _update_direction_lines() -> void:
    if not travelling or path == null or len(path) == 0:
        direction_line.clear_points()
        direction_line_partial.clear_points()
        return

    direction_line.clear_points()
    direction_line.add_point(global_position)
    for i in range(nodes_travelled, len(path)):
        var graph_point = path[i]
        var point_path: Path2D = player.graph_to_path[graph_point][0]
        var path_point: int = player.graph_to_path[graph_point][1]
        var global_pt = point_path.to_global(point_path.curve.get_point_position(path_point))
        direction_line.add_point(global_pt)

    direction_line_partial.clear_points()
    if nodes_travelled < len(path):
        direction_line_partial.add_point(global_position)
        var next_point = path[nodes_travelled]
        var next_point_path: Path2D = player.graph_to_path[next_point][0]
        var next_path_point: int = player.graph_to_path[next_point][1]
        var next_global_pt = next_point_path.to_global(next_point_path.curve.get_point_position(next_path_point))
        direction_line_partial.add_point(next_global_pt)
\end{lstlisting}

% -------------------------------------------------------------------------
\subsubsection{Kinematic Progress Loop and Geopolitical Boundary Incursions}
Kinematic advancement along curve resources is regulated by the periodic global tick (\texttt{\_on\_global\_timer\_timeout()}). During each execution cycle, the method evaluates the army's travel state:

\begin{enumerate}
    \item \textbf{Intra-Segment Progress:} While the army has not reached the boundary of its current \texttt{Path2D} node ($\text{progress\_ratio} \notin \{0.0, 1.0\}$ depending on direction), its spatial offset is advanced:
    \begin{equation}
        \text{progress}_{t+1} = \text{progress}_t + (v_{\text{current}} \times \text{direction})
    \end{equation}
    where $v_{\text{current}}$ represents the territory-adjusted speed scalar.
    \item \textbf{Inter-Segment Node Transition:} When the unit reaches the terminal vertex of a curve, \texttt{nodes\_travelled} is incremented by one. If all waypoints have been traversed ($\text{nodes\_travelled} = |\text{path}|$), locomotion halts, the path is cleared, and \texttt{\_on\_destination\_reached()} is triggered. Otherwise, the division dynamically re-parents to the next \texttt{Path2D} node in the sequence, recomputes the entry offset along the new curve, and updates movement direction.
    \item \textbf{Territorial Incursion and War Declaration Logic:} During transit, if an army crosses into a province that is neither owned by nor allied with its sovereign nation, diplomatic protocols execute automatically. If the territorial owner is neutral ($\text{owner} \notin \text{Allies} \land \text{owner} \notin \text{WarEnemies}$), the nation immediately issues a formal war declaration (\texttt{declare\_war(owner)}). Furthermore, if the incursion occurs against an active enemy nation whose provincial seat is currently undefended, the army triggers immediate territorial annexation via \texttt{game.conquer\_province(province\_location, player)}.
\end{enumerate}

\begin{lstlisting}[language=python, caption={Periodic timer event handler managing kinematic progress, inter-segment re-parenting, and territorial incursion logic.}, label={lst:army_timer_timeout}]
func _on_global_timer_timeout() -> void:
    if travelling:
        if (self.progress_ratio != 1.0 and direction == 1) or (self.progress_ratio != 0.0 and direction == -1):
            self.progress += current_speed * direction
            _update_direction_lines()
        else:
            nodes_travelled += 1

            var next_pos = _get_next_path_position()
            if next_pos != Vector2.ZERO:
                emit_signal("movement_updated", next_pos)

            if nodes_travelled == len(path):
                travelling = false
                path = null
                emit_signal("movement_updated", Vector2.ZERO)
                _on_destination_reached()
                return

            # Intermediate Geopolitical Incursion Check
            if province_location != null and not player.get_province_owner(province_location) and province_location.get_city_position() == position:
                var owner = province_location.get_parent()
                if owner is Agent:
                    if owner not in player.player_in_alliance and owner not in player.player_in_war:
                        player.declare_war(owner)
                    if owner in player.player_in_war:
                        var game: GameController = get_tree().current_scene
                        game.conquer_province(province_location, player)

            # Re-parenting to Subsequent Path Leg
            var new_path = player.graph_to_path[path[nodes_travelled]][0]
            var point = player.graph_to_path[path[nodes_travelled]][1]
            var dir = new_path.get_point_direction(point)
            if get_parent() != new_path:
                reparent(new_path)
            var distance = new_path.curve.get_point_position(point).distance_to(position)
            if dir == -1:
                self.progress = distance
            else:
                self.progress = new_path.curve.get_point_position(point).distance_to(new_path.curve.get_point_position(0)) - distance

            direction = dir
            _update_direction_lines()
            if is_instance_valid(player):
                player.dispatch_my_army_update(self)
\end{lstlisting}

% -------------------------------------------------------------------------
\subsubsection{Destination Resolution and Force Consolidation}
When the unit completes the final leg of its planned trajectory, the private handler \texttt{\_on\_destination\_reached()} processes destination arrival:

\begin{itemize}
    \item \textbf{Domestic Consolidation:} If the destination province is owned by the nation, the arriving army scans all domestic units stationed at that provincial seat. If another stationary division is present, \texttt{merge\_armies(self)} is invoked to consolidate the two detachments into a single unified army division.
    \item \textbf{Hostile / Neutral Destination Resolution:} If the destination province belongs to a foreign nation, the diplomatic and conquest validations are executed: automatic war declaration against neutral entities, or territorial conquest if the enemy territory lacks garrisoned defense forces.
    \item \textbf{State Clearing and Telemetry Dispatch:} The \texttt{target\_province} attribute is reset to \texttt{null}, trajectory lines are cleared, and the updated stationary state is serialized to the BDI reasoning layer via \texttt{dispatch\_my\_army\_update(self)}.
\end{itemize}

\begin{lstlisting}[language=python, caption={Destination arrival handler managing domestic consolidation and conquest checks in GDScript.}, label={lst:destination_reached}]
func _on_destination_reached() -> void:
    target_province = null
    _update_direction_lines()
    if province_location == null: return

    # Domestic Unit Consolidation
    for el in player.armies:
        if el != self and is_instance_valid(el) and el.province_location == province_location and not el.travelling:
            el.merge_armies(self)
            return

    # Foreign Destination Resolution
    if not player.get_province_owner(province_location):
        var owner = province_location.get_parent()
        if owner is Agent:
            if owner not in player.player_in_alliance and owner not in player.player_in_war:
                player.declare_war(owner)
            if owner in player.player_in_war:
                var game: GameController = get_tree().current_scene
                game.conquer_province(province_location, player)

    if is_instance_valid(player):
        player.dispatch_my_army_update(self)
\end{lstlisting}

% =========================================================================
\subsection{Combat Engagement, Dynamic Resolution, and Force Consolidation}
\label{subsec:combat_system}

Spatial unit collisions in the simulation space are intercepted via the physical \texttt{Area2D} collision trigger on each army division. When an overlap event occurs, the signal handler \texttt{\_on\_area\_entered(area)} is invoked. The system first evaluates the geopolitical status of the underlying territory: if the encounter occurs in a province owned by a neutral third party ($\text{owner} \notin \text{Allies} \land \text{owner} \notin \text{WarEnemies}$), a diplomatic war declaration is automatically triggered against that nation.

Following territorial validation, the algorithm queries the owning nation (\texttt{player}) of both divisions. If the interacting armies belong to the same sovereign nation or formal alliance coalition, the execution routes to \texttt{\_on\_friendly\_collision(other)}; otherwise, if a hostile war status is confirmed, it routes to \texttt{\_on\_enemy\_collision(other)}.

% -------------------------------------------------------------------------
\subsubsection{Friendly Unit Merging and Force Consolidation}
To prevent excessive node allocations and optimize pathfinding traversals, divisions belonging to the same sovereign nation merge upon arrival at a shared provincial seat. The routine \texttt{\_on\_friendly\_collision(other)} enforces structural constraints before consolidation:

\begin{itemize}
    \item Formal alliance units belonging to different sovereign nations do not merge.
    \item Merging occurs strictly between a moving division and a stationary division stationed at the destination provincial seat ($\text{nodes\_travelled} = |\text{path}| - 1$).
    \item The unit's arrival vertex along the active \texttt{Curve2D} resource must match the exact 2D position vector of the stationary army.
\end{itemize}

When validated, \texttt{other.merge\_armies(self)} accumulates the soldier headcount and remaining life points into the stationary unit, while invoking \texttt{destroy\_self()} on the arriving division to free its node hierarchy.

\begin{lstlisting}[language=python, caption={Friendly unit collision evaluation and force merging routines in GDScript.}, label={lst:friendly_collision_merge}]
func _on_friendly_collision(other: Army) -> void:
    if other.player in player.player_in_alliance: return
    if not travelling or other.travelling: return
    if nodes_travelled < len(path) - 1: return
    
    var path_points = get_parent().curve
    if (direction == 1 and path_points.get_point_position(path_points.point_count - 1) == other.position) or \
       (direction == -1 and path_points.get_point_position(0) == other.position):
        other.merge_armies(self)

func merge_armies(other: Army) -> void:
    soldier_number += other.soldier_number
    life_points += other.life_points
    other.destroy_self()
    queue_redraw()
    if is_instance_valid(player):
        player.dispatch_my_army_update(self)
\end{lstlisting}

% -------------------------------------------------------------------------
\subsubsection{Hostile Interception and Combat Role Assignment}
When opposing hostile divisions collide, \texttt{\_on\_enemy\_collision(other)} halts locomotion (\texttt{travelling = false}) and establishes offensive/defensive tactical roles:

\begin{itemize}
    \item \textbf{Asymmetric Encounter:} If one army is stationary while the other is in transit, the stationary division assumes the role of \textit{Defender}, while the moving unit acts as \textit{Attacker} by calling \texttt{begin\_attack(other)}.
    \item \textbf{Head-On Interception:} If both divisions are actively traveling toward each other along the same transit curve, both units simultaneously act as attackers and defenders. Two symmetrical engagement instances are initiated via \texttt{begin\_attack(other)} and \texttt{other.begin\_attack(self)}.
\end{itemize}

\begin{lstlisting}[language=python, caption={Hostile division collision handler and combat initiation in GDScript.}, label={lst:enemy_collision}]
func _on_enemy_collision(other: Army) -> void:
    if not travelling: return
    travelling = false
    if not other.travelling:
        begin_attack(other)
    else:
        begin_attack(other)
        other.begin_attack(self)
\end{lstlisting}

% -------------------------------------------------------------------------
\subsubsection{Engagement Lifecycle and Asynchronous Combat Rounds}
Upon combat initialization via \texttt{begin\_attack(other)}, a dedicated, one-shot \texttt{Timer} instance (\texttt{battle\_timer}) configured with a round duration of $t_{\text{round}} = 60.0\text{ s}$ (\texttt{ATTACK\_TIME}) is attached as a child node. The flag \texttt{day\_battle} is asserted on both combatants, and the initial strike round is resolved immediately:

\begin{lstlisting}[language=python, caption={Combat timer initialization and destruction signal binding in GDScript.}, label={lst:begin_attack}]
func begin_attack(other: Army) -> void:
    if not is_instance_valid(other): return
    var battle_timer := Timer.new()
    battle_timer.wait_time = ATTACK_TIME
    battle_timer.one_shot = true
    battle_timer.timeout.connect(_on_battle_timer_timeout.bind(other, battle_timer))
    add_child(battle_timer)
    other.army_destroyed.connect(_on_enemy_destroyed.bind(battle_timer))
    _on_battle_timer_timeout(other, battle_timer)
    battle_timer.start()
    day_battle = true
    other.day_battle = true
\end{lstlisting}

\paragraph{Mathematical Combat Formulations}
During each combat tick (\texttt{\_on\_battle\_timer\_timeout}), nominal offensive and defensive strengths are computed as a function of soldier count, base unit ratings ($\text{AtkUnit} = 4.0$, $\text{DefUnit} = 5.0$), and the division's current health ratio ($\frac{\text{HP}}{N \times \text{HPUnit}}$, with $\text{HPUnit} = 20.0$):
\begin{align}
    S_{\text{attacker}} &= (N_{\text{atk}} \times \text{AtkUnit}) \times \frac{\text{HP}_{\text{atk}}}{N_{\text{atk}} \times \text{HPUnit}} \\
    S_{\text{defender}} &= (N_{\text{def}} \times \text{DefUnit}) \times \frac{\text{HP}_{\text{def}}}{N_{\text{def}} \times \text{HPUnit}}
\end{align}

To introduce non-deterministic tactical friction—simulating variations in local battlefield terrain, weather anomalies, tactical leadership variance, and fluctuating unit morale—both strength outputs are modulated by stochastic multipliers sampled from a continuous uniform distribution $\mu \sim \mathcal{U}(0.5, 1.5)$:
\begin{align}
    \text{Damage}_{\text{inflicted}} &= S_{\text{attacker}} \times \mu_{\text{atk}} \\
    \text{Damage}_{\text{absorbed}} &= S_{\text{defender}} \times \mu_{\text{def}}
\end{align}

\paragraph{Attrition and Casualty Resolution}
The applied damage values directly reduce the respective division's hit points. Concurrently, discrete soldier casualties and cumulative death counters are derived from health depletion:
\begin{align}
    \Delta N_{\text{casualties}} &= \left\lfloor \frac{\text{Damage}}{\text{HPUnit}} \right\rfloor \\
    N_{t+1} &= N_t - \Delta N_{\text{casualties}} \\
    \text{HP}_{t+1} &= \text{HP}_t - \text{Damage}
\end{align}

If a division's soldier count ($N \le 0$) or life points ($\text{HP} \le 0$) drop to zero, the unit is destroyed via \texttt{destroy\_self()}, halting and freeing the \texttt{battle\_timer}. The state update is dispatched to the BDI agents through VEsNA. If both combatants survive the exchange, the \texttt{battle\_timer} restarts for another round, iterating until total annihilation of one force.

\paragraph{March Resumption Post-Victory}
When an enemy division is destroyed, the signal handler \texttt{\_on\_enemy\_destroyed()} frees the active combat timer, resets \texttt{day\_battle = false}, and inspects the remaining waypoint array. If the victorious division maintains an incomplete transit itinerary ($\text{path} \neq \emptyset$), it re-asserts \texttt{travelling = true}, reconstructs its trajectory vectors via \texttt{\_update\_direction\_lines()}, and resumes its march toward the target destination.

\begin{lstlisting}[language=python, caption={Stochastic combat resolution, attrition math, and march resumption handlers in GDScript.}, label={lst:battle_resolution}]
func _on_battle_timer_timeout(other: Army, battle_timer: Timer) -> void:
    var random_attacker_multiplier: float = rng.randf_range(0.5, 1.5)
    var random_defender_multiplier: float = rng.randf_range(0.5, 1.5)
    var attacker_strength: float = 0.0
    if soldier_number > 0:
        attacker_strength = (soldier_number * ATTACK_PER_UNIT) * life_points / (soldier_number * LIFE_PER_UNIT)
    var defener_strength: float = 0.0
    if is_instance_valid(other) and other.soldier_number > 0:
        defener_strength = (other.soldier_number * other.DEFENSE_PER_UNIT) * other.life_points / (other.soldier_number * other.LIFE_PER_UNIT)
    var final_attack: float = attacker_strength * random_attacker_multiplier
    var final_defense: float = defener_strength * random_defender_multiplier
    
    life_points -= final_defense
    soldier_number -= int(final_defense / LIFE_PER_UNIT)
    death_count += int(final_defense / LIFE_PER_UNIT)
    
    if is_instance_valid(other):
        other.life_points -= final_attack
        other.soldier_number -= int(final_attack / other.LIFE_PER_UNIT)
        other.death_count += int(final_attack / other.LIFE_PER_UNIT)

    var battle_finished: bool = false
    if not is_instance_valid(other) or other.soldier_number <= 0 or other.life_points <= 0:
        if is_instance_valid(other):
            other.destroy_self()
        battle_finished = true
    else:
        other.queue_redraw()
        if is_instance_valid(other.player):
            other.player.dispatch_my_army_update(other)
    
    if (soldier_number <= 0 or life_points <= 0) and not battle_finished:
        destroy_self()
        battle_finished = true
    else:
        queue_redraw()
        if is_instance_valid(player):
            player.dispatch_my_army_update(self)
    
    if battle_finished or not is_instance_valid(other):
        battle_timer.stop()
        battle_timer.queue_free()
        if path != null and len(path) > 0 and not destroyed:
            travelling = true
            _update_direction_lines()
            if is_instance_valid(player):
                player.dispatch_my_army_update(self)
    else:
        battle_timer.start()

func _on_enemy_destroyed(destroyed_army: Army, battle_timer: Timer) -> void:
    if is_instance_valid(battle_timer):
        battle_timer.stop()
        battle_timer.queue_free()
    
    if not destroyed:
        day_battle = false
        if path != null and len(path) > 0 and not travelling:
            travelling = true
            _update_direction_lines()
            if is_instance_valid(player):
                player.dispatch_my_army_update(self)
\end{lstlisting}

% =========================================================================
\subsection{Division Maintenance, Vision Management, and Dynamic Force Splitting}
\label{subsec:army_lifecycle_vision_split}

To complete the operational specification of the military subsystem, this section details the mechanics governing daily health regeneration, engine lifecycle notifications (dynamic speed modulation and asynchronous deletion dispatching), multi-tiered sensory field interception, spectator view switching, and dynamic force splitting.

% -------------------------------------------------------------------------
\subsubsection{Daily Division Healing and Combat Idle Gating}
Army survivability is regulated across long-term cycles via the daily clock event (\texttt{\_on\_global\_day\_timer\_timeout()}). In grand strategy simulations, garrisoned troops stationed in peaceful domestic provinces gradually recover combat readiness, provided they are not actively engaged in hostilities.

The system tracks daily combat engagement using the boolean state flag \texttt{day\_battle}, which is asserted whenever a division participates in a combat round. When \texttt{GlobalDayTimer} times out:
\begin{enumerate}
    \item If \texttt{day\_battle} is \texttt{true} or if the army is stationed in a province not owned by its sovereign nation, health recovery is aborted, and the flag is reset to \texttt{false} for the subsequent day.
    \item If the division is peacefully stationed in friendly territory, its current life points increment proportionally based on local provincial morale $M \in [0, 100]$:
    \begin{equation}
        \text{HP}_{t+1} = \begin{cases}
            \text{HP}_t \times 1.05, & \text{if } M > 80\% \quad (+5\% \text{ high-morale recovery}) \\
            \text{HP}_t \times 1.01, & \text{if } M < 20\% \quad (+1\% \text{ low-morale recovery}) \\
            \text{HP}_t \times 1.02, & \text{otherwise} \quad (+2\% \text{ baseline recovery})
        \end{cases}
    \end{equation}
\end{enumerate}

\begin{lstlisting}[language=python, caption={Daily health regeneration routine based on domestic morale in GDScript.}, label={lst:army_day_timer}]
func _on_global_day_timer_timeout() -> void:
    if day_battle or not player.get_province_owner(province_location):
        day_battle = false
        return
    
    if province_location.province_morale > 80:
        life_points += life_points * 0.05
    elif province_location.province_morale < 20:
        life_points += life_points * 0.01
    else:
        life_points += life_points * 0.02
\end{lstlisting}

% -------------------------------------------------------------------------
\subsubsection{Engine Lifecycle Notifications: Dynamic Speed Scaling and Destruction Telemetry}
The \texttt{\_notification(what)} virtual method intercepts low-level engine lifecycle notifications to synchronize internal simulation parameters and external agent cognitive beliefs:

\begin{itemize}
    \item \textbf{\texttt{NOTIFICATION\_PARENTED}:} Triggered whenever an army changes its parent node (e.g., when transitioning from one \texttt{Path2D} curve to another across provincial borders). The division updates its spatial reference \texttt{province\_location} and recalibrates its traversal velocity \texttt{current\_speed}:
    \begin{equation}
        v_{\text{current}} = v_{\text{base}} \times \mu_{\text{territory}}
    \end{equation}
    where $\mu_{\text{territory}}$ is resolved dynamically based on territorial jurisdiction:
    \begin{itemize}
        \item \textbf{Domestic or Allied Territory:} Scaled by \texttt{ALLY\_TERRITORY\_SPEED\_MULTIPLIER} ($0.70\times$).
        \item \textbf{Maritime Corridors (\texttt{Background}):} Scaled by \texttt{SEA\_SPEED\_MULTIPLIER} ($0.10\times$).
        \item \textbf{Hostile / Foreign Territory:} Scaled by \texttt{ENEMY\_TERRITORY\_SPEED\_MULTIPLIER} ($0.35\times$).
    \end{itemize}
    \item \textbf{\texttt{NOTIFICATION\_PREDELETE}:} Intercepts the imminent deallocation of the \texttt{Army} node from system memory (RAM). Before the node hierarchy is purged, the division sets \texttt{destroyed = true}, emits the \texttt{army\_destroyed} signal, and constructs an asynchronous JSON telemetry frame dispatched to VEsNA:
    \begin{equation}
        \text{Packet} = \big\{ \text{"type"}: \text{"sight"}, \, \text{"type"}: \text{"my\_army\_destroyed"}, \, \text{"id"}: \text{army\_id} \big\}
    \end{equation}
    This prevents deliberative BDI reasoning loops from dispatching orders or assigning military targets to divisions that have been eliminated in combat.
\end{itemize}

\begin{lstlisting}[language=python, caption={Engine lifecycle event interceptor managing territorial velocity scaling and pre-deletion telemetry dispatch.}, label={lst:army_notifications}]
func _notification(what: int) -> void:
    if what != NOTIFICATION_PARENTED and what != NOTIFICATION_PREDELETE: return

    if what == NOTIFICATION_PARENTED:
        province_location = get_parent().get_parent()
        if player.get_province_owner(province_location) or (province_location.get_parent() is Agent and province_location.get_parent() in player.player_in_alliance):
            current_speed = SPEED * ALLY_TERRITORY_SPEED_MULTIPLIER
        elif province_location is Background:
            current_speed = SPEED * SEA_SPEED_MULTIPLIER
        else:
            current_speed = SPEED * ENEMY_TERRITORY_SPEED_MULTIPLIER
    elif what == NOTIFICATION_PREDELETE:
        destroyed = true
        emit_signal("army_destroyed", self)
        if is_instance_valid(player):
            var msg: Dictionary = {
                "type": "sight",
                "sender": player.name,
                "receiver": "vesna",
                "data": {
                    "type": "my_army_destroyed",
                    "id": army_id
                }
            }
            player.send_vesna_message(msg)
\end{lstlisting}

% -------------------------------------------------------------------------
\subsubsection{Multi-Tiered Sensory Interception and Fog of War Management}
Environmental sensing is implemented through signal bindings attached to the \texttt{PerfectVision} and \texttt{ObfuscatedVision} circular areas. When foreign units intersect these boundaries, the division mutates the parent agent's cognitive collections:

\begin{itemize}
    \item \textbf{Obfuscated Field Transitions:} The handlers \texttt{\_on\_obfuscated\_vision\_entered} and \texttt{\_on\_obfuscated\_vision\_exited} manage entry and departure from the outer $135\text{ px}$ radius, invoking \texttt{player.add\_known\_army(other)} and \texttt{player.remove\_know\_army\_reference(other)}. This registers the existence and spatial coordinates of the foreign division in the agent's \texttt{known\_armies} data structure without exposing troop numbers.
    \item \textbf{Perfect Field Transitions:} The handlers \texttt{\_on\_perfect\_vision\_entered} and \texttt{\_on\_perfect\_vision\_exited} manage entry and departure from the inner $90\text{ px}$ radius, calling \texttt{player.add\_perfectly\_known\_army(other)} and \texttt{player.remove\_perfectly\_know\_army\_reference(other)}. This exposes full operational metrics (soldier count, remaining hit points, and unique ID) within the agent's \texttt{perfectly\_known\_armies} repository, synchronizing the corresponding BDI beliefs via VEsNA.
\end{itemize}

\begin{lstlisting}[language=python, caption={Concentric vision radius signal handlers managing Fog of War telemetry updates in GDScript.}, label={lst:vision_signals}]
func _on_obfuscated_vision_entered(area: Area2D) -> void:
    if area.name != "Area2D": return
    var other: Node = area.get_parent()
    if other is Army and player != other.player:
        player.add_known_army(other)

func _on_obfuscated_vision_exited(area: Area2D) -> void:
    if area.name != "Area2D": return
    var other: Node = area.get_parent()
    if other is Army and player != other.player:
        player.remove_know_army_reference(other)

func _on_perfect_vision_entered(area: Area2D) -> void:
    if area.name != "Area2D": return
    var other: Node = area.get_parent()
    if other is Army and player != other.player:
        player.add_perfectly_known_army(other)

func _on_perfect_vision_exited(area: Area2D) -> void:
    if area.name != "Area2D": return
    var other: Node = area.get_parent()
    if other is Army and player != other.player:
        player.remove_perfectly_know_army_reference(other)
\end{lstlisting}

% -------------------------------------------------------------------------
\subsubsection{Spectator UI Utilities and Dynamic View Filtering}
To maintain visual clarity for human observers during real-time multi-agent execution, the \texttt{Army} class provides specialized rendering and input handlers:

\begin{itemize}
    \item \textbf{Perspective-Aware Trajectory Filtering (\texttt{\_on\_view\_changed}):} When the spectator switches active perspectives between sovereign agents, trajectory line visibility toggles dynamically. If observing the division's owning nation (\texttt{view == player}) or in neutral observer mode, the full trajectory (\texttt{direction\_line}) is displayed. If observing an opposing sovereign nation, the global plan is hidden, and only the immediate partial vector (\texttt{direction\_line\_partial}) is rendered.
    \item \textbf{Morale Ratio Extraction (\texttt{\_get\_morale\_percentage}):} Computes the current health ratio relative to initial maximum capacity, factoring in cumulative battlefield casualties:
    \begin{equation}
        M_{\text{army}} = \frac{\text{HP}}{(N_{\text{current}} + N_{\text{deaths}}) \times \text{HPUnit}} \times 100\%
    \end{equation}
    \item \textbf{Camera-Relative UI Rescaling (\texttt{\_input}):} Intercepts zoom inputs to dynamically invert the scale of the \texttt{Control} telemetry panel ($\text{scale} = \frac{1.0}{\text{camera.zoom}}$), preserving readable font sizes and UI dimensions across varying zoom levels.
\end{itemize}

\begin{lstlisting}[language=python, caption={Perspective filtering, morale calculation, and UI rescaling routines in GDScript.}, label={lst:army_utilities}]
func _on_view_changed(view: Agent) -> void:
    if view == null or view == player:
        direction_line.show()
        direction_line_partial.hide()
    else:
        direction_line.hide()
        direction_line_partial.show()

func _get_morale_percentage() -> float:
    return life_points / ((soldier_number + death_count) * LIFE_PER_UNIT) * 100.0

func _input(_event: InputEvent) -> void:
    if Input.is_action_just_pressed("camera_zoom_in") or Input.is_action_just_pressed("camera_zoom_out"):
        control.scale = Vector2(1, 1) / get_viewport().get_camera_2d().zoom
\end{lstlisting}

% -------------------------------------------------------------------------
\subsubsection{Dynamic Tactical Force Splitting}
To execute the tactical force allocation, border garrisoning, and precision offensive maneuvers formulated in the cognitive design chapter, army divisions implement dynamic division splitting via \texttt{split(soldiers\_to\_split)}:

\begin{enumerate}
    \item \textbf{Validation and Ratio Partitioning:} The method verifies that $0 < N_{\text{split}} < N_{\text{current}}$. It computes the exact proportional partition ratio $r = \frac{N_{\text{split}}}{N_{\text{current}}}$ and derives the split hit point allocation:
    \begin{equation}
        \text{HP}_{\text{split}} = \text{HP}_{\text{current}} \times r
    \end{equation}
    \item \textbf{Origin Division Reduction:} The origin division decrements its soldier count by $N_{\text{split}}$ and its hit points by $\text{HP}_{\text{split}}$, redrawing its UI and notifying the BDI engine via \texttt{dispatch\_my\_army\_update(self)}.
    \item \textbf{Detachment Instantiation and Kinematic State Inheritance:} A new \texttt{Army} node is dynamically instantiated from \texttt{army.tscn}, assigned a sequential unique ID (\texttt{player.next\_army\_id++}), and initialized with $N_{\text{split}}$ and $\text{HP}_{\text{split}}$. It is attached as a sibling to the active parent \texttt{Path2D}, inheriting the identical spatial position, curve offset (\texttt{progress}), and normalized progress ratio (\texttt{progress\_ratio}).
    \item \textbf{Itinerary Cloning and Signal Wiring:} If the original army is actively traveling, the detachment clones the active \texttt{path} array, current \texttt{nodes\_travelled}, and travel \texttt{direction}, asserting \texttt{travelling = true} and initializing its own trajectory line renderers. Signal connections for destruction and sensory discovery are wired to the \texttt{GameController}, the dynamic Fog of War vision circle is registered, and the new division is serialized to the BDI reasoning architecture.
\end{enumerate}

\begin{lstlisting}[language=python, caption={Dynamic division force splitting implementation in GDScript.}, label={lst:army_split}]
func split(soldiers_to_split: int) -> Army:
    if soldiers_to_split <= 0 or soldiers_to_split >= soldier_number:
        return null
        
    var ratio = float(soldiers_to_split) / float(soldier_number)
    var split_life_points = life_points * ratio

    soldier_number -= soldiers_to_split
    life_points -= split_life_points
    queue_redraw()
    if is_instance_valid(player):
        player.dispatch_my_army_update(self)

    var army_scene = load("res://scenes/army.tscn")
    var new_army = army_scene.instantiate()
    new_army.army_id = player.next_army_id
    player.next_army_id += 1
    new_army.soldier_number = soldiers_to_split
    new_army.life_points = split_life_points
    new_army.player = player
    player.armies.append(new_army)

    get_parent().add_child(new_army)
    new_army.position = position
    new_army.progress = self.progress
    new_army.progress_ratio = self.progress_ratio

    if travelling:
        new_army.path = path
        new_army.nodes_travelled = nodes_travelled
        new_army.direction = direction
        new_army.travelling = true
        new_army.direction_line.default_color = MOVE_LINE_COLOR
        new_army.direction_line_partial.default_color = ATTACK_LINE_COLOR
        new_army._update_direction_lines()

    var game: GameController = get_tree().current_scene as GameController
    new_army.army_destroyed.connect(game._on_army_destroyed)
    new_army.became_known.connect(game._on_army_became_known)
    new_army.became_perfectly_known.connect(game._on_army_became_perfectly_known)
    new_army.became_unknown.connect(game._on_army_became_unknown)

    if new_army.player == game.view or game.view == null:
        game._fog_add_circle(new_army, new_army.position, new_army.VISION_RANGE)

    if new_army.player != game.view and game.view != null:
        new_army.hide()

    new_army.queue_redraw()
    if is_instance_valid(new_army.player):
        new_army.player.dispatch_my_army_update(new_army)
    return new_army
\end{lstlisting}


% =========================================================================
\section{The Sovereign Nation Subsystem: Architecture and Agent Integration}
\label{sec:nation_subsystem}

Within the simulation framework, each sovereign geopolitical power is modeled by an instance of the \texttt{Agent} class. Rather than functioning solely as an organizational container, a nation node embodies the complete administrative, economic, diplomatic, and perceptual state of a player. It coordinates territorial ownership, manages bi-directional network communication with the external VEsNA middleware, handles dynamic cartographic typography, and serializes environmental telemetry into structured JSON envelopes for BDI cognitive deliberation.

% -------------------------------------------------------------------------
\subsection{Territorial Node Hierarchy and Dynamic Label Placement}
\label{subsec:nation_hierarchy_labels}

At the scene-tree level, a nation acts as the direct parent node for all \texttt{Province} instances under its sovereign jurisdiction. When a province is conquered, territorial transfer is executed through engine scene re-parenting (\texttt{reparent()}), dynamically reattaching the province node to the conquering nation. This structural hierarchy ensures that territorial queries (such as resource aggregation and provincial counting) are resolved directly via standard node iteration.

To visually demarcate sovereign territory for the spectator interface, each nation manages a dedicated \texttt{Control} label container (\texttt{player\_name\_control}). Because the geometric footprint of a nation evolves dynamically as provinces are annexed or ceded, static text positioning is insufficient. Instead, the routine \texttt{update\_name\_position()} computes the optimal centroid and orientation of the national territory:

\begin{enumerate}
    \item \textbf{Area-Weighted Centroid Extraction:} The method \texttt{\_get\_nation\_center()} iterates across all owned provinces, extracting their collision polygon vertices. For each province $k$, the local polygon centroid $\mathbf{C}_k$ and signed planar area $A_k$ are calculated via the shoelace formula:
    \begin{equation}
        A_k = \frac{1}{2} \sum_{i=0}^{N-1} (x_i y_{i+1} - x_{i+1} y_i), \quad \mathbf{C}_k = \frac{1}{6 A_k} \sum_{i=0}^{N-1} (\mathbf{P}_i + \mathbf{P}_{i+1})(x_i y_{i+1} - x_{i+1} y_i)
    \end{equation}
    The overall national centroid $\mathbf{C}_{\text{nation}}$ is derived as the area-weighted mean:
    \begin{equation}
        \mathbf{C}_{\text{nation}} = \frac{\sum_k |A_k| \, \mathbf{C}_k}{\sum_k |A_k|}
    \end{equation}
    If $\mathbf{C}_{\text{nation}}$ falls within non-contiguous space (e.g., in a maritime gulf between territories), the algorithm queries the nearest owned provincial capital to ensure the text anchor resides within domestic territory.
    \item \textbf{Principal Inertial Orientation:} The method \texttt{\_get\_nation\_orientation()} calculates the spatial covariance tensor of all global perimeter vertices relative to their mean vector:
    \begin{equation}
        \text{Cov}_{xx} = \frac{1}{N}\sum \Delta x^2, \quad \text{Cov}_{yy} = \frac{1}{N}\sum \Delta y^2, \quad \text{Cov}_{xy} = \frac{1}{N}\sum \Delta x \Delta y
    \end{equation}
    The principal orientation angle $\theta$ is derived and clamped ($\theta \in [-1.2, 1.2]\text{ rad}$):
    \begin{equation}
        \theta = \frac{1}{2} \operatorname{atan2}(2 \, \text{Cov}_{xy}, \, \text{Cov}_{xx} - \text{Cov}_{yy})
    \end{equation}
    \item \textbf{Iterative Boundary Fitting:} Starting from an initial bounding-box scale factor, the label's eight perimeter control points are evaluated via \texttt{\_is\_point\_in\_nation()}. If any corner intersects foreign or maritime space, the scale factor is iteratively decayed by $15\%$ per step until the text fits within the nation's contiguous landmass.
\end{enumerate}

% -------------------------------------------------------------------------
\subsection{Bi-Directional WebSocket Bridge and Asynchronous Message Queueing}
\label{subsec:websocket_bridge}

The nation node serves as the communication conduit between the continuous Godot simulation loop and the discrete AgentSpeak BDI reasoning layer. Network transport is implemented via a dedicated \texttt{TCPServer} listening on an exported port, paired with a \texttt{WebSocketPeer} instance:

\begin{lstlisting}[language=python, caption={WebSocket network infrastructure declarations in Agent.gd.}, label={lst:network_vars}]
var tcp_server := TCPServer.new()
var ws := WebSocketPeer.new()
var end_communication: bool = true
var _last_ws_state = WebSocketPeer.STATE_CLOSED
var _message_queue: Array[Dictionary] = []
\end{lstlisting}

\subsubsection{Asynchronous Connection Lifecycle and Queue Flushing}
Because cognitive BDI agents execute in external Java processes (via JaCaMo/VEsNA) and may initialize with variable startup latency relative to the Godot engine, sensory state packets emitted during early simulation frames cannot be transmitted immediately. 

To eliminate data loss and prevent simulation stalls, the nation implements a message queue (\texttt{\_message\_queue}). The dispatch routine \texttt{send\_vesna\_message()} inspects the connection state: if \texttt{ws.get\_ready\_state() == STATE\_OPEN}, the payload is serialized to a JSON string and transmitted over the socket; otherwise, the dictionary is appended to \texttt{\_message\_queue}.

During each simulation frame (\texttt{\_process()}), the connection is polled:
\begin{itemize}
    \item Pending incoming TCP connections are accepted via \texttt{ws.accept\_stream(conn)}.
    \item When a state transition to \texttt{STATE\_OPEN} is detected, \texttt{\_flush\_message\_queue()} exhausts the pending queue in first-in, first-out (FIFO) order, followed by the transmission of the full initial environment snapshot via \texttt{send\_initial\_state()}.
    \item Incoming network packets are read, deserialized from JSON, and dispatched to the intention handler \texttt{manage(intention)}.
\end{itemize}

\begin{lstlisting}[language=python, caption={WebSocket polling, FIFO queue flushing, and packet dispatching loop in GDScript.}, label={lst:network_process}]
func send_vesna_message(msg: Dictionary) -> void:
    if ws.get_ready_state() == WebSocketPeer.STATE_OPEN:
        ws.send_text(JSON.stringify(msg))
    else:
        _message_queue.append(msg)

func _flush_message_queue() -> void:
    while not _message_queue.is_empty():
        var msg = _message_queue.pop_front()
        ws.send_text(JSON.stringify(msg))

func _process(_delta: float) -> void:
    while tcp_server.is_connection_available():
        var conn: StreamPeerTCP = tcp_server.take_connection()
        assert(conn != null)
        ws.accept_stream(conn)
        
    ws.poll()
    var current_state = ws.get_ready_state()
    if current_state == WebSocketPeer.STATE_OPEN:
        if _last_ws_state != WebSocketPeer.STATE_OPEN:
            _flush_message_queue()
            send_initial_state()
        while ws.get_available_packet_count():
            var msg: String = ws.get_packet().get_string_from_ascii()
            var intention: Dictionary = JSON.parse_string(msg)
            manage(intention)
    _last_ws_state = current_state
\end{lstlisting}

\subsubsection{Action Dispatching and Intention Execution}
The \texttt{manage(message)} routine processes incoming BDI action requests, mapping high-level declarative intentions to concrete simulation operations:
\begin{itemize}
    \item \textbf{\texttt{move\_army}:} Resolves the matching \texttt{Army} division and issues a march command via \texttt{army.order(target\_province)}.
    \item \textbf{\texttt{split\_army}:} Executes dynamic force division via \texttt{army.split(soldiers)}, immediately commanding the newly instantiated detachment toward its assigned objective.
    \item \textbf{\texttt{declare\_war}:} Appends the target power to \texttt{player\_in\_war}, dissolves active alliances, and cascades war declarations across the target's formal coalition members.
    \item \textbf{Diplomatic Negotiations (\texttt{request\_peace}, \texttt{accept\_peace}, \texttt{request\_alliance}, \texttt{accept\_alliance}):} Resolves bilateral diplomatic treaties. For multi-member coalitions, peace requests are stored in \texttt{accepted\_peace\_requests}, ratifying formal peace treaties only when all allied coalition members confirm acceptance.
\end{itemize}

% -------------------------------------------------------------------------
\subsection{Telemetry Serialization and Decoupled Environment State}
\label{subsec:telemetry_serialization}

A core architectural principle of this system is the complete decoupling of the BDI reasoning engine from the simulation environment. Rather than hardcoding topological layouts, provincial properties, and game rules as static beliefs within AgentSpeak source files, the simulation dynamically serializes its complete world state into structured JSON envelopes. 

This design mirrors human player perception: agents dynamically discover circuit layouts, border relationships, resource outputs, and operational constraints through external sensory telemetry at runtime, allowing map layouts and rule sets to be modified without requiring alterations to agent logic scripts.

Upon connection establishment and at the start of each daily cycle (\texttt{\_on\_global\_day\_timer\_timeout()}), the nation executes \texttt{send\_initial\_state()}, compiling a comprehensive state dictionary:

\begin{lstlisting}[caption={Structure of the comprehensive state update JSON packet emitted by Agent.gd to VEsNA.}, label={lst:state_update_json}]
{
  "type": "sight",
  "sender": "body",
  "receiver": "PlayerName",
  "data": {
    "type": "state_update",
    "my_armies": [
      {
        "id": 1,
        "x": 120.5,
        "y": 340.2,
        "life": 100.0,
        "attack_value": 20.0,
        "location": "ProvinceA",
        "distances": { "ProvinceA": 0.0, "ProvinceB": 150.4 },
        "travelling": false,
        "target_province": ""
      }
    ],
    "allied_provinces": [
      {
        "id": "ProvinceA",
        "morale": 95.0,
        "x": 100.0,
        "y": 200.0,
        "is_capital": true,
        "is_double": false,
        "resource": "grain"
      }
    ],
    "enemy_provinces": [
      {
        "id": "ProvinceB",
        "owner": "EnemyNation",
        "x": 300.0,
        "y": 400.0,
        "has_army": true,
        "defense_value": 45.0,
        "bordering_allies": ["ProvinceA"],
        "is_capital": false,
        "is_double": true,
        "resource": "oil"
      }
    ],
    "nation_strengths": { "PlayerName": 60, "EnemyNation": 40 },
    "active_wars": 1,
    "neighbor_nations": {
      "EnemyNation": { "grain": 120, "oil": 50 }
    },
    "day": 5,
    "alliances": [["Player1", "Player2"]],
    "game_rules": {
      "riot_threshold": 33.0,
      "riot_breakpoint": 30.0,
      "riot_base_defense": 5.0,
      "riot_penalty_multiplier": 6.0,
      "life_per_unit": 20.0,
      "attack_per_unit": 4.0,
      "defense_per_unit": 5.0,
      "max_alliances": 1,
      "conquered_morale": 25.0
    }
  }
}
\end{lstlisting}

% =========================================================================
\subsection{Army Event Interception and Sensory Streaming}
\label{subsec:nation_army_events}

To maintain an accurate mental model within the BDI deliberative layer, the \texttt{Agent} node intercepts dynamic military events occurring across the simulation and streams updated perceptual frames to the VEsNA bridge. These event handlers process domestic force updates, foreign troop trajectory tracking, casualty notifications, and reference-counted Fog of War state transitions.

% -------------------------------------------------------------------------
\subsubsection{Domestic Telemetry Dispatching: \texttt{dispatch\_my\_army\_update}}
The method \texttt{dispatch\_my\_army\_update(army)} is invoked externally across multiple game events (such as force splitting, recruitment completion, or path leg advancement) to synchronize the operational metrics of domestic divisions. 

The method computes nominal offensive strength based on current unit life points:
\begin{equation}
    \text{AttackValue} = (N_{\text{soldiers}} \times \text{AttackPerUnit}) \times \frac{\text{LifePoints}}{N_{\text{soldiers}} \times \text{LifePerUnit}}
\end{equation}
Additionally, it queries the nation's \texttt{GPA2D} routing instance to precompute pathfinding network distances from the army's current coordinates to all domestic provincial capitals, compiling the data into a structured telemetry payload:

\begin{lstlisting}[caption={Structure of the domestic army telemetry JSON packet dispatched to VEsNA.}, label={lst:my_army_update_json}]
{
  "type": "sight",
  "sender": "body",
  "receiver": "PlayerName",
  "data": {
    "type": "my_army_update",
    "id": 1,
    "x": 120.5,
    "y": 340.2,
    "life": 100.0,
    "attack_value": 20.0,
    "location": "ProvinceA",
    "distances": {
      "ProvinceA": 0.0,
      "ProvinceB": 150.4
    },
    "travelling": true,
    "target_province": "ProvinceB"
  }
}
\end{lstlisting}

This serialized dictionary is transmitted to the reasoning engine via \texttt{send\_vesna\_message()}, allowing the cognitive layer to evaluate garrison assignments and tactical reachability in real time.

% -------------------------------------------------------------------------
\subsubsection{Foreign Threat Incursion Tracking: \texttt{\_on\_watched\_army\_movement}}
To enable proactive defensive interception, the nation monitors the trajectories of visible non-allied forces. Whenever a foreign unit advances along its active path, \texttt{\_on\_watched\_army\_movement(army, next\_pos)} is triggered.

The method evaluates the directional vector toward \texttt{next\_pos} (representing the terminal coordinates of the division's immediate path leg). It queries domestic territory polygons via point-in-polygon checks (\texttt{Geometry2D.is\_point\_in\_polygon()}) to identify whether the foreign unit is advancing toward an allied province. The projected threat telemetry is serialized into a \texttt{foreign\_army\_moving} envelope:

\begin{lstlisting}[caption={Structure of the foreign threat trajectory JSON packet.}, label={lst:foreign_army_moving_json}]
{
  "type": "sight",
  "sender": "body",
  "receiver": "PlayerName",
  "data": {
    "type": "foreign_army_moving",
    "id": 2,
    "owner": "EnemyNation",
    "x": 210.0,
    "y": 180.5,
    "next_x": 245.0,
    "next_y": 200.0,
    "target_province": "BorderProvinceA",
    "attack_value": 35.0
  }
}
\end{lstlisting}

This payload provides the BDI agent with the foreign unit's calculated offensive strength and estimated target province, allowing defensive response plans to mobilize interception forces before the foreign division reaches domestic borders.

% -------------------------------------------------------------------------
\subsubsection{Foreign Division Destruction: \texttt{\_on\_army\_destroyed}}
When a foreign division is annihilated in combat or deleted from memory, \texttt{\_on\_army\_destroyed(army)} cleans up internal tracking caches, disconnects movement signals, and emits a notification packet to the cognitive layer:

\begin{lstlisting}[caption={Structure of the foreign army destruction JSON packet.}, label={lst:foreign_army_destroyed_json}]
{
  "type": "sight",
  "sender": "body",
  "receiver": "PlayerName",
  "data": {
    "type": "foreign_army_destroyed",
    "id": 2,
    "owner": "EnemyNation"
  }
}
\end{lstlisting}

\paragraph{Disambiguation via National Ownership}
Because division identifiers (\texttt{army\_id}) are generated sequentially starting from zero within each independent nation, IDs are not globally unique across the simulation. Consequently, two opposing nations can simultaneously field units sharing the same numeric index (e.g., ID $0$). To eliminate ambiguity, all foreign army telemetry packets explicitly pair the \texttt{id} attribute with the owning nation designation (\texttt{owner}). This compound key ensures that the BDI agent purges only the corresponding foreign belief without corrupting identical IDs belonging to other powers.

% -------------------------------------------------------------------------
\subsubsection{Reference-Counted Fog of War Management}
The nation manages sensory tracking pools through four complementary methods: \texttt{add\_known\_army}, \texttt{remove\_know\_army\_reference}, \texttt{add\_perfectly\_known\_army}, and \texttt{remove\_perfectly\_know\_army\_reference}.

Because multiple friendly divisions can simultaneously project overlapping sensory fields over the same foreign unit, sensory tracking cannot rely on simple boolean visibility flags. If one friendly army moves away while another remains in visual contact, a binary flag would prematurely discard visibility.

To resolve this, the nation employs a reference-counting mechanism across two tracking dictionaries: \texttt{known\_armies} (obfuscated outer vision) and \texttt{perfectly\_known\_armies} (perfect inner vision):
\begin{itemize}
    \item \textbf{Field Ingress (\texttt{add\_known\_army} / \texttt{add\_perfectly\_known\_army}):} When a foreign unit enters a sensory perimeter, its reference counter in the corresponding dictionary is incremented. If the counter transitions from $0$ to $1$, signal bindings for movement and destruction are attached, and visibility state signals (\texttt{became\_known} or \texttt{became\_perfectly\_known}) are emitted to update the BDI Belief Base.
    \item \textbf{Field Egress (\texttt{remove\_know\_army\_reference} / \texttt{remove\_perfectly\_know\_army\_reference}):} When a sensory circle departs, the reference counter is decremented. Only when the counter reaches zero is the foreign unit purged from the active collection. If the army is no longer covered by any sensory field, its signal connections are disconnected, and \texttt{became\_unknown} is emitted, prompting the BDI agent to retract the corresponding tactical beliefs.
\end{itemize}

\begin{lstlisting}[language=python, caption={Reference-counted sensory field tracking and signal management example in Agent.gd.}, label={lst:vision_tracking_example}]
func add_known_army(army: Army) -> void:
    if not army.movement_updated.is_connected(_on_watched_army_movement.bind(army)):
        army.movement_updated.connect(_on_watched_army_movement.bind(army))
        if army.travelling:
            var next_pos = army._get_next_path_position()
            if next_pos != Vector2.ZERO:
                _on_watched_army_movement(army, next_pos)
    
    if not army.army_destroyed.is_connected(_on_army_destroyed):
        army.army_destroyed.connect(_on_army_destroyed)
        
    var was_known = known_armies.has(army) and known_armies[army] > 0
    var is_perfectly_known = perfectly_known_armies.has(army) and perfectly_known_armies[army] > 0

    known_armies.get_or_add(army, 0)
    known_armies[army] += 1

    if not was_known and not is_perfectly_known:
        army.emit_signal("became_known", army, self)

func remove_know_army_reference(army: Army) -> void:
    if army not in known_armies: return
    known_armies[army] -= 1
    if known_armies[army] == 0:
        known_armies.erase(army)
        var is_perfectly_known = perfectly_known_armies.has(army) and perfectly_known_armies[army] > 0
        if not is_perfectly_known:
            _last_watched_army_destinations.erase(army)
            if army.movement_updated.is_connected(_on_watched_army_movement.bind(army)):
                army.movement_updated.disconnect(_on_watched_army_movement.bind(army))
            if army.army_destroyed.is_connected(_on_army_destroyed):
                army.army_destroyed.disconnect(_on_army_destroyed)
            army.emit_signal("became_unknown", army, self)
\end{lstlisting}

% =========================================================================
\subsection{Diplomatic Actions, Coalition War, and Consensus Peace Protocols}
\label{subsec:diplomatic_protocols}

The \texttt{Agent} class encapsulates the state-machine logic required to resolve macro-diplomatic interactions triggered by the BDI reasoning engine. Because sovereign agents operate in an open multi-agent environment where military alliances form cohesive power blocs, diplomatic transitions cannot be evaluated purely on isolated bilateral terms. Instead, the simulation enforces multi-party coalition mechanics, mutual defensive pact propagation, and unanimous consensus peace ratifications.

% -------------------------------------------------------------------------
\subsubsection{War Declaration and Coalition Escalation}
A state of war is initialized through the \texttt{declare\_war(target\_agent)} method. When invoked by the action dispatcher \texttt{manage()}, the routine executes several cascading checks:

\begin{itemize}
    \item \textbf{Idempotency and Alliance Guardrails:} If the target nation is already present in \texttt{player\_in\_war}, the call aborts. Although diplomatic betrayals are currently restricted by high-level BDI social conventions, the routine includes safety assertions that dissolve mutual alliance entries if a war state is declared against an ally.
    \item \textbf{Symmetric State Assignment:} The target power is added to the local \texttt{player\_in\_war} array, and the target reciprocally invokes \texttt{declare\_war(self)}, ensuring consistent bilateral tracking across both node states.
    \item \textbf{Coalition War Cascading:} The method inspects the target's active alliance roster (\texttt{target.player\_in\_alliance}). Any allied third-party nations not currently engaged in the conflict are recursively targeted with war declarations, enforcing mutual defensive pact obligations across entire diplomatic coalitions.
    \item \textbf{Telemetry Broadcasting:} A \texttt{war\_declared} sight frame is serialized and transmitted to VEsNA, allowing the agent's deliberative planning loop to re-evaluate defensive priorities and offensive allocations.
\end{itemize}

% -------------------------------------------------------------------------
\subsubsection{Peace Negotiation and Coalition Consensus Verification}
Peace settlements follow an asynchronous request-confirmation handshake. An agent initiates negotiations via \texttt{request\_peace(target)}, transmitting a \texttt{peace\_request\_received} message to the opposing power. The recipient may subsequently accept or decline the petition:

\begin{itemize}
    \item \textbf{Peace Rejection (\texttt{decline\_peace}):} Emits a \texttt{peace\_request\_declined} notification back to the proposing nation, maintaining the active war state.
    \item \textbf{Peace Ratification and Demilitarized Retreat (\texttt{make\_peace}):} Erases the mutual war record across both agents, emits the \texttt{peace\_established} telemetry event, and iterates through all active domestic military divisions. To prevent accidental border skirmishes following treaty ratification, any division marching toward a former enemy's province is redirected to the nearest friendly provincial capital via \texttt{find\_nearest\_allied\_province()}.
\end{itemize}

\paragraph{Unanimous Coalition Peace Consensus}
When a war involves multi-party alliances, an individual nation cannot unilaterally sign a separate peace treaty that exposes its coalition partner. All nations fighting within the same allied block must explicitly agree to cease hostilities before formal peace is enacted.

This consensus protocol is handled within the \texttt{manage()} message dispatcher upon receiving an \texttt{accept\_peace} action:

\begin{lstlisting}[language=python, caption={Coalition consensus peace evaluation in Agent.gd.}, label={lst:coalition_peace}]
var target_name = data.get("target")
var players = get_tree().get_nodes_in_group("players")
for p in players:
    if p.name == target_name:
        var alliance_block = [p]
        for ally in p.player_in_alliance:
            if ally in player_in_war:
                alliance_block.append(ally)
        
        if alliance_block.size() == 1:
            make_peace(p)
        else:
            accepted_peace_requests[p] = true
            var all_accepted = true
            for member in alliance_block:
                if not accepted_peace_requests.get(member, false):
                    all_accepted = false
                    break
            if all_accepted:
                for member in alliance_block:
                    make_peace(member)
                    accepted_peace_requests.erase(member)
        break
\end{lstlisting}

When an \texttt{accept\_peace} signal arrives for a member of a coalition:
\begin{enumerate}
    \item The algorithm constructs the \texttt{alliance\_block} containing the target agent and all of its formal allies currently engaged in the war.
    \item If the block contains only a single nation ($\text{size} = 1$), \texttt{make\_peace(p)} is executed immediately.
    \item If an alliance block exists, the accepting member is registered in the \texttt{accepted\_peace\_requests} tracking dictionary. The algorithm then iterates across all members of the coalition: if every allied member has registered an acceptance ($\text{all\_accepted} = \text{true}$), formal peace is ratified with all coalition partners, and the temporary consensus cache is purged.
\end{enumerate}

% -------------------------------------------------------------------------
\subsubsection{Bilateral Alliance System and Power Balance Constraints}
To prevent the emergence of unbreakable hegemonic coalitions and preserve game balance, formal alliances are constrained by a strict cap:
\begin{equation}
    |\text{player\_in\_alliance}| \le \text{MAX\_ALLIANCES}, \quad \text{where } \text{MAX\_ALLIANCES} = 1
\end{equation}

The alliance lifecycle is governed by three primary methods:
\begin{itemize}
    \item \textbf{Alliance Proposal (\texttt{request\_alliance}):} Verifies that neither the proposing nation nor the target nation has exceeded the maximum alliance threshold ($\text{MAX\_ALLIANCES} = 1$). If valid, a \texttt{alliance\_request\_received} message is dispatched to the target; otherwise, an immediate decline notification is returned.
    \item \textbf{Alliance Rejection (\texttt{decline\_alliance}):} Transmits a \texttt{alliance\_request\_declined} packet to the proposing power without altering domestic state.
    \item \textbf{Alliance Ratification (\texttt{make\_alliance}):} Binds both powers into a mutual non-aggression and defensive partnership. The method appends the partner to \texttt{player\_in\_alliance} and executes a symmetric handshake invocation (\texttt{target.make\_alliance(self, true)}), broadcasting an \texttt{alliance\_established} event to synchronize the respective BDI Belief Bases.
\end{itemize}


% =========================================================================
\section{The Game Controller Subsystem}
\label{sec:game_controller_subsystem}

Within the simulation architecture, the \texttt{GameController} acts as the master orchestrator, supervising global lifecycle events, victory condition evaluations, macro-economic day progression, and global visual rendering layers. Among its primary graphical responsibilities is managing the dynamic Fog of War visual system, which provides human observers and spectators with an accurate, real-time representation of battlefield intelligence and partial sensory observability.

% -------------------------------------------------------------------------
\subsection{GPU Shading Architecture and Fog of War Implementation}
\label{subsec:gpu_shaders_fog}

To render the multi-tiered Fog of War without incurring severe CPU performance bottlenecks, the system leverages Godot's hardware-accelerated shading pipeline. In the Godot Engine, shaders are specialized programs written in \textbf{GDShader}—a high-performance shading language syntactically derived from GLSL ES 3.0. 

Rather than iterating sequentially through screen pixels on the CPU, GDShader programs are compiled into parallelized GPU execution kernels. By executing concurrently across thousands of streaming multiprocessors, shaders compute per-vertex deformations and per-pixel color transformations at high frame rates, enabling visual post-processing effects such as dynamic chromatic aberration, real-time atmospheric blur, volumetric glows, and non-linear coordinate distortions.

For 2D canvas rendering and screen-space overlays, Godot utilizes shaders of type \texttt{canvas\_item}. The execution pipeline exposes two fundamental entry-point functions:

\begin{itemize}
    \item \textbf{\texttt{vertex()}:} Invoked for every geometric vertex of a mesh or sprite. It permits programmatic manipulation of vertex coordinates, UV texture mappings, and vertex colors in local space prior to rasterization. This stage is commonly employed for geometric transformations and dynamic mesh warping, such as simulating physical wind displacement across foliage or undulating terrain meshes.
    \item \textbf{\texttt{fragment()}:} Executed for every individual rasterized pixel (fragment) covered by the visual canvas item. Its primary role is computing the final fragment output color (\texttt{COLOR}) by sampling textures, evaluating lighting gradients, performing spatial distance tests, and blending multi-channel color data. It is widely used to generate procedural noise patterns, screen transitions, and per-pixel visibility cutouts.
\end{itemize}

\subsubsection{Shader Uniforms and Per-Fragment Distance Evaluation}
To render the multi-radius visibility zones projected by deployed military divisions, the Fog of War system utilizes a custom \texttt{fragment()} shader attached to a dedicated full-screen \texttt{Sprite2D} node, designated as \texttt{FogSprite}, via a \texttt{ShaderMaterial}.

The shader receives dynamic runtime data from the simulation loop through \texttt{uniform} parameters. In GDShader, uniform variables represent parameters that remain constant across all fragments in a single render pass but can be updated dynamically from GDScript:
\begin{itemize}
    \item \textbf{\texttt{circle\_count}:} An integer specifying the active number of sensory emitters currently rendered.
    \item \textbf{\texttt{pos[64]}:} A statically allocated uniform array of \texttt{vec2} coordinates holding the pixel-space centers of all active vision circles.
    \item \textbf{\texttt{small\_radius[64]}:} A uniform float array storing the inner, perfect-vision radii associated with each respective division.
    \item \textbf{\texttt{big\_radius\_multiplier}:} A global scalar (set to $1.5$) that scales the inner radius to establish the outer, obfuscated-vision perimeter.
    \item \textbf{\texttt{color\_background}, \texttt{color\_small}, \texttt{color\_big}:} Color vectors defining the visual shroud (dense unobserved fog), perfect clarity, and partial obfuscation tinting, respectively.
\end{itemize}

\begin{lstlisting}[language=C, caption={The CanvasItem fragment shader implementing multi-tiered Fog of War in GDShader.}, label={lst:fog_shader_code}]
shader_type canvas_item;

uniform int circle_count = 0;
uniform vec2 pos[64];
uniform float small_radius[64];
uniform float big_radius_multiplier = 1.5;

uniform vec4 color_background : source_color = vec4(0.0, 0.0, 0.0, 0.75);
uniform vec4 color_small : source_color = vec4(0.0, 0.0, 0.0, 0.0);
uniform vec4 color_big : source_color = vec4(0.0, 0.0, 0.0, 0.35);

void fragment() {
    vec2 p = UV * vec2(textureSize(TEXTURE, 0));

    bool inside_big = false;
    bool inside_small = false;

    for (int i = 0; i < circle_count; i++) {
        float d = distance(p, pos[i]);
        if (d < small_radius[i]) {
            inside_small = true;
        }
    }

    for (int i = 0; i < circle_count; i++) {
        float d = distance(p, pos[i]);
        if (d < small_radius[i] * big_radius_multiplier) {
            inside_big = true;
        }
    }

    vec4 col = color_background;

    if (inside_small) {
        col = color_small;
    } else if (inside_big) {
        col = color_big;
    }

    COLOR = col;
}
\end{lstlisting}

\subsubsection{Fragment Execution Logic}
The fragment computation executes through the following sequence:
\begin{enumerate}
    \item \textbf{Pixel-Space Coordinate Transformation:} The normalized texture coordinate vector $\mathbf{UV} \in [0, 1]^2$ is multiplied by the absolute texture dimensions $\mathbf{S}_{\text{tex}} = \text{textureSize}(\text{TEXTURE}, 0)$, producing the fragment's discrete pixel coordinate $\mathbf{p} = (x, y)$.
    \item \textbf{Concentric Distance Queries:} The fragment evaluates its Euclidean distance $d_i = \|\mathbf{p} - \mathbf{pos}_i\|$ against every active emitter $i \in \{0, \dots, \text{circle\_count}-1\}$:
    \begin{align}
        \text{inside\_small} &= \exists \, i : d_i < R_{\text{small}, i} \\
        \text{inside\_big} &= \exists \, i : d_i < (R_{\text{small}, i} \times \mu_{\text{big}})
    \end{align}
    \item \textbf{Hierarchical Color Assignment:} The output color is resolved hierarchically:
    \begin{equation}
        \mathbf{C}_{\text{final}} = \begin{cases}
            \mathbf{C}_{\text{small}} \quad (\text{fully transparent}), & \text{if } \text{inside\_small} = \text{true} \\
            \mathbf{C}_{\text{big}} \quad (\text{semi-transparent}), & \text{else if } \text{inside\_big} = \text{true} \\
            \mathbf{C}_{\text{background}} \quad (\text{opaque shroud}), & \text{otherwise}
        \end{cases}
    \end{equation}
    The resulting color vector is assigned directly to the engine output register \texttt{COLOR}.
\end{enumerate}

This approach decouples perceptual rendering from CPU memory operations, allowing the simulation to update and render tens of dynamic vision radii simultaneously at 60 frames per second without degrading the execution speed of the underlying BDI cognitive agents.

% =========================================================================
\subsection{Fog of War Management Pipeline in \texttt{GameController}}
\label{subsec:fog_controller_management}

The GPU shader uniform parameters described in Section~\ref{subsec:gpu_shaders_fog} are managed dynamically by the \texttt{GameController} singleton. The controller acts as the orchestrator between high-level game entities (such as moving armies or camera perspective switches) and the low-level rendering material, ensuring that the sensory shroud updates consistently across all simulation ticks.

To achieve this, \texttt{GameController} exposes dedicated lifecycle, registration, and update routines.

% -------------------------------------------------------------------------
\subsubsection{Canvas Allocation and Material Setup: \texttt{\_create\_fog}}
During the initial scene composition phase, the controller configures the underlying texture canvas and binds the baseline shader constants via \texttt{\_create\_fog()}:

\begin{enumerate}
    \item \textbf{Dynamic Canvas Sizing:} The world-space dimensions $\mathbf{S}_{\text{world}} = (W, H)$ are extracted from the background collision geometry via \texttt{get\_collision\_rect\_size(background)}.
    \item \textbf{Image and Texture Generation:} An uncompressed \texttt{Image} buffer formatted as \texttt{FORMAT\_RGBA8} is created with matching pixel dimensions, filled with a semi-transparent black base tint, and wrapped into an \texttt{ImageTexture} bound to \texttt{fog\_sprite.texture}.
    \item \textbf{Uniform Parameter Initialization:} Baseline color vectors and the outer radius scalar ($\mu_{\text{big}} = 1.5$) are injected directly into the \texttt{ShaderMaterial} instance using \texttt{set\_shader\_parameter()}.
\end{enumerate}

\begin{lstlisting}[language=python, caption={Canvas allocation and initial shader uniform configuration in GDScript.}, label={lst:create_fog}]
func _create_fog() -> void:
    var world_dimensions: Vector2 = get_collision_rect_size(background)
    
    fog_image = Image.create(int(world_dimensions.x), int(world_dimensions.y), false, Image.FORMAT_RGBA8)
    fog_image.fill(Color(0, 0, 0, 0.5))
    var fog_texture: ImageTexture = ImageTexture.create_from_image(fog_image)
    fog_sprite.texture = fog_texture
    
    fog_material.set_shader_parameter("color_background", FOG_OF_WAR_BACKGROUND_COLOR)
    fog_material.set_shader_parameter("color_big", FOG_OF_WAR_OUTER_CIRCLE_COLOR)
    fog_material.set_shader_parameter("color_small", FOG_OF_WAR_INNER_CIRCLE_COLOR)
    fog_material.set_shader_parameter("big_radius_multiplier", FOG_OF_WAR_OUTER_CIRCLE_RADIUS_MULTIPLIER)
\end{lstlisting}

% -------------------------------------------------------------------------
\subsubsection{Sensory Circle Registration and Deletion}
Dynamic unit lifecycle events—such as unit recruitment, force splitting, or battlefield destruction—require immediate synchronization with the active uniform arrays:

\begin{itemize}
    \item \textbf{Emitter Registration (\texttt{\_fog\_add\_circle}):} When a new division enters the active view, its world position vector $\mathbf{P}_{\text{global}}$ is mapped into the texture's local pixel-coordinate space:
    \begin{equation}
        \mathbf{P}_{\text{texture}} = \mathbf{P}_{\text{global}} + \frac{\mathbf{S}_{\text{image}}}{2} - \mathbf{P}_{\text{sprite\_global}}
    \end{equation}
    The entry is stored within the \texttt{fog\_circles} dictionary keyed by the \texttt{Army} reference, and the packed coordinate and radius arrays are dispatched to the GPU.
    \item \textbf{Emitter Removal (\texttt{\_fog\_del\_circle}):} When a division is eliminated or moves out of the active spectator scope, \texttt{\_fog\_del\_circle(army)} removes the dictionary entry, rebuilds the condensed coordinate arrays, and updates the shader uniforms.
    \item \textbf{Buffer Clearing (\texttt{\_fog\_clear\_circles}):} Flushes all entries from \texttt{fog\_circles} and resets the \texttt{circle\_count} uniform to zero.
\end{itemize}

\begin{lstlisting}[language=python, caption={Dynamic registration and deletion of sensory circle emitters.}, label={lst:fog_add_del}]
func _fog_add_circle(army: Army, pos: Vector2, radius: float) -> void:
    fog_circles[army] = {
        "position": pos + Vector2(fog_image.get_size()) / 2 + (fog_sprite.global_position * Vector2(-1.0, -1.0)),
        "radius": radius
    }
    
    var circles_position: Array[Vector2] = []
    var circles_radius: Array[float] = []
    for el in fog_circles:
        circles_position.append(fog_circles[el]["position"])
        circles_radius.append(fog_circles[el]["radius"])
    
    fog_material.set_shader_parameter("circle_count", len(fog_circles))
    fog_material.set_shader_parameter("pos", circles_position)
    fog_material.set_shader_parameter("small_radius", circles_radius)

func _fog_del_circle(army: Army) -> void:
    fog_circles.erase(army)
    
    var circles_position: Array[Vector2] = []
    var circles_radius: Array[float] = []
    for el in fog_circles:
        circles_position.append(fog_circles[el]["position"])
        circles_radius.append(fog_circles[el]["radius"])
    
    fog_material.set_shader_parameter("circle_count", len(fog_circles))
    fog_material.set_shader_parameter("pos", circles_position)
    fog_material.set_shader_parameter("small_radius", circles_radius)
\end{lstlisting}

% -------------------------------------------------------------------------
\subsubsection{Kinematic Tracking and Perspective Remapping}
During continuous execution, the controller updates emitter positions and handles observer perspective changes:

\paragraph{Continuous Coordinate Synchronization (\texttt{\_fog\_update\_circle\_position})}
As units advance along path curves, \texttt{\_fog\_update\_circle\_position(army, new\_pos)} recalibrates the texture-space coordinate entry for the division and pushes the updated \texttt{pos} and \texttt{small\_radius} uniform arrays directly to the shader material without reallocating arrays.

\begin{lstlisting}[language=python, caption={Per-frame coordinate updates for traveling army divisions.}, label={lst:fog_update_pos}]
func _fog_update_circle_position(army: Army, new_pos: Vector2) -> void:
    if not fog_circles.has(army): return
    fog_circles[army]["position"] = new_pos + Vector2(fog_image.get_size()) / 2 + (fog_sprite.global_position * Vector2(-1.0, -1.0))
    
    var circles_position: Array[Vector2] = []
    var circles_radius: Array[float] = []
    for el in fog_circles:
        circles_position.append(fog_circles[el]["position"])
        circles_radius.append(fog_circles[el]["radius"])
    
    fog_material.set_shader_parameter("pos", circles_position)
    fog_material.set_shader_parameter("small_radius", circles_radius)
\end{lstlisting}

\paragraph{Spectator Perspective Remapping (\texttt{\_fog\_map\_circles})}
When a spectator switches between observing a specific sovereign agent (\texttt{view = player}) or adopting an omniscient global spectator mode (\texttt{view = null}), \texttt{\_fog\_map\_circles()} clears all active emitters and repopulates \texttt{fog\_circles}:
\begin{itemize}
    \item If \texttt{view == null}, it queries all divisions registered in the global \texttt{"armies"} node group.
    \item If a specific nation is focused, it isolates only that agent's domestic \texttt{view.armies} collection, configuring the shader to render the precise sensory visibility available to that player.
\end{itemize}


% =========================================================================
\subsection{Perspective-Aware Visualization, Diplomatic Shading, and Unit Mobilization}
\label{subsec:view_perspective_and_mobilization}

Beyond managing the GPU shader parameters, the \texttt{GameController} regulates perspective-dependent visual filtering for human observers and coordinates the centralized initialization and commissioning of military units.

% -------------------------------------------------------------------------
\subsubsection{Spectator Perspective and Vision State Dispatching}
Visual rendering of military divisions is strictly coupled to the active observer perspective (\texttt{view}). An observer may either adopt an omniscient spectator perspective ($\texttt{view} = \text{null}$) or inspect the simulation through the partial perceptual scope of a specific sovereign nation ($\texttt{view} = \text{Agent}$).

When a foreign army transitions between sensory boundaries, the controller intercepts these visibility updates and modulates unit visibility based on the active perspective:
\begin{itemize}
    \item \textbf{Obfuscated Discovery (\texttt{\_on\_army\_became\_known}):} If the detecting nation matches the currently observed perspective, the foreign army is displayed on the map (\texttt{army.show()}), but its soldier headcount label is concealed (\texttt{soldier\_number\_visible = false}), displaying an obscured marker.
    \item \textbf{Perfect Discovery (\texttt{\_on\_army\_became\_perfectly\_known}):} If the detecting nation matches the active view, the unit is made visible alongside full tactical telemetry (\texttt{soldier\_number\_visible = true}).
    \item \textbf{Sensory Loss (\texttt{\_on\_army\_became\_unknown}):} If a foreign unit departs the observing nation's sensory perimeters, it is hidden from the spectator viewport (\texttt{army.hide()}).
    \item \textbf{Combat Destruction (\texttt{\_on\_army\_destroyed}):} Purges the unit reference from the owning agent's military roster and deregisters its sensory emitter via \texttt{\_fog\_del\_circle(army)}.
\end{itemize}

% -------------------------------------------------------------------------
\subsubsection{Dynamic Diplomatic Territorial Shading}
To provide immediate visual clarity regarding geopolitical alignments, \texttt{GameController} dynamically alters provincial surface mesh tints (\texttt{province.fill.color}) whenever the observer perspective shifts or when diplomatic pacts are updated:

\begin{itemize}
    \item \textbf{Omniscient Spectator Mode ($\texttt{view} = \text{null}$):} All provinces are rendered in a uniform neutral tone (\texttt{VIEW\_NEUTRAL\_COLOR}), eliminating subjective alignment biases during multi-agent evaluations.
    \item \textbf{Sovereign Player Mode ($\texttt{view} = \text{Agent}$):} The visual mesh shifts to an egocentric diplomatic palette:
    \begin{equation}
        \mathbf{C}_{\text{province}} = \begin{cases}
            \mathbf{C}_{\text{agent}} \quad (\text{Green}), & \text{if owned by } \text{view} \\
            \mathbf{C}_{\text{ally}} \quad (\text{Blue}), & \text{else if owner } \in \text{view.alliances} \\
            \mathbf{C}_{\text{enemy}} \quad (\text{Red}), & \text{else if owner } \in \text{view.wars} \\
            \mathbf{C}_{\text{neutral}} \quad (\text{Gray}), & \text{otherwise}
        \end{cases}
    \end{equation}
\end{itemize}

Whenever a perspective switch occurs via \texttt{\_on\_view\_changed(player)}, the controller updates the active \texttt{view} reference, sweeps all military divisions to recalibrate visibility labels and node display states, invokes \texttt{\_fog\_map\_circles()} to remap sensory emitters, and recomputes provincial fill colors. Similarly, when diplomatic treaties evolve, \texttt{\_on\_relations\_changed()} triggers color updates and redraws unit telemetry panels across the scene tree.

% -------------------------------------------------------------------------
\subsubsection{Centralized Army Commissioning: \texttt{add\_army\_to\_province}}
The creation and initial physical placement of military forces are centralized within the \texttt{GameController} via \texttt{add\_army\_to\_province(province, army, sn)}. Centralizing force commissioning in the controller decouples unit instantiation from individual provincial nodes, ensuring that unique global identifier sequences (\texttt{next\_army\_id}), kinematic path graph parenting, signal wiring, and sensory registrations are executed deterministically.

\begin{lstlisting}[language=python, caption={Centralized army instantiation and path graph anchoring in GameController.gd.}, label={lst:add_army_to_province}]
func add_army_to_province(province: Province, army: Army = null, sn: int = 10) -> void:
    var agent: Agent = province.get_parent()
    if army == null:
        var paths: Array[Node] = get_tree().get_nodes_in_group("paths")
        var army_scene = preload("res://scenes/army.tscn")
        army = army_scene.instantiate()
        army.army_id = agent.next_army_id
        agent.next_army_id += 1
        army.soldier_number = sn
        army.life_points = army.LIFE_PER_UNIT * sn * (float(province.province_morale) / 100.0)
        army.player = agent
        agent.armies.append(army)
        if is_instance_valid(agent):
            agent.dispatch_my_army_update(army)

        for el in paths:
            if el is Path2D:
                var count: int = el.curve.point_count
                var found: bool = false
                for j in range(count):
                    if el.curve.get_point_position(j) == province.get_city_position():
                        el.add_child(army)
                        if j != 0:
                            army.progress_ratio = 1.0
                        found = true
                        break
                if found:
                    break

    army.army_destroyed.connect(_on_army_destroyed)
    army.became_known.connect(_on_army_became_known)
    army.became_perfectly_known.connect(_on_army_became_perfectly_known)
    army.became_unknown.connect(_on_army_became_unknown)

    if army.player == view or view == null:
        _fog_add_circle(army, army.global_position, army.VISION_RANGE)

    if agent != view and view != null:
        army.hide()
    army.queue_redraw()
\end{lstlisting}

The routine performs the following operations:
\begin{enumerate}
    \item \textbf{Unit Instantiation and Initialization:} If no pre-existing army instance is passed, it instantiates \texttt{army.tscn}, assigns a unique sequential \texttt{army\_id} from the owning agent, configures the soldier count (defaulting to $N = 10$), and initializes starting life points proportional to provincial morale:
    \begin{equation}
        \text{HP}_{\text{init}} = N_{\text{soldiers}} \times \text{HPUnit} \times \frac{M_{\text{province}}}{100.0}
    \end{equation}
    The newly commissioned division is appended to the agent's roster and registered via \texttt{dispatch\_my\_army\_update(army)}.
    \item \textbf{Kinematic Path Anchoring:} The division scans all provincial \texttt{Path2D} nodes to find a trajectory whose anchor vertex coincides with the provincial capital coordinates. It attaches the army as a child of that \texttt{Path2D} node, configuring \texttt{progress\_ratio} ($1.0$ if attached to the terminal curve vertex, $0.0$ otherwise) to anchor the division at the provincial seat.
    \item \textbf{Signal Binding and Fog of War Registration:} Signal hooks for destruction and multi-tiered sensory discovery are bound to the controller's event handlers. If the unit belongs to the active spectator perspective, its sensory emitter is registered in \texttt{fog\_circles}; otherwise, the unit is hidden from view until discovered.
\end{enumerate}

% =========================================================================
\subsection{Territorial Conquest, Civil Riots, and Interactive Selection}
\label{subsec:territorial_conquest_selection}

Beyond global visual updates, the \texttt{GameController} coordinates macro-geopolitical transitions when territories change sovereignty. Reassigning a province involves modifying the engine scene-graph hierarchy, updating border contours, managing troop allegiances, recalculating national territorial labels, and broadcasting perceptual updates to all BDI agents.

% -------------------------------------------------------------------------
\subsubsection{Territorial Transfer Mechanics: \texttt{conquer\_province}}
Territory reassignment is handled via \texttt{conquer\_province(province, agent, is\_riot)}. The method takes the target province, the recipient nation, and a boolean flag indicating whether the transfer was triggered by an internal civil rebellion or military capture:

\begin{enumerate}
    \item \textbf{Scene-Tree Re-Parenting and Color Realignment:} The province is re-parented to the conquering nation node (\texttt{province.reparent(agent)}). The surface polygon fill color (\texttt{province.fill.color}) is immediately updated based on the relationship between the new owner and the active spectator view.
    \item \textbf{Topological Boundary Regeneration:} Frontier lines must reflect the new sovereignty. The controller invokes \texttt{set\_borders()} on the transferred province and iteratively on all adjacent neighboring provinces, causing domestic boundaries to merge and foreign borders to render.
    \item \textbf{National Centroid and Label Realignment:} The area-weighted territorial centroids and label positions of both the losing power and the acquiring power are recalibrated via \texttt{update\_name\_position()}.
    \item \textbf{Infrastructure and Morale Transition:}
    \begin{itemize}
        \item If an administrative seat structure (\texttt{Enums.Buildings.CAPITAL}) was present, it is dismantled upon occupation (\texttt{buildings.erase(CAPITAL)}).
        \item \textbf{Military Conquest ($\text{is\_riot} = \text{false}$):} Provincial morale is reset to a depleted baseline (\texttt{CONQUERED\_PROVINCE\_MORALE}), and the \texttt{conquered} state flag is asserted to prevent immediate day-one healing.
        \item \textbf{Civil Rebellion ($\text{is\_riot} = \text{true}$):} Provincial morale is preserved without resetting. Furthermore, all stationary military units garrisoned at the provincial capital defect to the acquiring nation: their owning player reference (\texttt{army.player}) is reassigned to the recipient agent, transferring operational military assets as a direct consequence of the civil uprising.
    \end{itemize}
    \item \textbf{Global Perception Synchronization:} A \texttt{province\_owner\_update} telemetry frame is constructed and dispatched to every active nation in the game via \texttt{player.send\_vesna\_message(msg)}, updating their BDI Belief Bases regarding the new ownership, morale, defensive strength, and adjacent allied border configurations.
    \item \textbf{Victory Evaluation:} The controller triggers \texttt{check\_victory\_conditions()} to evaluate whether the territorial shift satisfies match completion thresholds.
\end{enumerate}

\begin{lstlisting}[caption={Structure of the province ownership update JSON telemetry frame.}, label={lst:province_owner_update_json}]
{
  "type": "sight",
  "sender": "body",
  "receiver": "PlayerName",
  "data": {
    "type": "province_owner_update",
    "id": "ProvinceName",
    "owner": "NewOwnerName",
    "morale": 50,
    "x": 140.0,
    "y": 220.0,
    "has_army": true,
    "defense_value": 40.0,
    "bordering_allies": ["AlliedProvinceA"],
    "is_capital": false,
    "is_double": true,
    "resource": "oil"
  }
}
\end{lstlisting}

\begin{lstlisting}[language=python, caption={Territorial transfer execution and state re-parenting in GameController.gd.}, label={lst:conquer_province_code}]
func conquer_province(province: Province, agent: Agent, is_riot: bool = false) -> void:
    var previous_player: Agent = province.get_parent()
    province.reparent(agent)

    _set_provinces_color()
    province.set_borders()
    for el in province.neighbors:
        if el is Province:
            el.set_borders()
            
    previous_player.update_name_position()
    agent.update_name_position()

    if Enums.Buildings.CAPITAL in province.buildings:
        province.buildings.erase(Enums.Buildings.CAPITAL)

    if not is_riot:
        province.province_morale = CONQUERED_PROVINCE_MORALE
        province.conquered = true
    else:
        for army in previous_player.armies:
            if army.global_position == province.get_city_global_position():
                army.player = agent
                previous_player.armies.erase(army)
                agent.armies.append(army)

    _broadcast_province_owner_update(province)
    check_victory_conditions()
\end{lstlisting}

% -------------------------------------------------------------------------
\subsubsection{Interactive Province Selection and HUD Coupling}
To support spectator inspection, the controller provides interactive selection routines that link visual feedback directly to the Heads-Up Display (HUD):

\begin{itemize}
    \item \textbf{Province Selection (\texttt{select\_province}):} Toggles the active selection state of a target province. When selected, the polygon mesh modulation is highlighted (\texttt{PROVINCE\_SELECTED\_MODULATION}), the provincial name label is forced visible (if not already displayed as a national capital), and detailed socioeconomic statistics (daily yield, morale modifiers, and active buildings) are sent to \texttt{hud.show\_province\_details(province)}. If the same province is clicked again, it deselects, restoring a neutral white mesh modulation and hiding the details panel.
    \item \textbf{Selection Clearing (\texttt{on\_clear\_province}):} Resets any active selection, restoring standard canvas item modulation and closing the provincial HUD panel.
\end{itemize}

% =========================================================================
\subsection{Master Game Loop, Victory Conditions, and Auxiliary Lifecycle Routines}
\label{subsec:game_controller_lifecycle}

To supervise global match execution, the \texttt{GameController} encapsulates master initialization routines, real-time leaderboard accounting, automated victory condition evaluation, temporal scaling, and match termination mechanics.

% -------------------------------------------------------------------------
\subsubsection{Match Initialization and Scene Graph Configuration: \texttt{\_ready}}
When the main game scene is loaded, the \texttt{\_ready()} virtual method coordinates system-wide initialization across all environmental and multi-agent subsystems:

\begin{enumerate}
    \item \textbf{Engine Temporal Scaling:} The engine's global execution speed is adjusted via \texttt{Engine.time\_scale = time\_multiplier}. This configuration enables high-speed headless training sessions and rapid benchmarking runs for BDI reasoning cycles.
    \item \textbf{Entity Discovery and Group Registration:} The controller populates its internal tracking structures by scanning node groups: \texttt{provinces} are cached into the \texttt{provinces} array, while sovereign players are identified from the \texttt{players} group.
    \item \textbf{Leaderboard Setup and Diplomatic Event Wiring:} Each registered player is initialized into the \texttt{leaderboard} array with their starting score points (\texttt{player.points}). The controller connects to every agent's diplomatic signals (\texttt{war\_declared}, \texttt{peace\_established}, \texttt{alliance\_established}) to trigger dynamic color and UI updates across all provincial meshes.
    \item \textbf{Capital Uniqueness Enforcement:} The controller sweeps each player's territorial nodes. If a nation is configured with multiple capital buildings across its starting provinces, redundant instances are stripped, enforcing the single-capital constraint.
    \item \textbf{HUD Layer and Fog of War Instantiation:} An isolated \texttt{CanvasLayer} is instantiated to host the UI interface (\texttt{hud.tscn}), linking spectator view-switching signals. Subsequently, \texttt{\_create\_fog()}, \texttt{\_fog\_map\_circles()}, and \texttt{\_set\_provinces\_color()} initialize the global sensory and diplomatic canvas.
\end{enumerate}

% -------------------------------------------------------------------------
\subsubsection{Leaderboard Scoring and Point Calculation}
The simulation tracks national dominance through a real-time leaderboard sorted in descending order of accumulated points via \texttt{\_sort\_leaderboard(a, b)}. 

Points are not arbitrary values; rather, they are computed dynamically by each \texttt{Agent} based on the domestic morale of its provinces:
\begin{equation}
    \text{Points}_{\text{nation}} = \sum_{P \in \text{OwnedProvinces}} f(M_P)
\end{equation}
where $M_P \in [0, 100]$ represents the current morale of province $P$, and $f(M_P)$ is a piecewise scoring function:
\begin{equation}
    f(M_P) = \begin{cases}
        0.0, & \text{if } M_P < 20.0 \\
        \min\left(10.0, \, \frac{M_P - 20.0}{8.0}\right), & \text{if } M_P \ge 20.0
    \end{cases}
\end{equation}

Under this formulation, each stabilized high-morale province ($M_P = 100$) contributes a maximum of $10$ points toward the national score. During every global simulation tick (\texttt{\_on\_global\_timer\_timeout()}), the controller synchronizes point values across the \texttt{leaderboard} array, re-sorts standings, updates dynamic Fog of War emitter positions, and evaluates victory criteria.

% -------------------------------------------------------------------------
\subsubsection{Victory Evaluation and Match Termination}
Match progression and termination are governed by \texttt{check\_victory\_conditions()}, \texttt{\_on\_global\_day\_timer\_timeout()}, and \texttt{end\_match()}:

\begin{itemize}
    \item \textbf{Score-Based Hegemony:} At every evaluation pass, \texttt{check\_victory\_conditions()} verifies whether any sovereign power has reached or exceeded the maximum victory threshold:
    \begin{equation}
        \text{Points}_{\text{nation}} \ge \text{MAX\_POINTS}, \quad \text{where } \text{MAX\_POINTS} = 200
    \end{equation}
    If a player achieves this score, \texttt{end\_match()} is triggered immediately.
    \item \textbf{Temporal Horizon Expiration:} At the end of each daily cycle, \texttt{day\_cycle} increments by one. If the elapsed days exceed the global lifespan limit:
    \begin{equation}
        \text{day\_cycle} \ge \text{game\_length}, \quad \text{where } \text{game\_length} = 30 \text{ simulation days}
    \end{equation}
    the match concludes by time expiration, and the power holding the highest leaderboard score is declared the winner.
    \item \textbf{Match Termination Handling (\texttt{end\_match}):} Upon termination, the master timers (\texttt{global\_timer}, \texttt{global\_hour\_timer}, and \texttt{global\_day\_timer}) are halted, freezing simulation time, and the \texttt{match\_over} signal is broadcast to all active agents and provincial subsystems.
\end{itemize}

\begin{lstlisting}[language=python, caption={Victory condition evaluation and match termination routines in GameController.gd.}, label={lst:victory_conditions}]
func check_victory_conditions() -> void:
    var players = get_tree().get_nodes_in_group("players")
    for player: Agent in players:
        if player.points >= MAX_POINTS:
            end_match()
            break

func _on_global_day_timer_timeout() -> void:
    day_cycle += 1
    if day_cycle > game_length:
        day_cycle = game_length
        end_match()

func end_match() -> void:
    if not global_timer.is_stopped():
        global_timer.stop()
        global_hour_timer.stop()
        global_day_timer.stop()
        emit_signal("match_over")
\end{lstlisting}

% -------------------------------------------------------------------------
\subsubsection{Geometric World Dimension Resolution: \texttt{get\_collision\_rect\_size}}
To configure the full-screen canvas dimensions required by the Fog of War rendering pipeline (\texttt{\_create\_fog()}), the controller utilizes the helper function \texttt{get\_collision\_rect\_size(poly)}.

The routine inspects the quad vertex coordinates of the background \texttt{Polygon2D} node, computing the absolute axis-aligned bounding box dimensions:
\begin{equation}
    \mathbf{S}_{\text{world}} = \big(\max(x) - \min(x), \; \max(y) - \min(y)\big)
\end{equation}
This vector defines the buffer size used by the engine to instantiate the dynamic \texttt{Image} texture that matches the world canvas.

% =========================================================================
\section{The Heads-Up Display (HUD) Interface}
\label{sec:hud_interface}

The \texttt{Hud} component serves as the primary visual interface presented to human observers and spectators. Structured as an overlay attached to an isolated \texttt{CanvasLayer}, it continuously translates internal game-state variables, multi-agent standings, temporal metrics, and provincial telemetry into an interactive, real-time dashboard.

% -------------------------------------------------------------------------
\subsection{HUD Architecture and Telemetry Overview}
\label{subsec:hud_overview}

The interface aggregates critical macro- and micro-level simulation metrics:
\begin{itemize}
    \item \textbf{Temporal Indicators:} Displays the current simulation day (\texttt{day\_cycle}) and provides controls to dynamically manipulate the engine's execution time scale (\texttt{Engine.time\_scale}).
    \item \textbf{Dynamic Leaderboard:} Renders the ranked standings of all active sovereign nations, continuously sorted by points accrued through provincial morale and territorial dominance.
    \item \textbf{The Gazette / Strategic Journal:} A multi-nation statistical overview presenting macro-demographic data, global military capacity, industrial commodity yields, and active diplomatic wars/alliances across the map.
    \item \textbf{Provincial Inspection Canvas:} Dedicated side-panels that render detailed socioeconomic statistics (daily commodity extraction, target morale formulas, active buildings, and defense ratings) when a territory is selected.
    \item \textbf{Sovereign View Switcher:} An interactive dropdown selector that allows observers to seamlessly switch between an omniscient spectator view and the subjective perspective of any specific nation.
\end{itemize}

% -------------------------------------------------------------------------
\subsection{Perspective Switching and Viewport Redirection}
\label{subsec:hud_view_selector}

A core visual feature of the HUD is the \texttt{agent\_selector} control, which allows observers to inspect the simulation from the perspective of any active nation at runtime. 

\begin{figure}[H]
    \centering
    \includegraphics[width=0.2\textwidth]{imgs/view_selector.png} 
    \caption{Graphical component of the view selector}
    \label{fig:view_selector}
\end{figure}

When an observer selects a nation from the dropdown, the handler \texttt{\_on\_agent\_selector\_item\_selected(index)} triggers the custom signal \texttt{view\_changed(target\_agent)}. This signal propagates directly to the \texttt{GameController} and the underlying canvas nodes, re-routing:
\begin{enumerate}
    \item \textbf{Fog of War Visibility:} Regenerates sensory circles via \texttt{\_fog\_map\_circles()}, displaying only the exact visual intelligence available to the chosen nation.
    \item \textbf{Diplomatic Color Realignment:} Updates provincial surface tints to reflect the selected nation's alliances (blue), active wars (red), and sovereign territory (green).
    \item \textbf{Resource Bar Redirection:} Binds the top-bar commodity counters (\texttt{MATERIAL\_TEXTURES}) directly to the selected agent's treasury reserves.
\end{enumerate}

\begin{lstlisting}[language=python, caption={Perspective selection and view switching logic in Hud.gd.}, label={lst:hud_view_selector}]
func _on_agent_selector_item_selected(index: int) -> void:
    if index == 0:
        view = null
        nation_details_panel.hide()
    else:
        var player_name: String = agent_selector.get_item_text(index)
        view = _get_agent_by_name(player_name)
        nation_details_panel.show()
        _update_nation_details(view)
        
    emit_signal("view_changed", view)
    _update_resource_display()
\end{lstlisting}

% -------------------------------------------------------------------------
\subsection{Nation Details Dashboard and Resource Formatting}
\label{subsec:hud_nation_details}

When the observer perspective is focused on a specific nation ($\texttt{view} \neq \text{null}$), the HUD displays a dedicated bottom statistics bar detailing the sovereign power's internal metrics.

\begin{figure}[H]
    \centering
    \includegraphics[width=1.0\textwidth]{imgs/nation_info.png} 
    \caption{Graphical component of the nation details}
    \label{fig:nation_info}
\end{figure}

The panel visualizes:
\begin{itemize}
    \item \textbf{Sovereign Identification:} Displays the official name and diplomatic flag of the active power.
    \item \textbf{Competitive Standing:} Displays the nation's current accumulated score points and its rank on the global leaderboard.
    \item \textbf{Economic Reserves Formatting:} Visualizes accumulated stockpiles across all strategic resources (Wood, Iron, Oil, Gas, Coal, Grain, Fish, and Gold). To ensure scannability without cluttering the UI canvas, numeric totals are parsed into compact metric suffixes (e.g., $15.4\text{k}$ for thousands, $1.2\text{M}$ for millions).
\end{itemize}

To guarantee real-time fidelity, the national details dashboard is refreshed via two complementary triggers:
\begin{enumerate}
    \item \textbf{On Perspective Change:} Refreshed immediately upon selecting a new agent via \texttt{\_on\_agent\_selector\_item\_selected()}.
    \item \textbf{On Global Timer Timeout:} Refreshed cyclically during \texttt{\_on\_timer\_timeout()}, ensuring that hourly economic deductions, commodity extraction yields, and leaderboard position shifts are updated on screen.
\end{enumerate}

\begin{lstlisting}[language=python, caption={National statistics and treasury formatting routine in Hud.gd.}, label={lst:hud_nation_details}]
func _update_nation_details(agent: Agent) -> void:
    if agent == null or not is_instance_valid(agent): return

    nation_name_label.text = agent.name
    score_label.text = "Points: %d (Rank #%d)" % [agent.points, _get_leaderboard_rank(agent)]

    for res in agent.resources:
        var amount: int = agent.resources[res]
        var formatted_text: String = _format_resource_amount(amount)
        resource_labels[res].text = formatted_text

func _format_resource_amount(val: int) -> String:
    if val >= 1_000_000:
        return "%.1fM" % (float(val) / 1_000_000.0)
    elif val >= 1_000:
        return "%.1fk" % (float(val) / 1_000.0)
    return str(val)
\end{lstlisting}

% =========================================================================
\subsection{Leaderboard Architecture and Comparative National Analytics}
\label{subsec:hud_leaderboard_journal}

To provide players and spectators with comparative strategic intelligence, the \texttt{Hud} incorporates a centralized leaderboard and a multi-dimensional statistical journal inspired by the gazette mechanics of \textit{Supremacy 1914}.

\begin{figure}[H]
    \centering
    \includegraphics[width=0.5\textwidth]{imgs/leaderboard.png} 
    \caption{Graphical component of the leaderboard}
    \label{fig:leaderboard}
\end{figure}

% -------------------------------------------------------------------------
\subsubsection{Dynamic Leaderboard Display and Camera Lock}
The global leaderboard ranks sovereign powers based on accrued victory points calculated from provincial stability. 

While point calculation and sorting routines are executed on the \texttt{GameController} backend during every global timer tick, the \texttt{Hud} handles the visual presentation of the standings:
\begin{itemize}
    \item \textbf{Decoupled Sorting and View Rendering:} The HUD reads the pre-sorted \texttt{leaderboard} array from the controller, populating a scrollable vertical container with rank badges, player names, and current point tallies.
    \item \textbf{Camera Pan Interception:} When the leaderboard modal window is opened on screen, camera panning and world-navigation inputs are suspended. This input lock prevents accidental map translations or zoom shifts while interacting with UI elements.
\end{itemize}

% -------------------------------------------------------------------------
\subsubsection{The Strategic Journal and Secondary Metric Rankings}
In addition to the primary victory score, the HUD generates three secondary comparative rankings at the end of each simulation day (\texttt{GlobalDayTimer}). These metrics provide agents and human observers with macro-level intelligence regarding foreign industrial capacities, military parity, and domestic stability.

\paragraph{1. Most Powerful Armies Index (\texttt{\_update\_most\_powerful\_armies})}
Calculates each nation's share of global military personnel. Let $N_{k}$ denote the total active soldiers fielded by agent $k$, and $N_{\text{global}} = \sum_j N_j$ the global military headcount:
\begin{equation}
    P_{\text{military}, k} = \left( \frac{N_k}{N_{\text{global}}} \right) \times 100\%
\end{equation}
The list ranks powers in descending order of troop percentages. Because this metric tracks headcount rather than division hit points, it provides an estimate of mobilization scale while obscuring the actual operational health of foreign forces.

\begin{lstlisting}[language=python, caption={Military strength aggregation and percentage ranking routine in Hud.gd.}, label={lst:hud_powerful_armies}]
func _update_most_powerful_armies() -> void:
    var nation_data: Array[Dictionary] = []
    var global_soldiers: int = 0
    
    for agent in agent_list:
        if not is_instance_valid(agent): continue
        var nation_soldiers: int = 0
        for army in agent.armies:
            if is_instance_valid(army):
                nation_soldiers += army.soldier_number
        global_soldiers += nation_soldiers
        nation_data.append({"agent": agent, "soldiers": nation_soldiers})
        
    nation_data.sort_custom(func(a, b): return a["soldiers"] > b["soldiers"])
    for entry in most_powerful_armies_list.get_children():
        entry.queue_free()
        
    for i in range(len(nation_data)):
        var entry = nation_data[i]
        var agent = entry["agent"]
        var soldiers = entry["soldiers"]
        var power: float = (float(soldiers) / global_soldiers * 100.0) if global_soldiers > 0 else 0.0
        
        var container: HBoxContainer = HBoxContainer.new()
        var label_name: Label = Label.new()
        var label_value: Label = Label.new()
        label_name.text = "%d - %s:" % [i + 1, agent.name.capitalize()]
        label_value.text = "%.1f%%" % power
        label_value.size_flags_horizontal = Control.SIZE_EXPAND_FILL
        label_value.horizontal_alignment = HORIZONTAL_ALIGNMENT_RIGHT
        container.add_child(label_name)
        container.add_child(label_value)
        most_powerful_armies_list.add_child(container)
\end{lstlisting}

\paragraph{2. Economic Productivity Index (\texttt{\_update\_best\_economies})}
Aggregates the total raw commodity extraction rate across all sovereign territories owned by an agent:
\begin{equation}
    Y_{\text{total}, k} = \sum_{P \in \text{Provinces}_k} \text{get\_daily\_resource\_count}(P)
\end{equation}
This index evaluates industrial expansion and raw resource throughput. Values are rendered using metric tonnage labels (e.g., $\text{t}$ or $\text{k t}$), allowing players to gauge foreign industrial strength without exposing specific stockpile compositions.

\begin{lstlisting}[language=python, caption={Economic commodity throughput aggregation in Hud.gd.}, label={lst:hud_best_economies}]
func _update_best_economies() -> void:
    var nation_data: Array[Dictionary] = []
    for agent in agent_list:
        if not is_instance_valid(agent): continue
        var daily_prod: int = 0
        for child in agent.get_children():
            if child is Province:
                daily_prod += child.get_daily_resource_count()
        nation_data.append({"agent": agent, "production": daily_prod})
        
    nation_data.sort_custom(func(a, b): return a["production"] > b["production"])
    for entry in best_economies_list.get_children():
        entry.queue_free()
        
    for i in range(len(nation_data)):
        var entry = nation_data[i]
        var agent = entry["agent"]
        var prod = entry["production"]
        var prod_str = "%.1fk t" % (prod / 1000.0) if prod > 999 else str(prod) + " t"
        
        var container: HBoxContainer = HBoxContainer.new()
        var label_name: Label = Label.new()
        var label_value: Label = Label.new()
        label_name.text = "%d - %s:" % [i + 1, agent.name.capitalize()]
        label_value.text = prod_str
        label_value.size_flags_horizontal = Control.SIZE_EXPAND_FILL
        label_value.horizontal_alignment = HORIZONTAL_ALIGNMENT_RIGHT
        container.add_child(label_name)
        container.add_child(label_value)
        best_economies_list.add_child(container)
\end{lstlisting}

\paragraph{3. National Stability Index (\texttt{\_update\_most\_popular\_nations})}
Evaluates domestic cohesion by computing the unweighted arithmetic mean of provincial morale across all territories controlled by an agent:
\begin{equation}
    \bar{M}_k = \frac{1}{|\text{Provinces}_k|} \sum_{P \in \text{Provinces}_k} M_P
\end{equation}
This ranking exposes internal unrest, war weariness, or territorial overextension, highlighting powers susceptible to civil riots or administrative collapse.

\begin{lstlisting}[language=python, caption={National morale and stability index calculation in Hud.gd.}, label={lst:hud_popular_nations}]
func _update_most_popular_nations() -> void:
    var nation_data: Array[Dictionary] = []
    for agent in agent_list:
        if not is_instance_valid(agent): continue
        var total_morale: int = 0
        var province_count: int = 0
        for child in agent.get_children():
            if child is Province:
                total_morale += child.province_morale
                province_count += 1
        var mean_morale: float = (float(total_morale) / province_count) if province_count > 0 else 0.0
        nation_data.append({"agent": agent, "morale": mean_morale})
        
    nation_data.sort_custom(func(a, b): return a["morale"] > b["morale"])
    for entry in most_popular_nations_list.get_children():
        entry.queue_free()
        
    for i in range(len(nation_data)):
        var entry = nation_data[i]
        var agent = entry["agent"]
        var morale = entry["morale"]
        
        var container: HBoxContainer = HBoxContainer.new()
        var label_name: Label = Label.new()
        var label_value: Label = Label.new()
        label_name.text = "%d - %s:" % [i + 1, agent.name.capitalize()]
        label_value.text = "%.1f%%" % morale
        label_value.size_flags_horizontal = Control.SIZE_EXPAND_FILL
        label_value.horizontal_alignment = HORIZONTAL_ALIGNMENT_RIGHT
        container.add_child(label_name)
        container.add_child(label_value)
        most_popular_nations_list.add_child(container)
\end{lstlisting}

\subsubsection{Temporal Updating Cadence}
These secondary rankings are recomputed exclusively at the start of each daily cycle (\texttt{GlobalDayTimer}) rather than continuously on every tick. This periodic update interval prevents cognitive BDI agents and human players from relying on high-frequency intelligence, simulating asynchronous reconnaissance and preserving strategic uncertainty.

\begin{figure}[H]
    \centering
    \includegraphics[width=1.0\textwidth]{imgs/journal.png} 
    \caption{Graphical component of the journal}
    \label{fig:journal}
\end{figure}

% =========================================================================
\subsection{Temporal Tracking and Dynamic Engine Time-Scale Modulation}
\label{subsec:hud_time_controls}

The final component of the HUD interface manages macro-temporal visibility and interactive simulation scaling, providing human observers and researchers with fine-grained control over match progression.

% -------------------------------------------------------------------------
\subsubsection{Day Counter and Match Horizon Indicator}
A dedicated on-screen day counter displays the current simulation cycle (\texttt{day\_cycle}), synchronized directly with the \texttt{GameController} master timer. 

This counter serves as an essential indicator of overall match progression against the global temporal limit ($T_{\text{max}} = 30\text{ simulation days}$). When the day counter reaches this threshold without any power attaining point hegemony, the match concludes automatically, awarding victory to the nation leading the leaderboard standings. Real-time day tracking allows observers to assess whether trailing BDI agents are shifting toward aggressive high-risk expansion or entrenching their positions as the temporal deadline approaches.

% -------------------------------------------------------------------------
\subsubsection{Dynamic Time-Scale Slider and Simulation Acceleration}
To facilitate both granular inspection and rapid multi-day evaluation, the HUD integrates an interactive linear time-scale slider bound to the engine's execution frequency multiplier:
\begin{equation}
    \text{Engine.time\_scale} \in [1.0\times, \, 50.0\times]
\end{equation}

This dynamic modulation provides dual experimental capabilities:
\begin{itemize}
    \item \textbf{Real-Time Tactical Inspection ($1.0\times$):} Standard baseline speed allows researchers and observers to analyze frame-by-frame kinematics, micro-combat damage exchanges, collision intercept triggers, and dynamic Fog of War shader transitions.
    \item \textbf{Accelerated Evaluation Runs (up to $50.0\times$):} Rapid execution bypasses prolonged idle periods during early macro-economic accumulation and troop transit. This high-speed capability is particularly critical for running empirical multi-agent benchmarks, testing long-term BDI diplomatic convergence, and conducting automated match simulations within compressed experimental timeframes.
\end{itemize}





% ==============================================================================
\section{Architecture of the Reasoning Engine and Global Knowledge Base}
\label{sec:impl_knowledge_base}

The cognitive knowledge base of the sovereign agent is founded upon a declarative belief base representing the geopolitical landscape, static simulation constants, dynamic accessors, and recursive list manipulation predicates. The agent reasons by continuously evaluating these logical predicates against incoming environmental perceptions.

\subsection{Parametric Game Rules and Dynamic Accessors}
\label{sec:impl_game_rules}

To ensure robustness against environmental packet loss and to establish baseline operational invariants, the agent initializes a tuple of global simulation constants in its initial belief base. These constants govern civil riot probabilities, combat resolution ratios, diplomatic alliance limits, and post-conquest territorial stabilization metrics.

\begin{lstlisting}[language=prolog, caption={Default game rules and extraction inference rules in \texttt{vesna.asl}.}, label={lst:game_rules_asl}]
target_limit(1).
active_wars(0).

// Default fallback game rules: RiotThreshold, RiotBreakpoint, RiotBaseDefense,
// RiotPenaltyMultiplier, LifePerUnit, AttackPerUnit, DefensePerUnit, MaxAlliances, ConqueredMorale
game_rules(33, 30, 5.0, 6.0, 20.0, 4.0, 5.0, 1, 25).

get_riot_threshold(Val)          :- game_rules(Val, _, _, _, _, _, _, _, _).
get_riot_breakpoint(Val)         :- game_rules(_, Val, _, _, _, _, _, _, _).
get_riot_base_defense(Val)        :- game_rules(_, _, Val, _, _, _, _, _, _).
get_riot_penalty_multiplier(Val) :- game_rules(_, _, _, Val, _, _, _, _, _).
get_life_per_unit(Val)           :- game_rules(_, _, _, _, Val, _, _, _, _).
get_attack_per_unit(Val)         :- game_rules(_, _, _, _, _, Val, _, _, _).
get_defense_per_unit(Val)        :- game_rules(_, _, _, _, _, _, Val, _, _).
get_max_alliances(Val)           :- game_rules(_, _, _, _, _, _, _, Val, _).
get_conquered_morale(Val)        :- game_rules(_, _, _, _, _, _, _, _, Val).
\end{lstlisting}

The foundational belief \texttt{target\_limit(1)} restricts the agent to focusing its offensive operations on at most one active campaign target province per enemy nation simultaneously. This constraint prevents force dispersion and ensures concentration of military power. The belief \texttt{active\_wars(0)} initializes the active war counter to zero.

The \texttt{game\_rules/9} predicate encapsulates the simulation parameters:
\begin{enumerate}[label=(\roman*)]
    \item $\theta_{\text{thresh}} = 33\%$: The critical provincial morale threshold below which civil riots become possible;
    \item $\theta_{\text{break}} = 30\%$: The severe morale breakpoint below which riot penalties escalate;
    \item $D_{\text{base}} = 5.0$: The baseline stationary defensive resistance required to suppress civil unrest;
    \item $\mu = 6.0$: The penalty multiplier applied to morale deficits below the breakpoint;
    \item $\text{HP}_{\text{unit}} = 20.0$: The hit point capacity comprising a single infantry soldier;
    \item $\alpha = 4.0$: The offensive combat attack strength delivered per soldier;
    \item $\delta = 5.0$: The defensive combat resistance provided per soldier;
    \item $\mathit{MAX\_ALLIANCES} = 1$: The maximum number of bilateral alliances the agent may maintain;
    \item $M_{\text{conq}} = 25\%$: The baseline morale percentage assigned to newly conquered enemy territory.
\end{enumerate}

To access these parameters cleanly without hardcoding positional indices throughout strategic plans, a series of accessor inference rules are defined. Each rule uses Jason's unification mechanism with anonymous placeholder variables (\texttt{\_}) to project the designated parameter into the output variable \texttt{Val}. When updated parameters arrive from Godot via environment configuration messages, the agent replaces the \texttt{game\_rules/9} belief, and all accessor rules automatically reflect the updated values.

\subsection{Logical List Processing Predicates}
\label{sec:impl_list_utils}

Because AgentSpeak operates over first-order terms and lists, candidate selection and set filtering rely on declarative list operations. The membership predicate \texttt{list\_member/2} is defined using recursive pattern matching:

\begin{lstlisting}[language=prolog, caption={List membership inference rule in \texttt{vesna.asl}.}, label={lst:list_utils_asl}]
list_member(X, [X | _]).
list_member(X, [_ | T]) :-
    list_member(X, T).
\end{lstlisting}

The base clause unifies successfully if the queried element \texttt{X} is identical to the head of the list. If the head does not match, the recursive clause discards the head and evaluates membership across the tail \texttt{T}. This predicate provides the foundation for validating whether candidate provinces border allied territories or whether an enemy is part of an active coalition.

% ==============================================================================
\section{Economic Deficit Analysis and Strategic Campaign Planning}
\label{sec:impl_economic_campaign}

The economic engine of the sovereign agent is designed to maintain self-sufficiency across all eight primary commodities: Wood, Iron, Oil, Gas, Coal, Grain, Fish, and Gold. When domestic consumption or army upkeep exceeds daily output, structural deficits arise, requiring the agent to identify vulnerable foreign territories capable of restoring economic equilibrium.

\subsection{Resource Trend Tracking}
\label{subsec:resource_tracking}

Resource telemetry is received periodically from Godot through event perceptions. The agent compares incoming commodity quantities against stored historical balances to determine the current trend direction.

\begin{lstlisting}[language=prolog, caption={Resource update handling and trend classification in \texttt{vesna.asl}.}, label={lst:resource_update_asl}]
+resource_update(Type, Amount) : resource(Type, OldAmount) & OldAmount > Amount
  <- -resource(Type, OldAmount);
     +resource(Type, Amount);
     -resource_trend(Type, _);
     +resource_trend(Type, negative).

+resource_update(Type, Amount) : resource(Type, OldAmount) & OldAmount <= Amount
  <- -resource(Type, OldAmount);
     +resource(Type, Amount);
     -resource_trend(Type, _);
     +resource_trend(Type, positive).

+resource_update(Type, Amount) : not resource(Type, _)
  <- +resource(Type, Amount);
     +resource_trend(Type, positive).
\end{lstlisting}

As detailed in Listing~\ref{lst:resource_update_asl}, when a \texttt{+resource\_update(Type, Amount)} perception arrives, the plan's context evaluates the previous balance \texttt{OldAmount}. If $\mathit{OldAmount} > \mathit{Amount}$, the agent retracts the previous conviction and asserts \texttt{+resource\_trend(Type, negative)}. Conversely, if the stock is stable or increasing, \texttt{+resource\_trend(Type, positive)} is asserted. If no prior record exists, the commodity is initialized with a positive trend.

\subsection{Target Nation Identification and Strategic Need}
\label{subsec:target_nation}

Once a negative trend is registered for a resource, the agent evaluates all adjacent sovereign nations to determine which neighbor represents the most lucrative target for military conquest.

\begin{lstlisting}[language=prolog, caption={Inference rules for target nation evaluation in \texttt{vesna.asl}.}, label={lst:target_nation_asl}]
plans_to_attack(Nation) :-
    resource_trend(NeededResource, negative) &
    .findall(Amount, neighbor_resource(_, NeededResource, Amount), ListOfAmounts) &
    ListOfAmounts \== [] &
    .max(ListOfAmounts, MaxAmount) &
    neighbor_resource(Nation, NeededResource, MaxAmount).

needs_war_with(Enemy) :-
    resource_trend(NeededResource, negative) &
    neighbor_resource(Enemy, NeededResource, Amount) &
    Amount > 0.
\end{lstlisting}

The inference rule \texttt{plans\_to\_attack(Nation)} formalizes the target nation selection:
\begin{equation}
    \label{eq:target_nation}
    \mathit{TargetNation} = \arg\max_{N \in \text{Neighbors}} \Big( \mathit{Output}(N, \mathit{Res}^*) \Big)
\end{equation}
The rule queries all neighboring nations producing \texttt{NeededResource} using the internal action \texttt{.findall}, identifies the maximal daily output \texttt{MaxAmount} via \texttt{.max}, and unifies \texttt{Nation} with the sovereign power producing that maximum yield.

The auxiliary rule \texttt{needs\_war\_with(Enemy)} returns true if nation \texttt{Enemy} produces a non-zero quantity of any commodity for which the agent currently experiences a deficit. This rule is utilized extensively during diplomatic negotiations to decide whether a war should be continued or concluded.

\subsection{Multi-Factor Province Priority Scoring}
\label{subsec:target_scoring}

Within the borders of the identified target nation, the agent must select a specific province to conquer. Rather than attacking randomly, candidate border provinces are evaluated using a multi-factor priority score.

\begin{lstlisting}[language=prolog, caption={Inference rules for target province scoring in \texttt{vesna.asl}.}, label={lst:target_scoring_asl}]
capital_offset(true, 30).
capital_offset(false, 0).

double_offset(true, 20).
double_offset(false, 0).

resource_offset(Res, 15) :- resource_trend(Res, negative).
resource_offset(Res, 0)  :- not resource_trend(Res, negative).

target_score(ProvId, EnemyNation, Score) :-
    enemy_province_status(ProvId, EnemyNation, _, DefenseValue, _, _, BorderingAllies, IsCapital, IsDouble, Resource) &
    capital_offset(IsCapital, CapOffset) &
    double_offset(IsDouble, DoubleOffset) &
    resource_offset(Resource, ResOffset) &
    Score = DefenseValue - CapOffset - DoubleOffset - ResOffset.
\end{lstlisting}

The heuristic priority score $\mathit{Score}(P)$ is formulated mathematically as:
\begin{equation}
    \label{eq:target_score}
    \mathit{Score}(P) = \mathit{DefenseValue}(P) - \Omega_{\text{cap}}(P) - \Omega_{\text{double}}(P) - \Omega_{\text{res}}(P)
\end{equation}
where strategic reward offsets are defined as:
\begin{align}
    \Omega_{\text{cap}}(P) &= \begin{cases} 30, & \text{if } P \text{ is the Enemy National Capital} \\ 0, & \text{otherwise} \end{cases} \\
    \Omega_{\text{double}}(P) &= \begin{cases} 20, & \text{if } P \text{ has Double Resource Yield} \\ 0, & \text{otherwise} \end{cases} \\
    \Omega_{\text{res}}(P) &= \begin{cases} 15, & \text{if } \mathit{Resource}(P) \text{ is a Needed Deficit Resource} \\ 0, & \text{otherwise} \end{cases}
\end{align}

The scoring function balances defensive difficulty against strategic utility. A weakly defended capital that produces double quantities of a required commodity yields a substantial negative score, making it an optimal target for immediate conquest.

\subsection{Campaign Evaluation and Target Deliberation Plans}
\label{subsec:campaign_plans}

During each cognitive audit cycle, the agent iterates over active resource deficits and deliberates on potential war declarations and province target selections.

\begin{lstlisting}[language=prolog, caption={Target evaluation and province selection plans in \texttt{vesna.asl}.}, label={lst:campaign_plans_asl}]
+!evaluate_targets : true
  <- for (resource_trend(NeededResource, negative)) {
         .findall(Amount, neighbor_resource(Neighbor, NeededResource, Amount), ListOfAmounts);
         if (ListOfAmounts \== []) {
             .max(ListOfAmounts, MaxAmount);
             ?neighbor_resource(BestNeighbor, NeededResource, MaxAmount);
             if (not at_war_with(BestNeighbor) & not allied_with(BestNeighbor)) {
                 !decide_attack(BestNeighbor, NeededResource);
             }
         }
     }.

+!select_war_target(EnemyNation, BestProv)
  <- .findall([Score, ProvId], (enemy_province_status(ProvId, EnemyNation, _, _, _, _, BorderingAllies, _, _, _) & BorderingAllies \== [] & not war_target(EnemyNation, ProvId) & target_score(ProvId, EnemyNation, Score)), Options);
     if (Options \== []) {
         .min(Options, [MinScore, BestProv]);
         +war_target(EnemyNation, BestProv);
         .print("Selected target province ", BestProv, " of ", EnemyNation, " with priority score ", MinScore);
     }.
\end{lstlisting}

In \texttt{+!evaluate\_targets}, the agent executes a \texttt{for} loop across all active negative resource trends. For each deficit, it identifies \texttt{BestNeighbor} and invokes the goal \texttt{!decide\_attack(BestNeighbor, NeededResource)}.

When a war is active, the goal \texttt{+!select\_war\_target(EnemyNation, BestProv)} evaluates all candidate enemy provinces that directly border allied territory (\texttt{BorderingAllies \textbackslash== []}) and are not currently assigned as active targets. The agent gathers all \texttt{[Score, ProvId]} pairs via \texttt{.findall}, selects the pair with the minimal score using \texttt{.min}, and asserts \texttt{+war\_target(EnemyNation, BestProv)}.

% ==============================================================================
\section{Multi-Front Risk Assessment and War Deliberation}
\label{sec:impl_war_deliberation}

Declaring war against a sovereign nation carries existential risks, particularly if the adversary belongs to a powerful coalition or if the agent is already engaged on multiple military fronts. The agent implements a multi-front risk audit to ensure that it only initiates conflict when strategically and numerically advantageous.

\subsection{Coalition Power Aggregation and Representative Election}
\label{subsec:coalition_rules}

To accurately evaluate the balance of power, the agent calculates the combined military strength of nations and their formal alliance networks.

\begin{lstlisting}[language=prolog, caption={Alliance symmetry and coalition power aggregation rules in \texttt{vesna.asl}.}, label={lst:coalition_rules_asl}]
are_allied(X, Y) :- allies(X, Y).
are_allied(X, Y) :- allies(Y, X).

allied_with(Nation) :-
    .my_name(Me) &
    are_allied(Me, Nation).

is_war_block_representative(N) :-
    at_war_with(N) &
    not (at_war_with(O) & are_allied(N, O) & O < N).

alliance_strength(Nation, TotalStrength) :-
    .findall(Member, (Nation = Member | are_allied(Nation, Member)), RawMembers) &
    .union(RawMembers, [], Members) &
    .findall(Strength, (.member(M, Members) & nation_strength(M, Strength)), Strengths) &
    TotalStrength = math.sum(Strengths).

total_enemy_strength(TotalStrength) :-
    .findall(EnemyStrength, (is_war_block_representative(Enemy) & alliance_strength(Enemy, EnemyStrength)), Strengths) &
    TotalStrength = math.sum(Strengths).
\end{lstlisting}

The predicate \texttt{are\_allied(X, Y)} establishes bidirectional symmetry over alliance relations, while \texttt{allied\_with(Nation)} identifies whether a foreign nation is formally allied with the agent.

When aggregating total enemy coalition power across multiple active conflicts, a critical challenge arises: if two enemy nations at war with the agent are allied with each other, summing their alliance networks independently would double-count their mutual strength. To prevent this distortion, the rule \texttt{is\_war\_block\_representative(N)} elects a unique representative for each allied enemy coalition block by enforcing a strict lexicographical ordering constraint (\texttt{O < N}).

The total alliance strength $S_{\text{alliance}}(N)$ is computed in \texttt{alliance\_strength/2} as:
\begin{equation}
    \label{eq:alliance_strength}
    S_{\text{alliance}}(N) = S(N) + \sum_{A \in \text{Allies}(N)} S(A)
\end{equation}
The rule \texttt{total\_enemy\_strength(TotalStrength)} aggregates the alliance strengths exclusively from the elected block representatives, yielding the exact confronting enemy threat.

\subsection{War Supportability and Decision Execution}
\label{subsec:war_decision}

The agent formalizes the multi-front supportability constraint within the \texttt{can\_support\_new\_war} rule and executes war declarations through \texttt{+!decide\_attack}.

\begin{lstlisting}[language=prolog, caption={War supportability check and war declaration plan in \texttt{vesna.asl}.}, label={lst:war_plans_asl}]
can_support_new_war(Me, TargetNation) :-
    alliance_strength(Me, MyStrength) &
    total_enemy_strength(CurrentTotalEnemyStrength) &
    alliance_strength(TargetNation, TargetStrength) &
    CurrentTotalEnemyStrength + TargetStrength <= MyStrength * 1.5.

+!decide_attack(TargetNation, NeededResource)
  : not at_war_with(TargetNation)
    & not allied_with(TargetNation)
    & not war_cooldown
    & active_wars(W) & W < 1 
    & .my_name(Me)
    & can_support_new_war(Me, TargetNation)
    & needs_war_with(TargetNation)
  <- .my_name(Me);
     if (alliance_strength(Me, MyStrengthPercentage) & alliance_strength(TargetNation, TargetStrengthPercentage)) {
         if (MyStrengthPercentage >= TargetStrengthPercentage) {
             vesna.via.declare_war(TargetNation);
             +at_war_with(TargetNation);
             +war_cooldown;
             if (day_cycle(CurrentDay)) {
                 CooldownEndDay = CurrentDay + 1;
             } else {
                 CooldownEndDay = 2;
             };
             +war_cooldown_end_day(CooldownEndDay);
             .print("Declared war on ", TargetNation, " (MyStrength: ", MyStrengthPercentage, "%, TargetStrength: ", TargetStrengthPercentage, "%). Cooldown ends day ", CooldownEndDay);
         } else {
             .print("Decided NOT to declare war on ", TargetNation, ": enemy is stronger.");
         }
     }.

+!decide_attack(TargetNation, NeededResource) <- true.
\end{lstlisting}

A proactive war declaration on $\mathit{TargetNation}$ is executed only if all of the following conditions are simultaneously met in the plan context:
\begin{enumerate}[label=\alph*)]
    \item \textbf{Diplomatic Neutrality:} The target is neither an active enemy (\texttt{not at\_war\_with}) nor a formal ally (\texttt{not allied\_with});
    \item \textbf{Cooldown Inactive:} The agent is not on war cooldown (\texttt{not war\_cooldown});
    \item \textbf{Front Limitation:} Active war front count satisfies $W < 1$;
    \item \textbf{Coalition Capacity Rule:} The projected combined enemy power does not exceed $150\%$ of the agent's alliance strength:
    \begin{equation}
        \label{eq:coalition_capacity}
        \mathit{TotalEnemyStrength} + S_{\text{alliance}}(\mathit{TargetNation}) \le 1.5 \times S_{\text{alliance}}(Me)
    \end{equation}
    \item \textbf{Self-Confidence Requirement:} $S_{\text{alliance}}(Me) \ge S_{\text{alliance}}(\mathit{TargetNation})$.
\end{enumerate}

Upon satisfying these conditions, the agent executes the Java internal action \texttt{vesna.via.declare\_war(TargetNation)}, asserts \texttt{+at\_war\_with(TargetNation)}, activates \texttt{+war\_cooldown}, and schedules the cooldown expiration for $\mathit{CurrentDay} + 1$. The fallback clause \texttt{+!decide\_attack(\_, \_) <- true.} guarantees safe termination if any condition fails.

% ==============================================================================
\section{Offensive Tactical Operations and Force Splitting}
\label{sec:impl_offensive_tactics}

The offensive tactical engine translates high-level campaign objectives into granular division movements, precision force splitting, and multi-army coordinated assaults.

\subsection{Undefended Target Capture and Precision Garrison Splitting}
\label{subsec:undefended_tactics}

When attacking an undefended enemy province ($\mathit{DefenseValue} = 0$), deploying an entire army division is wasteful and leaves domestic bases vulnerable. Furthermore, newly conquered provinces drop to $M_{\text{conq}} = 25\%$ morale, creating immediate civil riot risk. The agent calculates the exact garrison required to stabilize the province upon capture and splits only that necessary force.

\begin{lstlisting}[language=prolog, caption={Undefended target force splitting in \texttt{vesna.asl}.}, label={lst:undefended_split_asl}]
// Segment of +!evaluate_attack_opportunities for undefended targets
if (EnemyDefenseValue == 0) {
    .max(AvailableArmies, [MaxAttackValue, BestArmyId]);
    ?my_army(BestArmyId, _, _, BestLife, _, _);
    ?get_riot_breakpoint(RiotBreakpoint);
    ?get_riot_base_defense(RiotBaseDefense);
    ?get_riot_penalty_multiplier(RiotPenaltyMultiplier);
    ?get_life_per_unit(LifePerUnit);
    ?get_defense_per_unit(DefensePerUnit);
    ?get_conquered_morale(ConqueredMorale);
    BestSoldiers = BestLife / LifePerUnit;
    RequiredGarrison = (RiotBaseDefense + (RiotBreakpoint - ConqueredMorale) * RiotPenaltyMultiplier) / DefensePerUnit;
    SplitCount = RequiredGarrison div 1;
    if (BestSoldiers > SplitCount + 1) {
        vesna.via.split_army(BestArmyId, SplitCount, ProvId);
        +targeted_province(ProvId);
        .print("Split army ", BestArmyId, " to send ", SplitCount, " soldiers to target province ", ProvId);
    } else {
        vesna.via.move_army(BestArmyId, ProvId);
        +my_army_travelling(BestArmyId);
        +my_army_target(BestArmyId, ProvId);
        +targeted_province(ProvId);
        .print("Sent army ", BestArmyId, " from border to target province ", ProvId);
    }
}
\end{lstlisting}

As formalized in Listing~\ref{lst:undefended_split_asl}, the required garrison $G_{\text{conq}}$ for a newly captured province ($M_{\text{conq}} = 25\%$) is derived as:
\begin{equation}
    \label{eq:conquered_garrison}
    G_{\text{conq}} = \frac{D_{\text{base}} + (\theta_{\text{break}} - M_{\text{conq}}) \times \mu}{\delta} = \frac{5.0 + (30 - 25) \times 6.0}{5.0} = 7.0 \text{ soldiers}
\end{equation}

If the available border army possesses $\mathit{BestSoldiers} > 8$ ($G_{\text{conq}} + 1$), the agent executes \texttt{vesna.via.split\_army(BestArmyId, 7, ProvId)}. Exactly 7 soldiers march to capture and stabilize the target, while the remaining force stays at the origin base for territorial defense. If the army contains 8 or fewer soldiers, the full unit is ordered to advance via \texttt{vesna.via.move\_army}.

\subsection{Defended Target Operations and Single-Army Splitting}
\label{subsec:defended_tactics}

Against fortified enemy positions ($\mathit{DefenseValue} > 0$), the agent enforces a $120\%$ decisive superiority threshold to guarantee victory while minimizing casualties.

\begin{lstlisting}[language=prolog, caption={Defended target force splitting in \texttt{vesna.asl}.}, label={lst:defended_split_asl}]
// Segment of +!evaluate_attack_opportunities for defended targets
?get_life_per_unit(LifePerUnit);
?get_attack_per_unit(AttackPerUnit);
TargetStrength = EnemyDefenseValue * 1.2;
TargetSoldiers = (TargetStrength / AttackPerUnit) div 1 + 1;

.findall([Soldiers, MyId], (.member([AttackVal, MyId], AvailableArmies) & my_army(MyId, _, _, Life, _, _) & Soldiers = Life / LifePerUnit & Soldiers > TargetSoldiers + 2), SplittableArmies);
if (SplittableArmies \== []) {
    .max(SplittableArmies, [MaxSoldiers, BestArmyId]);
    SplitCount = TargetSoldiers;
    vesna.via.split_army(BestArmyId, SplitCount, ProvId);
    +targeted_province(ProvId);
    .print("Split army ", BestArmyId, " to send ", SplitCount, " soldiers (Atk: ", SplitCount * AttackPerUnit, ") to target province ", ProvId, " (Def: ", EnemyDefenseValue, ")");
} else {
    !try_coordinated_attack(ProvId, EnemyDefenseValue, AvailableArmies);
}
\end{lstlisting}

The required offensive power $\mathit{TargetStrength}$ and corresponding soldier count $\mathit{TargetSoldiers}$ are computed as:
\begin{align}
    \mathit{TargetStrength} &= 1.2 \times \mathit{EnemyDefenseValue} \\
    \mathit{TargetSoldiers} &= \left\lfloor \frac{\mathit{TargetStrength}}{\alpha} \right\rfloor + 1
\end{align}

The agent inspects all adjacent divisions that are not traveling and not guarding active riots. If an adjacent army contains $\mathit{Soldiers} > \mathit{TargetSoldiers} + 2$, the agent splits off exactly $\mathit{TargetSoldiers}$ to execute the strike.

\subsection{Coordinated Multi-Army Convergence and Inland Mobilization}
\label{subsec:coordinated_attack}

If no single adjacent division is large enough to achieve decisive superiority on its own, the agent initiates a coordinated multi-army assault by accumulating multiple surrounding forces.

\begin{lstlisting}[language=prolog, caption={Coordinated multi-army accumulation and inland mobilization in \texttt{vesna.asl}.}, label={lst:coordinated_attack_asl}]
+!try_coordinated_attack(ProvId, EnemyDefenseValue, AvailableArmies)
  <- TargetStrength = EnemyDefenseValue * 1.2;
     !accumulate_attackers(ProvId, TargetStrength, AvailableArmies, 0, []).

+!accumulate_attackers(ProvId, TargetStrength, [[Atk, ArmyId] | Rest], AccumulatedStrength, ChosenArmies)
  : AccumulatedStrength <= TargetStrength
  <- !accumulate_attackers(ProvId, TargetStrength, Rest, AccumulatedStrength + Atk, [ArmyId | ChosenArmies]).

+!accumulate_attackers(ProvId, TargetStrength, _, AccumulatedStrength, ChosenArmies)
  : AccumulatedStrength > TargetStrength
  <- .print("Sending coordinated attack to ", ProvId, " with armies ", ChosenArmies, " (Total Atk: ", AccumulatedStrength, ")");
     +targeted_province(ProvId);
     for (.member(ArmyId, ChosenArmies)) {
         vesna.via.move_army(ArmyId, ProvId);
         +my_army_travelling(ArmyId);
         +my_army_target(ArmyId, ProvId);
     }.

+!accumulate_attackers(ProvId, TargetStrength, [], AccumulatedStrength, ChosenArmies)
  : AccumulatedStrength <= TargetStrength
  <- .print("Adjacent armies insufficient (", AccumulatedStrength, " vs ", TargetStrength, "). Mobilizing inland armies...");
     if (enemy_province_status(ProvId, _, _, _, _, _, BorderingAllies, _, _, _)) {
         .findall([Dist, MyId], (my_army(MyId, _, _, _, _, Loc) & not .member(Loc, BorderingAllies) & .member(AllyProv, BorderingAllies) & my_army_path_distance(MyId, AllyProv, Dist) & not is_defending_riot_province(MyId) & not my_army_travelling(MyId)), InlandArmies);
         if (InlandArmies \== []) {
             .min(InlandArmies, [MinDist, BestArmyId]);
             vesna.via.move_army(BestArmyId, ProvId);
             +my_army_travelling(BestArmyId);
             +my_army_target(BestArmyId, ProvId);
             +targeted_province(ProvId);
             .print("Sent closest inland army ", BestArmyId, " (distance: ", MinDist, ") to attack ", ProvId);
         }
     }.
\end{lstlisting}

The recursive plan \texttt{+!accumulate\_attackers} operates as follows:
\begin{enumerate}[label=\alph*)]
    \item \textbf{Recursive Accumulation:} While $\mathit{AccumulatedStrength} \le \mathit{TargetStrength}$, the plan pops \texttt{[Atk, ArmyId]} from the available list, adds \texttt{Atk} to the total, appends \texttt{ArmyId} to \texttt{ChosenArmies}, and recurses on the remainder of the list.
    \item \textbf{Coordinated Launch:} As soon as $\mathit{AccumulatedStrength} > \mathit{TargetStrength}$, the base plan iterates over \texttt{ChosenArmies} with a \texttt{for} loop, issuing simultaneous \texttt{vesna.via.move\_army(ArmyId, ProvId)} commands so that all selected divisions converge on the target at once.
    \item \textbf{Inland Fallback:} If all adjacent armies are exhausted and cumulative strength remains below $\mathit{TargetStrength}$, the agent queries non-border inland armies using \texttt{my\_army\_path\_distance}, selects the closest unit ($\min \mathit{Dist}$), and orders it to advance toward the front line to reinforce the staging base.
\end{enumerate}

% ==============================================================================
\section{Defensive Threat Interception Hierarchy}
\label{sec:impl_defense_hierarchy}

When hostile armies march toward domestic territory, the agent intercepts the invading forces before they reach provincial capital seats.

\subsection{Threat Detection and Defensive Garrison Sizing}
\label{subsec:threat_sizing}

The presence of an incoming hostile force is perceived through the \texttt{+foreign\_army\_moving} event.

\begin{lstlisting}[language=prolog, caption={Defensive threat detection and required force calculation in \texttt{vesna.asl}.}, label={lst:threat_detection_asl}]
province_under_threat(ProvId) :-
    enemy_army_moving(_, Nation, _, _, _, _, ProvId, _) &
    at_war_with(Nation).

+foreign_army_moving(ArmyId, Nation, PosX, PosY, NextX, NextY, TargetProv, AttackValue)
  <- .abolish(enemy_army_moving(ArmyId, Nation, _, _, _, _, _, _));
     if (TargetProv \== none) {
         +enemy_army_moving(ArmyId, Nation, PosX, PosY, NextX, NextY, TargetProv, AttackValue);
     };
     if (at_war_with(Nation) & TargetProv \== none & allied_province(TargetProv, _, _, _, _, _, _)) {
         .print("Enemy army ", ArmyId, " of ", Nation, " moving towards ", TargetProv, " (Atk: ", AttackValue, ")");
         if (not ordered_to_defend(_, TargetProv)) {
             ?get_life_per_unit(LifePerUnit);
             ?get_defense_per_unit(DefensePerUnit);
             RequiredSoldiers = (AttackValue * 1.2) / DefensePerUnit;
             SplitCount = RequiredSoldiers div 1 + 1;
             // Dispatch intercepting forces (Safe or Riot armies)...
         }
     };
     !evaluate_defenses.
\end{lstlisting}

The inference rule \texttt{province\_under\_threat(ProvId)} confirms whether the destination province is domestic and the marching force belongs to a nation at war. The required defensive garrison $\mathit{RequiredSoldiers}$ is calculated to provide $120\%$ resistance against the incoming attack power:
\begin{equation}
    \label{eq:required_defense}
    \mathit{RequiredSoldiers} = \left\lfloor \frac{1.2 \times \mathit{AttackValue}}{\delta} \right\rfloor + 1
\end{equation}

\subsection{Tiered Interception Selection and Defensive Release}
\label{subsec:tiered_interception}

To fulfill the defense requirement while preserving domestic civil order, the agent employs a two-tiered priority hierarchy.

\begin{lstlisting}[language=prolog, caption={Tiered unit selection and defense release in \texttt{vesna.asl}.}, label={lst:tiered_defense_asl}]
.findall([PathDist, MyId], (my_army(MyId, _, _, _, _, Loc) & my_army_path_distance(MyId, TargetProv, PathDist) & Loc \== TargetProv & not ordered_to_defend(MyId, _) & not is_defending_riot_province(MyId) & not my_army_travelling(MyId)), SafeArmies);
.findall([PathDist, MyId], (my_army(MyId, _, _, _, _, Loc) & my_army_path_distance(MyId, TargetProv, PathDist) & Loc \== TargetProv & not ordered_to_defend(MyId, _) & is_defending_riot_province(MyId) & not my_army_travelling(MyId)), RiotArmies);

if (SafeArmies \== []) {
    .min(SafeArmies, [MinDist, BestArmyId]);
    ?my_army(BestArmyId, _, _, Life, _, _);
    Soldiers = Life / LifePerUnit;
    if (Soldiers > SplitCount + 2) {
        vesna.via.split_army(BestArmyId, SplitCount, TargetProv);
        +ordered_to_defend(BestArmyId, TargetProv);
        +targeted_province(TargetProv);
    } else {
        +ordered_to_defend(BestArmyId, TargetProv);
        vesna.via.move_army(BestArmyId, TargetProv);
        +my_army_travelling(BestArmyId);
        +my_army_target(BestArmyId, TargetProv);
        +targeted_province(TargetProv);
    }
} else if (RiotArmies \== []) {
    // Fallback to riot-guarding armies if no safe armies exist...
    .min(RiotArmies, [MinDist, BestArmyId]);
    // Split or move BestArmyId as last resort...
}.

+!evaluate_defenses : true
  <- for (ordered_to_defend(MyId, ProvId)) {
         if (not province_under_threat(ProvId)) {
             .abolish(ordered_to_defend(MyId, ProvId));
             -my_army_target(MyId, ProvId);
             -my_army_travelling(MyId);
             .print("Threat to ", ProvId, " is gone. Released defending army ", MyId);
         }
     }.
\end{lstlisting}

The selection process evaluates:
\begin{itemize}
    \item \textbf{Tier 1 (Safe Armies):} Armies stationed in stable provinces ($M \ge 33\%$) that are not currently traveling or defending other threats. The agent identifies the unit minimizing graph path distance $\min \mathit{Dist}(A, \mathit{TargetProv})$. If the unit has surplus soldiers ($\mathit{Soldiers} > \mathit{SplitCount} + 2$), it splits off $\mathit{SplitCount}$ troops; otherwise, the full unit is deployed.
    \item \textbf{Tier 2 (Riot-Guarding Fallback):} If no safe armies exist across the nation, the agent queries units stationed at low-morale bases as an emergency measure, selecting the closest unit to intercept the existential threat.
    \item \textbf{Defense Release Cycle (\texttt{+!evaluate\_defenses}):} During every audit, the agent iterates across active \texttt{ordered\_to\_defend(MyId, ProvId)} convictions. If \texttt{not province\_under\_threat(ProvId)} (e.g., the enemy army changed course or was destroyed), the defensive order is retracted, immediately liberating the division for offensive or border-garrison duties.
\end{itemize}

% ==============================================================================
\section{Socio-Political Stability and Dynamic Riot Garrisoning}
\label{sec:impl_riot_control}

Provinces with low morale ($M < 33\%$) suffer from civil unrest and risk spontaneous rebellion. The agent stabilizes domestic territory by maintaining exact defensive garrisons and dynamically releasing surplus troops as morale recovers.

\begin{lstlisting}[language=prolog, caption={Riot defense predicate and garrison optimization plan in \texttt{vesna.asl}.}, label={lst:riot_plans_asl}]
is_defending_riot_province(MyId) :-
    my_army(MyId, _, _, _, _, Loc) &
    allied_province(Loc, Morale, _, _, _, _, _) &
    get_riot_threshold(RiotThreshold) &
    Morale < RiotThreshold.

is_defending_riot_province(MyId) :-
    ordered_to_defend(MyId, DefProv) &
    allied_province(DefProv, Morale, _, _, _, _, _) &
    get_riot_threshold(RiotThreshold) &
    Morale < RiotThreshold.

+!optimize_riot_garrisons : true
  <- ?get_riot_threshold(RiotThreshold);
     ?get_riot_breakpoint(RiotBreakpoint);
     ?get_riot_base_defense(RiotBaseDefense);
     ?get_riot_penalty_multiplier(RiotPenaltyMultiplier);
     ?get_life_per_unit(LifePerUnit);
     ?get_defense_per_unit(DefensePerUnit);
     for (my_army(MyId, _, _, Life, _, Loc)) {
         if (allied_province(Loc, Morale, _, _, _, _, _) & Morale < RiotThreshold) {
             if (Morale >= RiotBreakpoint) {
                 Required = RiotBaseDefense / DefensePerUnit;
             } else {
                 Required = (RiotBaseDefense + (RiotBreakpoint - Morale) * RiotPenaltyMultiplier) / DefensePerUnit;
             };
             Soldiers = Life / LifePerUnit;
             if (Soldiers > Required + 2) {
                 SplitCount = (Soldiers - Required) div 1;
                 !find_neighboring_allied_province(Loc, TargetProv);
                 if (TargetProv \== none) {
                     vesna.via.split_army(MyId, SplitCount, TargetProv);
                     .print("Split army ", MyId, " at low-morale province ", Loc, ": keeping ", Required, " soldiers, sending ", SplitCount, " to allied ", TargetProv);
                 }
             }
         }
     }.

+!find_neighboring_allied_province(Loc, TargetProv)
  <- ?allied_province(Loc, _, LocX, LocY, _, _, _);
     .findall([Dist, ProvId], (allied_province(ProvId, _, PX, PY, _, _, _) & ProvId \== Loc & Dist = (LocX - PX)**2 + (LocY - PY)**2), AlliedDists);
     if (AlliedDists \== []) {
         .min(AlliedDists, [MinDist, TargetProv]);
     } else {
         TargetProv = none;
     }.
\end{lstlisting}

The inference rule \texttt{is\_defending\_riot\_province(MyId)} flags any division stationed at or moving to a low-morale base, protecting it from being siphoned off for offensive or border-garrison tasks.

In \texttt{+!optimize\_riot\_garrisons}, the required stabilization garrison $G(M)$ is computed as:
\begin{equation}
    \label{eq:riot_formula}
    G(M) = \begin{cases}
    0, & \text{if } M \ge 33\% \\
    \frac{D_{\text{base}}}{\delta} = 1.0 \text{ soldier}, & \text{if } 30\% \le M < 33\% \\
    \frac{D_{\text{base}} + (\theta_{\text{break}} - M) \times \mu}{\delta}, & \text{if } M < 30\%
    \end{cases}
\end{equation}

If an army stationed at a low-morale territory contains $\mathit{Soldiers} > G(M) + 2$, the surplus $\mathit{SplitCount} = \lfloor \mathit{Soldiers} - G(M) \rfloor$ is split off via \texttt{vesna.via.split\_army} and deployed to reinforce the nearest domestic neighbor identified by \texttt{!find\_neighboring\_allied\_province}. As provincial morale recovers over time, $G(M)$ shrinks, automatically releasing garrison forces into active operational service.

% ==============================================================================
\section{Border Security and Deficit-Based Force Distribution}
\label{sec:impl_border_distribution}

To deter opportunistic foreign invasions, the agent continuously audits garrison strength along all non-allied frontiers and redistributes idle military forces to eliminate defensive deficits.

\subsection{Border Identification and Deficit Formulation}
\label{subsec:border_deficit}

An allied province is classified as a frontier post if it directly adjoins at least one enemy territory owned by a non-allied sovereign power.

\begin{lstlisting}[language=prolog, caption={Border classification, garrison requirements, and deficit rules in \texttt{vesna.asl}.}, label={lst:border_rules_asl}]
is_border_province(AlliedProv) :-
    allied_province(AlliedProv, _, _, _, _, _, _) &
    enemy_province_status(_, Owner, _, _, _, _, BorderingAllies, _, _, _) &
    .member(AlliedProv, BorderingAllies) &
    not allied_with(Owner).

nation_border_defense(BorderProv, Owner, SumDef) :-
    enemy_province_status(_, Owner, _, _, _, _, BorderingAllies, _, _, _) & 
    .member(BorderProv, BorderingAllies) & 
    not allied_with(Owner) &
    .findall(EnemyDef, (enemy_province_status(FProv, Owner, _, EnemyDef, _, _, BAllies, _, _, _) & .member(BorderProv, BAllies)), EnemyDefs) &
    SumDef = math.sum(EnemyDefs).

border_province_max_enemy_defense(BorderProv, MaxEnemyDef) :-
    enemy_province_status(_, Owner, _, _, _, _, BorderingAllies, _, _, _) & .member(BorderProv, BorderingAllies) & not allied_with(Owner) &
    .findall(SumDef, (enemy_province_status(_, Own, _, _, _, _, BAllies, _, _, _) & .member(BorderProv, BAllies) & not allied_with(Own) & nation_border_defense(BorderProv, Own, SumDef)), SumDefs) &
    .max(SumDefs, MaxEnemyDef).

target_garrison(Prov, Target) :-
    is_border_province(Prov) &
    border_province_max_enemy_defense(Prov, EnemyDef) &
    get_defense_per_unit(DefensePerUnit) &
    EnemyDef > 0 &
    Target = (EnemyDef * 1.2 / DefensePerUnit) div 1 + 1.

target_garrison(Prov, 1) :-
    is_border_province(Prov) &
    border_province_max_enemy_defense(Prov, EnemyDef) &
    EnemyDef <= 0.

target_garrison(Prov, 0) :-
    not is_border_province(Prov).

province_garrison_deficit(Prov, Deficit) :-
    province_soldiers(Prov, OurSoldiers) &
    target_garrison(Prov, Target) &
    OurSoldiers < Target &
    Deficit = Target - OurSoldiers.

province_excess_soldiers(Prov, Excess) :-
    province_soldiers(Prov, OurSoldiers) &
    target_garrison(Prov, Target) &
    OurSoldiers > Target &
    Excess = OurSoldiers - Target.
\end{lstlisting}

The target border garrison $\mathit{TargetGarrison}(P_{\text{border}})$ is formulated as:
\begin{equation}
    \label{eq:border_target}
    \mathit{TargetGarrison}(P_{\text{border}}) = \begin{cases}
    \left\lfloor \frac{1.2 \times \mathit{MaxEnemyDef}(P_{\text{border}})}{\delta} \right\rfloor + 1, & \text{if } \mathit{MaxEnemyDef} > 0 \\
    1 \text{ soldier}, & \text{if } \mathit{MaxEnemyDef} \le 0 \\
    0, & \text{if non-border}
    \end{cases}
\end{equation}

For each territory, the garrison deficit $\mathit{Deficit}(P)$ and surplus $\mathit{Excess}(P)$ are formulated as:
\begin{align}
    \mathit{Deficit}(P) &= \max\Big(0, \mathit{TargetGarrison}(P) - \mathit{OurSoldiers}(P)\Big) \\
    \mathit{Excess}(P) &= \max\Big(0, \mathit{OurSoldiers}(P) - \mathit{TargetGarrison}(P)\Big)
\end{align}

\subsection{Idle Force Redistribution Algorithm}
\label{subsec:border_redistribution}

During each audit, the agent identifies all idle military divisions (stationary, non-traveling, not ordered to defend active incursions, not guarding low-morale bases, and without active targets) and redistributes forces.

\begin{lstlisting}[language=prolog, caption={Border garrison redistribution plan in \texttt{vesna.asl}.}, label={lst:border_redistribution_asl}]
+!garrison_non_allied_borders : true
  <- ?get_life_per_unit(LifePerUnit);
     for (my_army(MyId, _, _, Life, _, Loc)) {
         if (not my_army_travelling(MyId) & 
             not ordered_to_defend(MyId, _) & 
             not is_defending_riot_province(MyId) & 
             not my_army_target(MyId, _)) {
             
             ?province_excess_soldiers(Loc, Excess);
             if (Excess > 0) {
                 .findall([Score, BorderProv], (is_border_province(BorderProv) & 
                                                province_garrison_deficit(BorderProv, Deficit) & 
                                                Deficit > 0 & 
                                                my_army_path_distance(MyId, BorderProv, Dist) & 
                                                Dist < 999999.0 & 
                                                Score = -Deficit + (Dist * 0.05)), OptionList);
                 if (OptionList \== []) {
                     .min(OptionList, [MinScore, SelectedProv]);
                     ?province_garrison_deficit(SelectedProv, DeficitVal);
                     ReqAmount = math.min(Excess, DeficitVal);
                     ArmySoldiers = Life / LifePerUnit;
                     Amount = math.min(ReqAmount, ArmySoldiers);
                     SplitCount = Amount div 1;
                     if (SplitCount > 0) {
                         if (SplitCount < ArmySoldiers) {
                             vesna.via.split_army(MyId, SplitCount, SelectedProv);
                             .print("Split idle army ", MyId, " sending ", SplitCount, " to deficit border ", SelectedProv);
                         } else {
                             vesna.via.move_army(MyId, SelectedProv);
                             +my_army_travelling(MyId);
                             +my_army_target(MyId, SelectedProv);
                             .print("Moving full idle army ", MyId, " to deficit border ", SelectedProv);
                         }
                     }
                 }
             } else if (not is_border_province(Loc)) {
                 .findall([Score, BorderProv], (is_border_province(BorderProv) & 
                                                my_army_path_distance(MyId, BorderProv, Dist) & 
                                                Dist < 999999.0 & 
                                                province_soldiers(BorderProv, OurSoldiers) & 
                                                Score = OurSoldiers + (Dist * 0.05)), AlliedOptionList);
                 if (AlliedOptionList \== []) {
                     .min(AlliedOptionList, [MinScore, SelectedProv]);
                     vesna.via.move_army(MyId, SelectedProv);
                     +my_army_travelling(MyId);
                     +my_army_target(MyId, SelectedProv);
                     .print("Moving idle inland army ", MyId, " to safe border ", SelectedProv);
                 }
             }
         }
     }.
\end{lstlisting}

The redistribution algorithm handles two spatial configurations:
\begin{enumerate}[label=\alph*)]
    \item \textbf{Redistribution from Excess Stations:} If an idle army is stationed in a base with $\mathit{Excess} > 0$, the agent evaluates all reachable border posts having $\mathit{Deficit} > 0$ using the cost metric:
    \begin{equation}
        \label{eq:deficit_score}
        \mathit{Score}(P_{\text{border}}) = -\mathit{Deficit}(P_{\text{border}}) + 0.05 \times \mathit{Dist}(\mathit{Army}, P_{\text{border}})
    \end{equation}
    The agent selects the post minimizing $\mathit{Score}$ (which prioritizes large deficits while penalizing excessive transit distance) and deploys $\min(\mathit{Excess}, \mathit{Deficit})$ soldiers via \texttt{split\_army} or \texttt{move\_army}.
    \item \textbf{Redistribution from Inland Positions:} If an idle army is located in a non-border inland province, the agent identifies the closest border post minimizing $\mathit{OurSoldiers} + 0.05 \times \mathit{Dist}$ and orders the entire inland division to relocate to the frontier.
\end{enumerate}

% ==============================================================================
\section{Diplomatic Bargaining and Coalition Protocols}
\label{sec:impl_diplomacy}

The diplomatic module enables the agent to negotiate bilateral peace treaties and establish strategic alliance coalitions to manage multi-front exposure.

\subsection{Bilateral Peace Negotiation Logic}
\label{subsec:peace_negotiation}

Peace negotiations operate under both reactive and proactive triggers.

\begin{lstlisting}[language=prolog, caption={Peace request processing and proactive peace proposals in \texttt{vesna.asl}.}, label={lst:peace_plans_asl}]
+peace_request_received(Sender)
  <- .my_name(Me);
     ?total_enemy_strength(Me, TotalEnemyStrength);
     ?alliance_strength(Me, MyStrength);
     
     if (TotalEnemyStrength > MyStrength * 1.5) {
         .print("Accepting peace from ", Sender, ": Total enemy strength too high (", TotalEnemyStrength, " vs ", MyStrength * 1.5, ")");
         vesna.via.accept_peace(Sender);
     } else {
         .findall(Rep, is_war_block_representative(Rep), RepList);
         .length(RepList, WarCount);
         ?alliance_strength(Sender, SenderStrength);
         if (WarCount > 1 & MyStrength - 10 < SenderStrength) {
             .print("Accepting peace from ", Sender, ": Multi-front exposure (WarCount=", WarCount, ")");
             vesna.via.accept_peace(Sender);
         } else {
             if (not needs_war_with(Sender)) {
                 .print("Accepting peace from ", Sender, ": Resource objectives fulfilled");
                 vesna.via.accept_peace(Sender);
             } else {
                 .print("Declining peace from ", Sender, ": Deficit resources still required");
                 vesna.via.decline_peace(Sender);
             }
         }
     }.

+!check_peace_requests : true
  <- .my_name(Me);
     ?total_enemy_strength(Me, TotalEnemyStrength);
     ?alliance_strength(Me, MyStrength);
     if (TotalEnemyStrength > MyStrength * 1.5) {
         .findall(Enemy, at_war_with(Enemy), EnemyList);
         !choose_enemy_for_peace(EnemyList, EnemyToRequest);
         if (EnemyToRequest \== none) {
             !send_peace_request_to_all(EnemyToRequest, "total enemy strength too high");
         }
     } else {
         for (at_war_with(Enemy)) {
             if (not needs_war_with(Enemy) & not sent_peace_request(Enemy)) {
                 !send_peace_request_to_all(Enemy, "resources not needed");
             } else if (alliance_strength(Enemy, EnemyStrengthPercentage)) {
                 if (EnemyStrengthPercentage > MyStrength * 1.5 & not sent_peace_request(Enemy)) {
                     !send_peace_request_to_all(Enemy, "enemy is a lot stronger");
                 }
             }
         }
     }.

+!choose_enemy_for_peace([], none).
+!choose_enemy_for_peace(EnemyList, EnemyToRequest)
  : EnemyList \== []
  <- .findall([Priority, MinusStrength, Enemy], 
              (.member(Enemy, EnemyList) & 
               get_enemy_peace_priority(Enemy, Priority) & 
               alliance_strength(Enemy, Strength) & 
               MinusStrength = -Strength), Options);
     if (Options \== []) {
         .min(Options, [_, _, EnemyToRequest]);
     } else {
         EnemyToRequest = none;
     }.
\end{lstlisting}

\subsubsection{Decision Breakdown}
\begin{enumerate}[label=\alph*)]
    \item \textbf{Incoming Peace Requests (\texttt{+peace\_request\_received}):} When a peace proposal is received from \texttt{Sender}, the agent evaluates acceptance through a prioritized decision tree:
    \begin{itemize}
        \item \textit{Emergency Acceptance:} If $\mathit{TotalEnemyStrength} > 1.5 \times \mathit{MyStrength}$, peace is accepted immediately via \texttt{vesna.via.accept\_peace(Sender)} to avoid defeat.
        \item \textit{Multi-Front Relief:} If active war count $\mathit{WarCount} > 1$ and $\mathit{MyStrength} - 10 < \mathit{SenderStrength}$, peace is accepted to consolidate against the primary front.
        \item \textit{Resource Saturation:} If the agent no longer requires resources from the sender ($\neg \text{needs\_war\_with}(\mathit{Sender})$), peace is accepted.
        \item \textit{Decline:} If none of these hold and resources are still needed, the proposal is declined via \texttt{vesna.via.decline\_peace(Sender)}.
    \end{itemize}
    \item \textbf{Proactive Peace Proposals (\texttt{+!check\_peace\_requests}):} If confronting excessive enemy coalitions ($> 1.5 \times \mathit{MyStrength}$), the agent executes \texttt{!choose\_enemy\_for\_peace}, which ranks active enemies by $\min(\mathit{Priority}, -\mathit{Strength}, \mathit{Enemy})$ where $\mathit{Priority} = 0$ for non-essential adversaries and $-\mathit{Strength}$ targets the strongest opposing bloc first.
\end{enumerate}

\subsection{Bilateral Alliance Formation Protocol}
\label{subsec:alliance_formation}

The agent maintains a strict alliance ceiling ($\mathit{MAX\_ALLIANCES} = 1$) to avoid conflicting coalition obligations.

\begin{lstlisting}[language=prolog, caption={Alliance request handling and proactive alliance proposals in \texttt{vesna.asl}.}, label={lst:alliance_plans_asl}]
+alliance_request_received(Sender)
  <- ?get_max_alliances(MaxAlliances);
     .count(allied_with(_), AllyCount);
     .count(sent_alliance_request(_), SentCount);
     .count(accepting_alliance(_), AcceptingCount);
     if (AllyCount < MaxAlliances & SentCount < MaxAlliances & AcceptingCount < MaxAlliances & not plans_to_attack(Sender) & not at_war_with(Sender) & not allied_with(Sender)) {
         +accepting_alliance(Sender);
         .print("Accepting alliance request from ", Sender);
         vesna.via.accept_alliance(Sender);
     } else {
         .print("Declining alliance request from ", Sender);
         vesna.via.decline_alliance(Sender);
     }.

+!check_alliance_opportunities : true
  <- ?get_max_alliances(MaxAlliances);
     .count(allied_with(_), AllyCount);
     .count(sent_alliance_request(_), SentCount);
     .count(accepting_alliance(_), AcceptingCount);
     if (AllyCount < MaxAlliances & SentCount < MaxAlliances & AcceptingCount < MaxAlliances) {
         .findall(Neighbor, (neighbor_resource(Neighbor, _, _) & not at_war_with(Neighbor) & not allied_with(Neighbor) & not plans_to_attack(Neighbor) & not sent_alliance_request(Neighbor) & not alliance_declined(Neighbor)), Neighbors);
         if (Neighbors \== []) {
             Neighbors = [BestNeighbor | _];
             +sent_alliance_request(BestNeighbor);
             vesna.via.request_alliance(BestNeighbor);
             .print("Sent alliance request to neighbor ", BestNeighbor);
         }
     }.
\end{lstlisting}

When an alliance request is received (\texttt{+alliance\_request\_received}), the agent accepts via \texttt{vesna.via.accept\_alliance(Sender)} only if its active alliances, sent requests, and accepting states are all strictly below $\mathit{MaxAlliances}$, and it does not plan to attack the sender.

Proactively (\texttt{+!check\_alliance\_opportunities}), if holding zero alliances, the agent queries peaceful neighboring sovereign states and transmits an alliance proposal via \texttt{vesna.via.request\_alliance}.

% ==============================================================================
\section{Belief Base Maintenance and Deliberation Loop Scheduling}
\label{sec:impl_lifecycle}

To preserve synchronization with the real-time Godot simulation, the agent maintains dedicated reactive update plans and a prioritized deliberative audit loop.

\subsection{Reactive Perception Synchronization}
\label{subsec:belief_sync}

When physical events occur in Godot (e.g., division destruction, territorial conquest, or diplomatic treaty establishment), sight packets update the agent's belief base.

\begin{lstlisting}[language=prolog, caption={Reactive state synchronization triggers in \texttt{vesna.asl}.}, label={lst:sync_plans_asl}]
+my_army_update(Id, X, Y, Life, AttackValue, Loc)
  <- .abolish(my_army(Id, _, _, _, _, _));
     +my_army(Id, X, Y, Life, AttackValue, Loc);
     if (ordered_to_defend(Id, Loc)) {
         .abolish(ordered_to_defend(Id, Loc));
         .print("Army ", Id, " has arrived at and successfully defended ", Loc);
     }.

+my_army_destroyed_update(Id)
  <- .abolish(my_army(Id, _, _, _, _, _));
     .abolish(my_army_path_distance(Id, _, _));
     .abolish(ordered_to_defend(Id, _));
     -my_army_travelling(Id);
     if (my_army_target(Id, TargetProv)) {
         -my_army_target(Id, TargetProv);
         if (not my_army_target(_, TargetProv)) {
             -targeted_province(TargetProv);
         }
     };
     .print("Army ", Id, " has been destroyed.").

+allied_province_owner_update(Id, Morale, X, Y, IsCapital, IsDouble, Resource)
  <- if (enemy_province_status(Id, PrevOwner, _, _, _, _, _, _, _, _)) {
         -war_target(PrevOwner, Id);
         if (at_war_with(PrevOwner)) {
             !select_war_target(PrevOwner, _);
         }
     };
     .abolish(allied_province(Id, _, _, _, _, _, _));
     .abolish(enemy_province_status(Id, _, _, _, _, _, _, _, _, _));
     -targeted_province(Id);
     +allied_province(Id, Morale, X, Y, IsCapital, IsDouble, Resource).

+match_over <-
     .drop_all_intentions;
     .drop_all_desires;
     .drop_all_events;
     .print("Match over! Interrupted all goals and intentions.").
\end{lstlisting}

As illustrated in Listing~\ref{lst:sync_plans_asl}:
\begin{itemize}
    \item \texttt{+my\_army\_update}: Updates division coordinates and combat stats. If an army has reached its defense destination, its defensive assignment is automatically retired.
    \item \texttt{+my\_army\_destroyed\_update}: Purges all memory references, routing tables, and target assignments linked to the destroyed division.
    \item \texttt{+allied\_province\_owner\_update}: Triggered upon capturing an enemy territory. It clears the completed \texttt{war\_target}, re-invokes \texttt{!select\_war\_target} to choose the next campaign objective, and incorporates the new territory into domestic beliefs.
    \item \texttt{+match\_over}: Immediately drops all active intentions, desires, and events to halt the cognitive cycle upon simulation end.
\end{itemize}

\subsection{Prioritized Deliberation Cycle Scheduling}
\label{subsec:deliberation_scheduling}

To prevent conflicting behavioral commands, the agent executes its cognitive evaluation loop following a strict priority hierarchy every $5000\text{ ms}$.

\begin{lstlisting}[language=prolog, caption={Prioritized periodic evaluation loop in \texttt{vesna.asl}.}, label={lst:evaluation_loop_asl}]
+active_wars(Count) : not evaluation_loop_started
  <- +evaluation_loop_started;
     !evaluate_defenses;
     !optimize_riot_garrisons;
     !evaluate_targets;
     !evaluate_attack_opportunities;
     !garrison_non_allied_borders;
     !check_peace_requests;
     !check_alliance_opportunities;
     !!periodic_evaluation.

+!periodic_evaluation : true
  <- !evaluate_defenses;
     !optimize_riot_garrisons;
     !evaluate_targets;
     !evaluate_attack_opportunities;
     !garrison_non_allied_borders;
     !check_peace_requests;
     !check_alliance_opportunities;
     .wait(5000);
     !!periodic_evaluation.
\end{lstlisting}

The deliberation loop executes a sequential pipeline:
\begin{equation}
    \label{eq:deliberation_pipeline}
    \text{Defenses} \longrightarrow \text{Riot Garrisons} \longrightarrow \text{Targets} \longrightarrow \text{Attacks} \longrightarrow \text{Border Garrisons} \longrightarrow \text{Peace} \longrightarrow \text{Alliances}
\end{equation}

By placing defensive threat evaluation and riot garrison stabilization ahead of offensive campaign operations, the agent guarantees that national security and internal order are fully preserved before military divisions are committed to foreign expansion.
